\documentclass[11pt,reqno]{amsart}

%==============================================================
% Energy-First Stochastic Calculus:
% Iterated Brownian Motion as a Test Case
%==============================================================

\usepackage[utf8]{inputenc}
\usepackage[T1]{fontenc}
\usepackage{lmodern}
\usepackage[margin=1.15in]{geometry}
\usepackage{amsmath,amssymb,amsthm,mathtools}
\usepackage{mathrsfs}
\usepackage{microtype}
\usepackage{hyperref}
\usepackage{cleveref}

% --- theorem environments
\theoremstyle{plain}
\newtheorem{theorem}{Theorem}[section]
\newtheorem{proposition}[theorem]{Proposition}
\newtheorem{lemma}[theorem]{Lemma}
\newtheorem{corollary}[theorem]{Corollary}

\theoremstyle{definition}
\newtheorem{definition}[theorem]{Definition}
\newtheorem{example}[theorem]{Example}
\newtheorem{assumption}[theorem]{Assumption}
\newtheorem{conjecture}[theorem]{Conjecture}
\newtheorem{program}[theorem]{Program}

\theoremstyle{remark}
\newtheorem{remark}[theorem]{Remark}
\newtheorem*{sketch}{Sketch}

% --- macros
\newcommand{\R}{\mathbb{R}}
\newcommand{\N}{\mathbb{N}}
\newcommand{\E}{\mathbb{E}}
\newcommand{\Prob}{\mathbb{P}}
\newcommand{\Hilb}{\mathcal{H}}
\newcommand{\Dom}{\mathrm{Dom}}
\newcommand{\Ran}{\mathrm{Ran}}
\newcommand{\Gra}{\mathrm{Gra}}
\newcommand{\Energy}{\mathcal{E}}
\newcommand{\Source}{\mathsf{S}}
\newcommand{\Obs}{\mathsf{Z}}
\renewcommand{\iota}{\mathfrak{i}}
\newcommand{\IBM}{Z}
\newcommand{\Carre}{\Gamma}
\DeclareMathOperator{\Var}{Var}
\DeclareMathOperator{\Cov}{Cov}

\title{Energy-First Stochastic Calculus:\\
       Iterated Brownian Motion as a Test Case}
\author{Ramiro Fontes}
\date{\today}

\begin{document}

\begin{abstract}
We begin from \emph{observable It\^o identities} --- evolution
formulas for $t\mapsto\E[f(X_t)]$ on a test class
$\mathcal T$ --- and ask when such an identity admits a pathwise
lift to $f(X_t)-f(X_0)=M_f(t)+C_f(t)$ and a further energy-first
lift in which $M_f$ is a divergence/integral and $C_f$ an
observable-energy contraction (\cref{def:obs-ito-identity}). On this
hierarchy, the classical semimartingale It\^o formula is the
quadratic-energy case at the energy-first level; it is one
specific lift, not the definition. Accordingly, we propose an
energy-first reorganisation of stochastic calculus: take a closed
densely-defined operator derivative $D$ on an energy space
$\Energy$ as the primitive datum and \emph{define} the divergence by
adjunction $\delta=D^{*}$, then construct an intrinsic observable
derivative $D_\Obs$ and projected divergence
$\delta_\Obs u=\E[\delta_\Source(Ju)\mid\mathcal F_\Obs]$ for
composite processes whose observable filtration is strictly smaller
than a source filtration (\cref{thm:compression,thm:proj-div}); a
conditional operator-It\^o transfer principle (\cref{thm:transfer})
shows how source-level It\^o expansions descend to the observable. The flagship
application of this framework is an It\^o formula for iterated
Brownian motion $\IBM_t=X_{|Y_t|}$, with two-sided Brownian scenery
$X$ and independent Brownian clock $Y$, which does not supply the
semimartingale quadratic-variation structure needed for the
classical It\^o formula in its own filtration. For the mollified IBM
$\IBM_t^{\,\varepsilon}=x+X^\varepsilon_{|Y_t|}$ with
$X^\varepsilon=\rho_\varepsilon*X$ ($\rho$ even, non-negative,
compactly supported), $f\in\mathscr S(\R)$, and
$g_\varepsilon(r)=f(x+X^\varepsilon_r)$, we prove the pathwise
identity
\begin{equation*}
  f(\IBM_t^{\,\varepsilon})-f(\IBM_0^{\,\varepsilon})
   \;=\;\delta_Y\bigl(B_{f,\varepsilon}(t)\bigr)
   \;+\;C_{\IBM,\varepsilon}(f;t),
\end{equation*}
where the clock Skorohod-divergence integrand is
$B_{f,\varepsilon}(t)(s)=g_\varepsilon'(|Y_s|)\,\mathrm{sgn}(Y_s)\,
\mathbf 1_{[0,t]}(s)$ and the mollified pathwise correction is
\begin{equation*}
  C_{\IBM,\varepsilon}(f;t)
   \;=\;g_\varepsilon'(0)\,L_t^{0}(Y)
   \;+\;\tfrac{1}{2}\int_0^{t}g_\varepsilon''(|Y_s|)\,ds
\end{equation*}
(\cref{thm:energy-first-ito-moll}). The expectation of this
correction reproduces the fourth-order Brownian-time PDE
\begin{equation*}
  \partial_t u
   \;=\;\tfrac{1}{8}\,\partial_x^{4}u
   \;+\;\frac{\partial_x^{2}f(x)}{2\sqrt{2\pi t}},
  \qquad u(t,x)=\E[f(x+X_{|Y_t|})],
\end{equation*}
in the sense
$\lim_{\varepsilon\to 0}\E[C_{\IBM,\varepsilon}(f;t)]=u(t,x)-f(x)$
(\cref{thm:bt-correction,thm:moll-consistency}). The identity is
pathwise for every $\varepsilon>0$ and fourth-order correct in
expectation as $\varepsilon\to 0$. The unmollified pathwise
correction does \emph{not} exist in classical $L^{2}(\Omega)$:
$\Var(\mathrm{Proj}_1^{X}\,C_{\IBM,\varepsilon}(f;t))\to+\infty$ as
$\varepsilon\to 0$ under the nondegeneracy condition
$(P_\eta f')(x)\not\equiv 0$ on $(0,\infty)$
(\cref{prop:var-divergence}), via an exact cancellation of leading
$\varepsilon^{-1}$ terms across the two summands of
$g_\varepsilon''$ and the Bouleau--Yor spatial quadratic variation
of the clock local time. The natural target is therefore the dual of
the quartic composite observable energy
$\Energy_\IBM^{(4)}=\Energy_\IBM^{(2)}\otimes\Energy_\IBM^{(2)}$ of
the rigged hierarchy \eqref{eq:rigged-4}. The proved IBM It\^o
formula is therefore mollified-pathwise and fourth-order consistent;
the remaining analytic problem is the unmollified correction in the
quartic observable-energy dual.

The clock It\^o formula behind the identity is second order in the
Brownian-time variable, but because its coefficients are derivatives
of Brownian scenery, the observable correction has fourth-order PDE
shadow. IBM is the canonical fourth-order witness for an
energy-order classification of It\^o corrections that supersedes
classical semimartingale quadratic variation; the candidate
second-order observable energy $\Energy_\IBM^{(2)}$
(\cref{def:second-order-ibm}) is a second \emph{operator} derivative
whose carr\'e-form, contracted against the clock, produces the
fourth-order spatial shadow.
\end{abstract}

\maketitle

\section{Introduction: the reversal}\label{sec:intro}

The classical construction of stochastic calculus, in nearly every textbook
treatment, proceeds in the following order:
\begin{enumerate}
\item Fix a Brownian motion $B$ on a probability space.
\item Construct the stochastic (It\^o or Skorohod) integral $\delta$ on
  a class of integrands by an isometry argument.
\item Define the Malliavin derivative $D$ either by chaos expansion or
  as the adjoint $D=\delta^{*}$, and verify duality.
\item Derive the It\^o/Tanaka formulae and identify the PDE shadow of
  the process as the action of the infinitesimal generator on smooth
  functions.
\end{enumerate}
For Brownian motion this ordering is harmless. The Gaussian structure
forces every reasonable derivative/divergence pair to coincide, the
energy space is $L^2(\R_+)$ with the Lebesgue inner product, and the
PDE shadow is the heat equation.

The companion paper~\cite{Fontes-OperatorDerivative} (\emph{The Operator
Derivative in Continuous Stochastic Calculus}) develops a divergence-first
operator framework: it takes the stochastic integral $\delta_X$ as primitive
on a Hilbert energy space $\mathcal H_X$ and defines the operator derivative
by adjointness $D_X\coloneqq\delta_X^{*}$ on its natural adjoint domain.
The representation layer (Clark--Ocone, GKW projection, variance identity)
is shown there to follow from adjointness alone; the calculus layer
(Leibniz, chain rule, It\^o formula) requires a cylindrical reduction with
process-dependent representers. The present paper inverts the order of
this construction.

\subsection*{The hidden distinction}
The ordering hides a structural choice. \emph{Three} pieces of data
are in play: an energy space $\Energy$, an operator derivative
$D\colon L^2(\Omega)\to\Energy$, and a divergence
$\delta\colon\Energy\to L^2(\Omega)$. In the Brownian regime any two
determine the third. Outside that regime --- fractional Brownian
motion, Volterra and ambit processes, stable L\'evy processes,
stochastic volatility models --- the choice of $\Energy$ becomes the
primary modelling decision, and postulating $\delta$ first is a
technical convenience rather than the conceptually correct order.

\subsection*{The reversal, and the compression construction}
We invert the ordering relative to~\cite{Fontes-OperatorDerivative}. The
primitive object is the pair $(\Energy, D)$. The divergence is
\emph{defined} as $\delta\coloneqq D^{*}$. For a composite process in
which an observable $\sigma$-algebra $\mathcal F_\Obs$ is strictly
contained in a richer source $\sigma$-algebra $\mathcal F_\Source$, the
construction goes one step further: a source energy datum
$(\Hilb_\Source,\Energy_\Source,D_\Source)$ induces an intrinsic
observable energy datum $(\Hilb_\Obs,\Energy_\Obs,D_\Obs)$ by
\emph{compression} --- restricting $D_\Source$ to functionals of the
observable $\sigma$-algebra and closing. The intrinsic divergence is the
observable conditional expectation of the source divergence. This is the
content of Section~\ref{sec:compat}.

\subsection*{The flagship example: It\^o for iterated Brownian motion}
The energy-first framework was developed against a concrete
obstruction: iterated Brownian motion
$\IBM_t=X_{|Y_t|}$,
with two-sided Brownian scenery $X$ and independent Brownian clock
$Y$, is not a semimartingale in its own filtration and does not
admit a classical semimartingale It\^o formula of the usual
quadratic-variation form
$\tfrac12\!\int f''(\IBM_s)\,d\langle\IBM\rangle_s$. The flagship payoff of the framework
(\cref{thm:energy-first-ito-moll}) is the following identity for
\emph{mollified} IBM
$\IBM_t^{\,\varepsilon}=x+X^\varepsilon_{|Y_t|}$, where
$X^\varepsilon=\rho_\varepsilon*X$ for an even non-negative
compactly supported $\rho$: for $f\in\mathscr S(\R)$ and
$g_\varepsilon(r)=f(x+X^\varepsilon_r)$,
\begin{equation}\label{eq:flagship-intro}
  f(\IBM_t^{\,\varepsilon})-f(\IBM_0^{\,\varepsilon})
  \;=\;\delta_Y\bigl(B_{f,\varepsilon}(t)\bigr)
   \;+\;C_{\IBM,\varepsilon}(f;t),
\end{equation}
where the clock Skorohod-divergence integrand is
$B_{f,\varepsilon}(t)(s)=g_\varepsilon'(|Y_s|)\,\mathrm{sgn}(Y_s)\,
\mathbf 1_{[0,t]}(s)$ and the mollified pathwise correction is
\begin{equation}\label{eq:flagship-C-intro}
  C_{\IBM,\varepsilon}(f;t)
   \;=\;g_\varepsilon'(0)\,L_t^{0}(Y)
   \;+\;\tfrac{1}{2}\int_0^{t}g_\varepsilon''(|Y_s|)\,ds.
\end{equation}
The clock derivatives are explicit:
$g_\varepsilon'(r)=f'(x+X^\varepsilon_r)\,\partial_r X^\varepsilon_r$
and
$g_\varepsilon''(r)=f''(x+X^\varepsilon_r)\,
(\partial_r X^\varepsilon_r)^{2}
+f'(x+X^\varepsilon_r)\,\partial_r^{2}X^\varepsilon_r$.

The expectation of $C_{\IBM,\varepsilon}(f;t)$ converges as
$\varepsilon\to 0$ to the fourth-order Brownian-time PDE correction
$\partial_t u=\tfrac{1}{8}\,\partial_x^{4}u+\partial_x^{2}f/
(2\sqrt{2\pi t})$ for $u(t,x)=\E[f(x+X_{|Y_t|})]$:
$\lim_{\varepsilon\to 0}\E[C_{\IBM,\varepsilon}(f;t)]=u(t,x)-f(x)$
(\cref{thm:bt-correction,thm:moll-consistency}). The identity
\eqref{eq:flagship-intro} is therefore \emph{pathwise for every
$\varepsilon>0$ and fourth-order correct in expectation as
$\varepsilon\to 0$}. The unmollified pathwise correction does
\emph{not} exist in classical $L^{2}(\Omega)$:
$\Var(\mathrm{Proj}_1^{X}\,C_{\IBM,\varepsilon}(f;t))\to+\infty$
for every $f$ satisfying the nondegeneracy condition
$(P_\eta f')(x)\not\equiv 0$ on $(0,\infty)$
(\cref{prop:var-divergence}), so the natural target is the dual of
the quartic composite observable energy
$\Energy_\IBM^{(4)}=\Energy_\IBM^{(2)}\otimes\Energy_\IBM^{(2)}$ of
the rigged hierarchy \eqref{eq:rigged-4}.

The clock It\^o formula behind \eqref{eq:flagship-intro} is second
order in the Brownian-time variable, but because its coefficients
are derivatives of Brownian scenery the observable correction has
fourth-order PDE shadow. IBM is therefore the canonical
fourth-order witness for an energy-order classification of It\^o
corrections that supersedes classical semimartingale quadratic
variation; the framework of this paper builds the operator/energy
structures needed to formulate, control, and classify such
corrections, and IBM is where it lands a proved theorem. The
proved IBM It\^o formula is mollified-pathwise and fourth-order
consistent; the remaining analytic problem is the unmollified
correction in the quartic observable-energy dual.

\subsection*{Why iterated Brownian motion}
We use iterated Brownian motion
$\IBM_t = X_{|Y_t|}$
--- with two-sided Brownian scenery $X$ and independent Brownian clock
$Y$ --- as the principal test case. Three features make IBM
diagnostically useful.
\begin{enumerate}
\item Its PDE shadow is the bi-Laplacian/local-time PDE signature
  identified by Allouba--Zheng~\cite{AllouZheng}, not the heat equation
  (IBM is not Markov in its own filtration).
\item It is naturally a \emph{composite} process: the source Malliavin
  calculi $(D_X,\delta_X)$ and $(D_Y,\delta_Y)$ live on the two
  independent Gaussian factors and provide an explicit atlas indexed by
  the chaos bidegree $(m,n)$, as developed in~\cite{Fontes-IBMchaos}.
\item It does not fall under the classical semimartingale It\^o
  theory in its own filtration --- it lacks the usual semimartingale
  quadratic-variation structure needed for the classical It\^o
  formula --- so there is no a priori canonical $\delta_\IBM$ to
  take as starting datum.
\end{enumerate}
Applied to IBM, the compression construction yields a minimal
first-order intrinsic derivative on $\mathcal C_\IBM$
(\cref{def:ibm-datum}), made concrete at the mollified occupation
level in \cref{thm:ibm-occ}. The bi-Laplacian/local-time PDE shadow
is organised by the analogous second-order observable energy, while
the principal open analytic target is the unmollified pathwise
correction in the associated quartic observable-energy dual.

\subsection*{Relation to companion papers}
\cite{Fontes-OperatorDerivative} develops the operator-first
representation/calculus split with $\delta$ as primitive; the present
note inverts that ordering. \cite{Fontes-IBMchaos} computes
source-resolved bidegree projections for IBM occupation functionals
and supplies the $(m,n)$-atlas and the spatial
factorisation~\eqref{eq:occ-factor} used in
Section~\ref{sec:ibm-occ}. \cite{Fontes-IBMnumerical} validates the
bidegree variance budgets numerically; not used here.
\cite{Fontes-PDEgenerators} treats PDE shadows of energy structures
generally and is the natural setting for the second-order
verification listed in Section~\ref{sec:tasks}.

\subsection*{Main-results ladder}
The paper organises the results into four groups:
\begin{itemize}
\item[\textbf{(A)}] \emph{General energy-first framework.}
  The observable-compression construction $D_\Obs$ and projected
  divergence formula
  $\delta_\Obs u=\E[\delta_\Source(Ju)\mid\mathcal F_\Obs]$
  (\cref{thm:compression,thm:proj-div}); a conditional operator-It\^o
  transfer principle (\cref{thm:transfer}); a concrete IBM first-order
  datum on the mollified occupation core
  (\cref{thm:ibm-occ,cor:unmollified}); and a candidate second-order
  observable energy $\Energy_\IBM^{(2)}$
  (\cref{def:second-order-ibm,prop:compression-2}) whose
  tensor square $\Energy_\IBM^{(4)}=
  \Energy_\IBM^{(2)}\otimes\Energy_\IBM^{(2)}$ is the home of the
  quartic composite observable energy.
\item[\textbf{(B)}] \emph{Flagship IBM theorem: mollified pathwise
  It\^o formula.} The energy-first It\^o formula for mollified IBM
  (\cref{thm:energy-first-ito-moll}): the pathwise identity
  \eqref{eq:flagship-intro} with the explicit mollified correction
  \eqref{eq:flagship-C-intro}, for every $\varepsilon>0$ and
  $t\ge 0$, and the operator-divergence reading of the martingale
  part. \emph{Proved.}
\item[\textbf{(C)}] \emph{Fourth-order PDE consistency.}
  At the deterministic level, $u(t,x)=\E[f(x+X_{|Y_t|})]$ satisfies
  the bi-Laplacian Brownian-time PDE
  $\partial_t u=\tfrac18\partial_x^{4}u+\partial_x^{2}f/(2\sqrt{2\pi t})$
  (\cref{thm:bt-correction}); the flagship correction of (B) is
  consistent with it after expectation:
  $\lim_{\varepsilon\to 0}\E[C_{\IBM,\varepsilon}(f;t)]=u(t,x)-f(x)$
  (\cref{thm:moll-consistency}). \emph{Proved.}
\item[\textbf{(D)}] \emph{$L^{2}$ obstruction and quartic-dual
  target.} The first-scenery-chaos $L^{2}$-divergence
  \[
    \Var\bigl(\mathrm{Proj}_1^{X}\,C_{\IBM,\varepsilon}(f;t)\bigr)
     \;\to\;+\infty\qquad(\varepsilon\to 0)
  \]
  is \textbf{proved} for every $f\in\mathscr S(\R)$ satisfying the
  nondegeneracy condition $(P_\eta f')(x)\not\equiv 0$ on
  $(0,\infty)$ (\cref{prop:var-divergence}); the Brownian local-time
  input is \cref{lem:lt-sobolev}. Ordinary Wick / Hermite
  renormalisation, which acts in the zeroth and second scenery
  chaoses, does not address the first-chaos singularity
  (\cref{prop:wick-insufficient}). The natural target for the
  unmollified pathwise correction $C_\IBM(f;t)$ is therefore the
  dual of the quartic composite observable energy
  $\Energy_\IBM^{(4)}$ in the rigged hierarchy
  \eqref{eq:rigged-4}; the schematic pathwise identity
  \cref{conj:ibm-ito} remains conditional on this construction.
  \emph{$L^{2}$-obstruction proved; quartic-dual construction is the
  open analytic target.}
\end{itemize}

\subsection*{Contents}
Section~\ref{sec:operator} sets up the energy-first formalism, and
Section~\ref{sec:compat} contains the abstract observable-compression
construction (the framework half of (A)). Section~\ref{sec:ibm-source}
specialises to IBM at the abstract level and sketches the
second-order energy. Section~\ref{sec:ibm-occ} carries out the first
concrete IBM instantiation on the mollified occupation core,
realising the framework piece of (A) with \cref{thm:ibm-occ} as the
first IBM theorem; the Fourier-mode diagnostic of
\cref{prop:fourier-2nd} sits there as a PDE-shadow check.
Section~\ref{sec:second-order} defines the candidate second-order
observable energy (the higher-order half of (A)).
Section~\ref{sec:ibm-ito} contains the operator-It\^o transfer
theorem and the deterministic Brownian-time fourth-order PDE
benchmark (part of (C)), together with the schematic pathwise target
underlying (D). Section~\ref{sec:moll-pathwise} states and proves
the flagship mollified IBM It\^o theorem (B) together with the
expectation-level fourth-order consistency (C).
Section~\ref{sec:toward-unmoll} records the first-chaos $L^{2}$
obstruction (D). Section~\ref{sec:tasks} lists the analytic tasks
that remain for (D).

\subsection*{Convention on IBM}
Outside Section~\ref{sec:ibm-occ}, the canonical example treated in
this paper is $\IBM_t=X_{|Y_t|}$. Section~\ref{sec:ibm-occ}, which
develops the spatial-factorisation atlas of the occupation functional
$F_T(\varphi)=\int_0^{T}\varphi(\IBM_s)\,ds$, uses the two-sided
composition $X_{Y_s}$ for that section only, since it produces the
cleaner $\R$-valued local-time decomposition \eqref{eq:occ-factor}
(see~\cite{Fontes-IBMchaos}); the canonical $X_{|Y_t|}$ form
corresponds to the half-line restriction, with $L^{x}(Y)$ on $\R$
replaced by $L^{x}(|Y|)$ on $\R_{+}$. This convention switch is
local to Section~\ref{sec:ibm-occ} and is recalled at its start.

\section{Energy-first operator calculus}\label{sec:operator}

We work throughout on a probability space $(\Omega,\mathcal F,\Prob)$
and write $L^2(\Omega)=L^2(\Omega,\mathcal F,\Prob)$.

\subsection*{Three levels of an It\^o formula}\label{sec:three-levels}

The paper is organised around the principle that an It\^o formula
should be read at three nested levels of structure: a baseline
\emph{observable} layer, a \emph{pathwise} lift, and an
\emph{energy-first / operator} lift. The classical semimartingale
It\^o formula is the quadratic-energy case of the third level; it
is one specific lift, not the definition.

\begin{definition}[Observable It\^o identity]\label{def:obs-ito-identity}
Let $\mathcal T\subseteq C(\R)$ be a class of test functions and
$X=(X_t)_{t\ge 0}$ an $\R$-valued process with
$\E|f(X_t)|<\infty$ for every $f\in\mathcal T$ and $t\ge 0$. We
say that $X$ admits an \emph{observable It\^o identity} on
$\mathcal T$ if, for each $f\in\mathcal T$, the map
$t\mapsto\E[f(X_t)]$ is locally absolutely continuous on
$(0,\infty)$ and there exist a family of linear operators
$\mathcal G_t$ on $\mathcal T$ and a family of linear functionals
$\mathcal B_t$ on $\mathcal T$ such that
\begin{equation}\label{eq:obs-ito-identity}
  \frac{d}{dt}\,\E[f(X_t)]
   \;=\;\E\bigl[(\mathcal G_t f)(X_t)\bigr]
   \;+\;\mathcal B_t f
\end{equation}
in the sense of distributions in $t$, or pointwise wherever the
left-hand-side derivative exists.
\end{definition}

We call $\mathcal G_t$ the \emph{observable correction operator}
and $\mathcal B_t$ the \emph{boundary/source term}. For Brownian
motion, $\mathcal G=\tfrac12\partial_x^{2}$ and $\mathcal B=0$; for
a general Markov process the observable correction operator is the
infinitesimal generator. We do not require $X$ to be Markovian, nor
to be a semimartingale: \cref{def:obs-ito-identity} is purely an
evolution identity for test-function expectations, and assumes no
filtration structure on $X$.

\paragraph{The hierarchy.} On a given class $\mathcal T$:
\begin{itemize}
\item[\textbf{Level 1.}] \emph{Observable It\^o identity}
  \eqref{eq:obs-ito-identity}: an evolution identity for
  $t\mapsto\E[f(X_t)]$. This is the weakest level; it depends only
  on the one-dimensional time marginals of $X$ paired with
  $\mathcal T$.
\item[\textbf{Level 2.}] \emph{Pathwise lift}: a decomposition
  \begin{equation}\label{eq:pathwise-lift}
    f(X_t)-f(X_0)\;=\;M_f(t)\;+\;C_f(t),
    \qquad \E[M_f(t)]=0,
  \end{equation}
  where $C_f$ is a process whose time-derivative expectation
  reproduces the right-hand side of \eqref{eq:obs-ito-identity}
  (so that integrating in time and taking expectation recovers
  Level~1), and $M_f$ is a centred fluctuation.
\item[\textbf{Level 3.}] \emph{Energy-first / operator lift}: a
  further lift of \eqref{eq:pathwise-lift} in which $M_f$ is
  represented by a divergence/integral $\delta(U_f)$ for an energy
  datum $(\Hilb,\Energy,D)$ in the sense of \cref{def:energy-datum}
  below, and $C_f$ is represented by an observable-energy
  contraction associated with $D$ or a higher-order analogue
  of $D$.
\end{itemize}

\paragraph{The lift is additional structure, not automatic.}
A Level-1 observable It\^o identity does not by itself produce a
Level-2 pathwise lift: such a lift requires $C_f$ to exist
pathwise in a suitable topology, which is a structural property of
$X$ relative to $\mathcal T$ and may fail when $X$ is not a
semimartingale or when $\mathcal G_t$ is of higher order than two.
Similarly, a Level-2 lift is not automatically Level-3: realising
$M_f$ as $\delta(U_f)$ and $C_f$ as an energy contraction requires
an energy datum whose divergence/contraction reproduces the
pathwise decomposition. The classical semimartingale It\^o formula
sits at Level~3 with $\Energy=L^{2}(\R_{+})$, $D$ the Malliavin
derivative, and the quadratic-energy contraction
$C_f(t)=\tfrac12\int_0^t f''(X_s)\,d\langle X\rangle_s$: it is one
specific energy-first lift, not the definition of an It\^o formula.

\paragraph{IBM as a fourth-order observable It\^o process.}
For iterated Brownian motion $\IBM_t=X_{|Y_t|}$ with
$\IBM_0=x$ and $f\in\mathscr S(\R)$, the half-normal integration
of the heat semigroup yields the observable It\^o identity
\begin{equation}\label{eq:obs-ito-ibm}
  \frac{d}{dt}\,\E[f(\IBM_t)]
   \;=\;\tfrac{1}{8}\,\E\bigl[f^{(4)}(\IBM_t)\bigr]
   \;+\;\frac{1}{2\sqrt{2\pi t}}\,f''(x),
  \qquad t>0,
\end{equation}
which is \cref{thm:bt-correction} written in test-function-expectation
form. Thus IBM is an observable It\^o process of \emph{fourth} order,
with observable correction operator
$\mathcal G_t=\tfrac{1}{8}\,\partial_x^{4}$ and boundary/source term
$\mathcal B_t f=f''(x)/(2\sqrt{2\pi t})$. The purpose of the
flagship mollified pathwise theorem of \cref{sec:moll-pathwise}
(\cref{thm:energy-first-ito-moll}) is to lift this fourth-order
observable identity to a Level-2/3 identity at scenery-mollified
level: \eqref{eq:flagship-intro} is the lift, and
\eqref{eq:efi-fourth} verifies that its expectation reproduces
\eqref{eq:obs-ito-ibm}. The unmollified Level-2/3 lift is not
constructed here as an $L^{2}(\Omega)$ object; the analytic
obstacle (\cref{prop:var-divergence}) shows that the correct
target topology is the dual of the quartic composite observable
energy
$\Energy_\IBM^{(4)}=\Energy_\IBM^{(2)}\otimes\Energy_\IBM^{(2)}$
of the rigged hierarchy \eqref{eq:rigged-4}.

\subsection*{Energy datum and the operator/divergence layer}

We now state the structural data underlying Level~3. The rest of
this section sets up the energy datum, the carr\'e du champ, and
the three-layer (R)/(P)/(I) decomposition internal to Level~3.

\begin{definition}[Energy datum]\label{def:energy-datum}
An \emph{energy datum} is a triple $(\Hilb,\Energy, D)$ where
$\Hilb\subseteq L^2(\Omega)$ is a closed subspace,
$\Energy$ is a separable Hilbert space, and
\[
  D\colon \Dom(D)\subseteq\Hilb \longrightarrow \Energy
\]
is a closed, densely defined linear operator. The
\emph{divergence} is the Hilbert-space adjoint
$\delta\coloneqq D^{*}\colon\Dom(\delta)\subseteq\Energy\to\Hilb$.
\end{definition}

The closedness assumption is what makes adjunction well-posed: by the
von Neumann adjoint theorem, $\delta=D^{*}$ is then also closed and
densely defined.

\begin{definition}[Carr\'e du champ]\label{def:carre}
On $\Dom(D)\times\Dom(D)$ define the bilinear form
$\Carre(F,G)\coloneqq\langle DF,\,DG\rangle_\Energy$.
We call $\Carre$ the \emph{energy bracket} associated to the datum.
\end{definition}

In the classical Wiener space $\Hilb=L^2(\Omega)$,
$\Energy=L^2(\R_+)$, $D$ the Malliavin derivative, the bracket reduces
to the standard Malliavin matrix entry; the present formalism does not
require this specialisation.

\paragraph{Layer structure.}
Once the datum is fixed, three layers are constructed in sequence:

\medskip
\noindent\textbf{(R) Representation layer.}
Given $(\Hilb,\Energy,D)$, the integration-by-parts formula
\begin{equation}\label{eq:ibp}
  \E\bigl[F\,\delta u\bigr]
  \;=\;\langle DF,\,u\rangle_{\Energy},\qquad
  F\in\Dom(D),\;u\in\Dom(\delta),
\end{equation}
holds by definition of adjunction. \emph{Existence of a chaos
expansion} adapted to $D$ is then an independent assumption on
$(\Hilb,\Energy,D)$, e.g.~that $D$ admits a Wiener chaos diagonalisation.

\medskip
\noindent\textbf{(P) Product/chain layer.}
A \emph{product rule} for $D$, of the form
$D(FG)=F\,DG+G\,DF$ on a suitable algebra, and a \emph{chain rule}
$D(f(F))=f'(F)\,DF$ on smooth $f$, are not consequences of
\cref{def:energy-datum}. They have to be imposed as further hypotheses
on the datum; when they hold, we call the datum a \emph{calculus datum}.

\medskip
\noindent\textbf{(I) It\^o layer.}
Once (R) and (P) are in place, an It\^o-type formula
\begin{equation}\label{eq:operator-ito}
  f(F_t)-f(F_0) \;=\; \delta\bigl(U_f(F_\bullet)\bigr) + C_f(F_\bullet),
\end{equation}
with $U_f$ a representation integrand built from $f',DF$ and
$C_f$ an \emph{energy correction} built from $f''$ and $\Carre(F,F)$,
becomes a candidate identity. We do not assert that
\cref{eq:operator-ito} holds for every calculus datum; we record it as
the scheme that must be verified in each instance.

\begin{remark}[Conditional, not automatic]\label{rem:conditional}
The energy-first framework does not produce \cref{eq:operator-ito}
for free. What it does produce is a sharp separation between the
two ingredients of an It\^o formula: a divergence representative
$U_f$ in $\Energy$, and a correction term whose form is dictated by
$\Carre$. For Gaussian/Brownian data the correction reduces to the
standard $\tfrac12 f''(F)\,d\langle F\rangle$; for non-Gaussian data
it need not.
\end{remark}

\begin{remark}[Two-layer separation, inherited]\label{rem:two-layer}
The representation/calculus separation
of~\cite{Fontes-OperatorDerivative} survives the reversal: the
integration-by-parts formula~\eqref{eq:ibp} and any chaos representation
adapted to $D$ form the representation layer, while a chain/product
rule for $D$ is the additional input needed for the calculus layer. Such
a chain rule is equivalent, in the sense
of~\cite[Thm.~B]{Fontes-OperatorDerivative}, to a \emph{cylindrical
reduction} --- the existence of a dense algebra of test functionals
$f(\xi_1,\dots,\xi_n)$ and representers $\kappa_i\in L^2(\Omega;\Energy)$
through which the integration-by-parts formula factors. The classical
Malliavin calculus is the special case in which the representers are the
deterministic Cameron--Martin elements; in the present setting it is
recovered as one cylindrical reduction among others, not as the
definition of $D$. The difference between the two orderings is one of
\emph{starting datum}: in the energy-first ordering, both layers are
read off $(\Energy, D)$ rather than off $(\mathcal H_X, \delta_X)$.
\end{remark}

\section{Construction of the intrinsic observable derivative}\label{sec:compat}

This section contains the central abstract content of the paper. Given
a closed densely defined source derivative $D_\Source$ and an observable
$\sigma$-algebra $\mathcal F_\Obs\subseteq\mathcal F_\Source$, we
construct a closed first-order intrinsic derivative $D_\Obs$ and its
adjoint $\delta_\Obs$. On the subspace of $u\in\Energy_\Obs$ for which
$Ju\in\Dom(\delta_\Source)$, the intrinsic divergence is given by the
projected formula $\delta_\Obs u=P_\Obs\delta_\Source(Ju)$; we do not
claim this characterises all of $\Dom(\delta_\Obs)$. Exact intertwining
$\iota\,\delta_\Obs=\delta_\Source J$ is recovered as a special case.

\subsection*{Setup}

Let $\mathcal F_\Source\supseteq\mathcal F_\Obs$ be sub-$\sigma$-algebras
of $\mathcal F$. Set
\[
  \Hilb_\Source\coloneqq L^2(\Omega,\mathcal F_\Source,\Prob),\qquad
  \Hilb_\Obs\coloneqq L^2(\Omega,\mathcal F_\Obs,\Prob),
\]
and let
\[
  \iota\colon \Hilb_\Obs\hookrightarrow\Hilb_\Source
\]
denote the canonical isometric inclusion. Its Hilbert-space adjoint is
the conditional expectation onto $\mathcal F_\Obs$,
\[
  P_\Obs\coloneqq\iota^{*}=\E[\,\cdot\mid\mathcal F_\Obs]\colon
  \Hilb_\Source\to\Hilb_\Obs.
\]
Let $(\Hilb_\Source,\Energy_\Source,D_\Source)$ be a source energy datum
in the sense of \cref{def:energy-datum}, with divergence
$\delta_\Source=D_\Source^{*}$.

\subsection*{The compression datum}

\begin{definition}[Observable smooth core]\label{def:obs-core}
The \emph{observable smooth core} induced by $D_\Source$ is
\[
  \mathcal C_\Obs
  \;\coloneqq\;
  \bigl\{\,F\in\Hilb_\Obs : \iota F\in\Dom(D_\Source)\,\bigr\}.
\]
\end{definition}

\begin{assumption}[Density of the observable core]\label{ass:dense}
$\mathcal C_\Obs$ is dense in $\Hilb_\Obs$.
\end{assumption}

\Cref{ass:dense} is a structural input. Its verification depends on
$\mathcal F_\Obs\hookrightarrow\mathcal F_\Source$ and on $D_\Source$; for
IBM it is among the open analytic tasks of Section~\ref{sec:tasks}.

\begin{definition}[Compression operator and observable energy]
\label{def:compression}
For $F\in\mathcal C_\Obs$, set
\[
  D_\Obs^{0} F \;\coloneqq\; D_\Source(\iota F) \;\in\; \Energy_\Source.
\]
Define the observable energy space
\[
  \Energy_\Obs
  \;\coloneqq\;
  \overline{D_\Obs^{0}\,\mathcal C_\Obs}^{\,\Energy_\Source}
  \;\subseteq\;\Energy_\Source,
\]
a closed subspace of $\Energy_\Source$, and let
$J\colon\Energy_\Obs\hookrightarrow\Energy_\Source$ denote the isometric
inclusion. Regard $D_\Obs^{0}$ henceforth as a linear operator
$\mathcal C_\Obs\to\Energy_\Obs$.
\end{definition}

\subsection*{Closability and the intrinsic derivative}

\begin{theorem}[Observable compression derivative]\label{thm:compression}
Under \cref{ass:dense}, the operator
$D_\Obs^{0}\colon\mathcal C_\Obs\subseteq\Hilb_\Obs\to\Energy_\Obs$ is
closed on the maximal pullback domain $\mathcal C_\Obs$. We denote it
$D_\Obs$ and call it the \emph{minimal observable derivative} induced by
$D_\Source$. Consequently, any restriction
$D_\Obs^{0}|_{\mathcal C'}$ to a subcore
$\mathcal C'\subseteq\mathcal C_\Obs$ that is dense in $\Hilb_\Obs$ is
closable, and the closure of any such restriction is contained in
$D_\Obs$.
\end{theorem}

\begin{proof}
Let $F_n\in\mathcal C_\Obs$ with $F_n\to F$ in $\Hilb_\Obs$ and
$D_\Obs^{0}F_n\to h$ in $\Energy_\Obs$. Since $\iota$ is an isometric
inclusion, $\iota F_n\to\iota F$ in $\Hilb_\Source$; since $J$ is an
isometric inclusion,
$J\,D_\Obs^{0}F_n=D_\Source(\iota F_n)\to Jh$ in $\Energy_\Source$.
Closedness of $D_\Source$ yields $\iota F\in\Dom(D_\Source)$, hence
$F\in\mathcal C_\Obs$, and $D_\Source(\iota F)=Jh$, whence
$D_\Obs^{0}F=h$. So the graph of $D_\Obs^{0}$ on $\mathcal C_\Obs$ is
closed.

The second statement is a direct consequence: if
$\mathcal C'\subseteq\mathcal C_\Obs$ is dense in $\Hilb_\Obs$, the same
graph-closure argument applied to sequences in $\mathcal C'$ converging
in $\Hilb_\Obs\times\Energy_\Obs$ to a graph point lands in
$\mathcal C_\Obs$ and in the graph of $D_\Obs$. In particular, the
closure of $D_\Obs^{0}|_{\mathcal C'}$ exists as a closed operator and
is contained in $D_\Obs$.
\end{proof}

The intrinsic divergence is now defined by adjunction.

\begin{definition}[Intrinsic divergence]\label{def:int-div}
\[
  \delta_\Obs\coloneqq D_\Obs^{*}\colon\Dom(\delta_\Obs)\subseteq\Energy_\Obs
  \longrightarrow\Hilb_\Obs.
\]
\end{definition}

\subsection*{The projected divergence formula}

\begin{theorem}[Projected divergence formula]\label{thm:proj-div}
Let $u\in\Energy_\Obs$ satisfy $Ju\in\Dom(\delta_\Source)$. Then
$u\in\Dom(\delta_\Obs)$ and
\begin{equation}\label{eq:proj-div}
  \delta_\Obs u \;=\; P_\Obs\,\delta_\Source(Ju).
\end{equation}
\end{theorem}

\begin{proof}
Let $F\in\mathcal C_\Obs$. By \cref{thm:compression} and the inclusion
isometry $J$,
\[
  \langle D_\Obs F,\,u\rangle_{\Energy_\Obs}
  \;=\;\langle J\,D_\Obs F,\,Ju\rangle_{\Energy_\Source}
  \;=\;\langle D_\Source(\iota F),\,Ju\rangle_{\Energy_\Source}.
\]
By source integration-by-parts (\cref{eq:ibp} applied to $D_\Source$)
and the hypothesis $Ju\in\Dom(\delta_\Source)$,
\[
  \langle D_\Source(\iota F),\,Ju\rangle_{\Energy_\Source}
  \;=\;\langle\iota F,\,\delta_\Source(Ju)\rangle_{\Hilb_\Source}.
\]
Applying $\iota^{*}=P_\Obs$,
\[
  \langle\iota F,\,\delta_\Source(Ju)\rangle_{\Hilb_\Source}
  \;=\;\langle F,\,P_\Obs\,\delta_\Source(Ju)\rangle_{\Hilb_\Obs}.
\]
Chaining,
\begin{equation}\label{eq:chain}
  \langle D_\Obs F,\,u\rangle_{\Energy_\Obs}
  \;=\;\langle F,\,P_\Obs\,\delta_\Source(Ju)\rangle_{\Hilb_\Obs}
  \qquad\text{for all }F\in\mathcal C_\Obs.
\end{equation}
The right-hand side of \cref{eq:chain} is a bounded linear functional of
$F\in\Hilb_\Obs$, with bound
$\|P_\Obs\,\delta_\Source(Ju)\|_{\Hilb_\Obs}$. Since
$\Dom(D_\Obs)=\mathcal C_\Obs$ by \cref{thm:compression}, the identity
holds on all of $\Dom(D_\Obs)$. Hence $u\in\Dom(D_\Obs^{*})=
\Dom(\delta_\Obs)$ and $\delta_\Obs u=P_\Obs\,\delta_\Source(Ju)$.
\end{proof}

\begin{remark}[Automaticity of the $L^{2}$ hypothesis]
\label{rem:l2-auto}
In more abstract versions of this construction one finds the additional
hypothesis $P_\Obs\delta_\Source(Ju)\in\Hilb_\Obs$. In our setup this is
automatic: $\delta_\Source(Ju)\in\Hilb_\Source$ by definition of
$\Dom(\delta_\Source)$, and $P_\Obs=\E[\cdot\mid\mathcal F_\Obs]$ is an
$L^{2}$-contraction, so the conclusion always lies in $\Hilb_\Obs$.
\end{remark}

\begin{remark}[Reduction to exact intertwining]\label{rem:exact-special}
If $\delta_\Source(Ju)$ is $\mathcal F_\Obs$-measurable --- equivalently,
fixed by $P_\Obs$ --- then \cref{eq:proj-div} reduces to the exact
intertwining
\[
  \iota\,\delta_\Obs u \;=\; \delta_\Source(Ju),
\]
recovering the special case stated below in
\cref{cor:exact-intertwining}. For general lifts the divergence carries
a $P_\Obs$, which is the content of the present construction.
\end{remark}

\subsection*{Exact intertwining as a special case}

\begin{definition}[Exact divergence intertwining]\label{def:dintertwine}
The lift $J$ \emph{intertwines the divergences exactly} if there is a
subspace $\mathcal C\subseteq\Dom(\delta_\Obs)$, dense in $\Energy_\Obs$,
such that for every $u\in\mathcal C$, $Ju\in\Dom(\delta_\Source)$,
$\delta_\Source(Ju)$ is $\mathcal F_\Obs$-measurable, and
\[
  \iota\,\delta_\Obs u \;=\; \delta_\Source(Ju).
\]
\end{definition}

\begin{corollary}[Adjoint compatibility, exact case]\label{cor:exact-intertwining}
Suppose $D_\Obs$ is the compression derivative of \cref{thm:compression}
and the inclusion $J$ intertwines the divergences exactly in the sense
of \cref{def:dintertwine}. Then for every $F\in\mathcal C_\Obs$,
\[
  D_\Obs F \;=\; J^{*}\,D_\Source(\iota F),
\]
which is the rigidity statement of the previous draft of this paper.
\end{corollary}

\begin{proof}
By \cref{thm:proj-div} together with exact intertwining,
$\delta_\Obs u = P_\Obs\delta_\Source(Ju) = \delta_\Source(Ju)$ on
$\mathcal C$. The identity for $D_\Obs F$ follows by taking the adjoint
of $J$ on $\mathcal C$: for $F\in\mathcal C_\Obs$ and $u\in\mathcal C$,
\[
  \langle D_\Obs F,u\rangle_{\Energy_\Obs}
  = \langle D_\Source(\iota F),Ju\rangle_{\Energy_\Source}
  = \langle J^{*}D_\Source(\iota F),u\rangle_{\Energy_\Obs},
\]
and density of $\mathcal C$ in $\Energy_\Obs$ concludes.
\end{proof}

\subsection*{Domain caveats}

\begin{remark}[Density of $\mathcal C_\Obs$ is non-trivial]
\label{rem:density-caveat}
Without \cref{ass:dense} the compression still produces a closable
operator on a non-dense domain, but $\delta_\Obs=D_\Obs^{*}$ is no longer
guaranteed to be densely defined and the duality machinery is lost. For
abstract Wiener spaces with a smooth source datum, density is typically
inherited from smooth cylindrical functionals; for IBM it is more
delicate and listed as an open task.
\end{remark}

\begin{remark}[Minimality]\label{rem:minimality}
$D_\Obs$ is minimal in the following sense. Any closed operator
$\widetilde D\colon\Dom(\widetilde D)\subseteq\Hilb_\Obs\to\Energy_\Source$
that extends $F\mapsto D_\Source(\iota F)$ on $\mathcal C_\Obs$ has
$\Dom(D_\Obs)\subseteq\Dom(\widetilde D)$ and
$\widetilde D|_{\Dom(D_\Obs)}=JD_\Obs$. Choosing a different source
derivative produces, in general, a different observable derivative; the
compression is a functor from source-data to observable-data, not a
canonical construction independent of the source.
\end{remark}

\begin{remark}[Energy norm transfer]\label{rem:energy-norm}
Since $J$ is an isometric inclusion, for every $F\in\Dom(D_\Obs)$,
\[
  \|D_\Obs F\|_{\Energy_\Obs}
  \;=\; \|D_\Source(\iota F)\|_{\Energy_\Source}.
\]
The intrinsic energy is strictly the source energy restricted to those
modes that arise from observable functionals --- not a coarsened or
contracted version of it.
\end{remark}

\section{IBM: the compression construction}\label{sec:ibm-source}

We now apply the abstract construction of Section~\ref{sec:compat} to
iterated Brownian motion. The output, under the stated density
hypothesis, is a first-order intrinsic observable derivative $D_\IBM$
on the maximal pullback core and a projected divergence formula for
$\delta_\IBM$ on the lifted source-divergence domain; the second-order
analogue, which is the natural home for the bi-Laplacian/local-time
PDE shadow, is sketched in Section~\ref{sec:ibm-source-2nd}.

\subsection*{The IBM source datum}

Let $X=(X_x)_{x\in\R}$ be a two-sided Brownian scenery and
$Y=(Y_t)_{t\ge 0}$ an independent Brownian clock, with
$\mathcal F_\Source=\mathcal F_X\vee\mathcal F_Y$ the source
$\sigma$-algebra and
$\mathcal F_\IBM=\sigma(\IBM_s:s\ge 0)\subseteq\mathcal F_\Source$
the observable $\sigma$-algebra, where $\IBM_t=X_{|Y_t|}$. The joint
Wiener chaos decomposition is
\[
  L^2(\Omega_X\times\Omega_Y)
  =\bigoplus_{m,n\ge 0}\Hilb^{X}_{m}\otimes\Hilb^{Y}_{n}.
\]
The partial Malliavin derivatives $D_X$, $D_Y$ on each Gaussian factor
combine into
\[
  D_\Source=(D_X,D_Y),\qquad
  \Dom(D_\Source)=\Dom(D_X)\cap\Dom(D_Y),
\]
mapping into the source energy
$\Energy_\Source=L^2(\R)\oplus L^2(\R_+)$. The source divergence acts
componentwise on $\Energy_\Source$,
\begin{equation}\label{eq:delta-S}
  \delta_\Source(u_X,u_Y) \;=\; \delta_X(u_X)+\delta_Y(u_Y),
\end{equation}
for $u_X\in\Dom(\delta_X)$, $u_Y\in\Dom(\delta_Y)$, where
$\delta_X=D_X^{*}$ and $\delta_Y=D_Y^{*}$ are the source Skorohod
divergences~\cite{Nualart,Janson}. The bidegree-$(m,n)$ chaos projection
is the explicit coordinate atlas of~\cite{Fontes-IBMchaos}.

\subsection*{The intrinsic IBM datum}

Applying \cref{def:obs-core,def:compression,thm:compression,thm:proj-div}
to $\mathcal F_\IBM\subseteq\mathcal F_\Source$ and $D_\Source$:

\begin{definition}[Intrinsic IBM datum]\label{def:ibm-datum}
\hfill
\begin{enumerate}
\item The observable smooth core of IBM is
\[
  \mathcal C_\IBM
  \;=\;\bigl\{\,F\in L^{2}(\Omega,\mathcal F_\IBM,\Prob):
        \iota F\in\Dom(D_X)\cap\Dom(D_Y)\,\bigr\}.
\]
\item The observable energy of IBM is
\[
  \Energy_\IBM
  \;=\;\overline{\bigl\{(D_X\iota F,\,D_Y\iota F):F\in\mathcal C_\IBM\bigr\}}
       ^{\,L^{2}(\R)\oplus L^{2}(\R_{+})}
  \;\subseteq\; L^{2}(\R)\oplus L^{2}(\R_{+}).
\]
\item The intrinsic IBM derivative is
\[
  D_\IBM F \;=\; (D_X\iota F,\,D_Y\iota F) \;\in\;\Energy_\IBM,
  \qquad F\in\mathcal C_\IBM,
\]
closed on $\mathcal C_\IBM$ by \cref{thm:compression}.
\item The intrinsic IBM divergence $\delta_\IBM=D_\IBM^{*}$ admits the
projected formula, by \cref{thm:proj-div},
\begin{equation}\label{eq:ibm-delta}
  \delta_\IBM u \;=\; \E\!\bigl[\,\delta_X u_X+\delta_Y u_Y\,\big|\,
   \mathcal F_\IBM\,\bigr]
\end{equation}
for $u=(u_X,u_Y)\in\Energy_\IBM$ with $u_X\in\Dom(\delta_X)$ and
$u_Y\in\Dom(\delta_Y)$.
\end{enumerate}
\end{definition}

We adopt as a standing hypothesis that $\mathcal C_\IBM$ is dense in
$L^{2}(\Omega,\mathcal F_\IBM,\Prob)$. The verification of this density
on natural concrete test classes (smooth bounded functions of finitely
many IBM increments, smooth integral functionals
$\int_{0}^{T}\varphi(\IBM_{s})\,ds$, etc.) is listed in
Section~\ref{sec:tasks}.

\Cref{def:ibm-datum} should be read as the assertion: \emph{the minimal
first-order intrinsic observable derivative for IBM, induced by the
source atlas, is a closed operator on the maximal pullback domain
$\mathcal C_\IBM$, and its adjoint $\delta_\IBM$ admits the projected
conditional-expectation formula \eqref{eq:ibm-delta} on the lifted
source-divergence domain}. No further input from the IBM chaos paper is
needed for the construction itself; the chaos paper becomes relevant
when computing $D_\IBM$ mode by mode.

\subsection*{What is constructed and what is not}

The construction is genuinely a construction, not a programme. It
\emph{does} produce a closed first-order operator on the IBM core
$\mathcal C_\IBM$, with the projected divergence
formula~\eqref{eq:ibm-delta} valid on
$\{u\in\Energy_\IBM:Ju\in\Dom(\delta_\Source)\}$. Density of
$\mathcal C_\IBM$ in $L^{2}(\mathcal F_\IBM)$ is a standing hypothesis,
not a theorem proved here; a concrete inner-product realisation of
$\Energy_\IBM$ is likewise not provided. It does \emph{not} produce,
on its own, the full first-order calculus of IBM: a chain or product
rule for $D_\IBM$ is not part of the definition and must be verified
via a cylindrical reduction (\cref{rem:two-layer}). The natural
candidate reduction inherits random representers from the joint
$(X,Y)$ Brownian Malliavin gradient projected onto $\mathcal F_\IBM$;
its domain of validity is again analytic, not formal.

\subsection*{Source bidegree compatibility}

For $F\in\mathcal C_\IBM$ with joint Wiener chaos expansion
$F=\sum_{m,n\ge 0}F_{m,n}\in\Hilb^{X}_{m}\otimes\Hilb^{Y}_{n}$, the
partial source derivatives $D_X F_{m,n}$ and $D_Y F_{m,n}$ are explicit
by source Malliavin calculus, and
\[
  D_\IBM F \;=\; \sum_{m,n\ge 0}\bigl(D_X F_{m,n},\,D_Y F_{m,n}\bigr)
  \;\in\;\Energy_\IBM.
\]
The chaos paper~\cite{Fontes-IBMchaos} supplies the bidegree kernels;
the present construction supplies the operator. Together they pin down
$D_\IBM$ mode by mode on $\mathcal C_\IBM$. We do not assert that this
mode-by-mode description gives a closed-form $\Energy_\IBM$;
identifying $\Energy_\IBM$ concretely is task~(2) of
Section~\ref{sec:tasks}.

\subsection{Second-order energy and the PDE shadow}\label{sec:ibm-source-2nd}

The bi-Laplacian/local-time PDE signature of~\cite{AllouZheng} is a
\emph{second-order} object: it lives in the Brownian-time PDE shadow
for the deterministic averaged observable $u(t,x)=\E[f(x+X_{|Y_t|})]$,
not in any first-order pathwise representation of $\IBM$ (which is
not Markov in its own filtration). The first-order datum of
\cref{def:ibm-datum} therefore should not be expected, on its own, to
display the $\partial_x^{4}$ correction.

We sketch the natural higher-order construction. Let
$D_\Source^{(2)}\coloneqq D_\Source\otimes D_\Source$ act on the
second-order Sobolev space of the source via iteration of the partial
Malliavin derivatives, with values in
$\Energy_\Source^{\otimes 2}=L^{2}(\R^{2})\oplus L^{2}(\R\times\R_{+})
\oplus L^{2}(\R_{+}^{2})$ (or a symmetrised version, depending on
convention). Define the second-order observable smooth core
\[
  \mathcal C_\IBM^{(2)}
  \;=\;\bigl\{\,F\in L^{2}(\Omega,\mathcal F_\IBM):
        \iota F\in\Dom(D_\Source^{(2)})\,\bigr\}
\]
and (assuming density) the second-order observable energy
\[
  \Energy_\IBM^{(2)}
  \;=\;\overline{\bigl\{D_\Source^{(2)}\iota F:F\in\mathcal C_\IBM^{(2)}\bigr\}}
       ^{\,\Energy_\Source^{\otimes 2}}.
\]
The same compression argument as \cref{thm:compression} closes the
operator $F\mapsto D_\Source^{(2)}(\iota F)$ on $\mathcal C_\IBM^{(2)}$,
yielding a closed second-order intrinsic operator
\[
  D_\IBM^{(2)}\colon\mathcal C_\IBM^{(2)}\subseteq L^{2}(\mathcal F_\IBM)
  \longrightarrow\Energy_\IBM^{(2)}.
\]
The formal development of $D_\IBM^{(2)}$ as a closed operator and its
relation to the source bidegrees and the Fourier PDE shadow is given
in Section~\ref{sec:second-order} (\cref{prop:compression-2,def:second-order-ibm}).
We do not introduce a second-order intrinsic divergence symbol; the
adjoint of $D_\IBM^{(2)}$ is a candidate object whose closed-form
action is part of the next-stage analysis, not constructed here.

\begin{remark}[Where the bi-Laplacian should appear]\label{rem:bilap-where}
The proposed mechanism is that the bi-Laplacian/local-time correction
of~\cite{AllouZheng} appears as a structural component of the
second-order $\Carre$-form built from $D_\IBM^{(2)}$, averaged against
the law of the Brownian clock $Y$. We do not prove this here. The
source-resolved chaos atlas of~\cite{Fontes-IBMchaos} provides the
expected consistency check: at bidegree $(m,n)$ the second-order energy
of $F\in\mathcal C_\IBM^{(2)}$ should be computable in closed form, and
matching against the Allouba--Zheng PDE signature becomes a finite
verification at each bidegree. We make no claim that this matching is
established.
\end{remark}

The first-order construction supplies the observable representation
layer. The second-order construction is the proposed route to the PDE
shadow. The It\^o identity of Section~\ref{sec:ibm-ito} sits between
them.

\section{First IBM theorem: the mollified occupation core}\label{sec:ibm-occ}

\paragraph{Scope.} The construction below is first-order and
core-level. It defines a \emph{mollified} occupation observable
derivative and its projected divergence on an explicit non-empty
subspace of $\mathcal C_\IBM$. The full unmollified intrinsic IBM
calculus --- including density of the occupation core, the unmollified
clock derivative as an $L^{2}(\R_{+})$-valued object, and the
second-order energy producing the Brownian-time PDE shadow --- remains
open.

We single out three levels of regularity for the IBM occupation
functional $F_T(\varphi)=\int_0^T\varphi(\IBM_s)\,ds$.

\begin{itemize}
\item[(A)] \emph{Mollified occupation core.}
  $F_{T,\varepsilon}(\varphi)$ in $\Dom(D_X)\cap\Dom(D_Y)$ for every
  $\varepsilon>0$, with $D_X F_{T,\varepsilon}$,
  $D_Y F_{T,\varepsilon}$ classical and
  $D_X D_Y F_{T,\varepsilon}$ classical to the extent the
  second-order source calculus allows. All theorem statements of
  this section operate at this level.
\item[(B)] \emph{Unmollified scenery derivative.}
  $D_X F_T(\varphi)$ is rigorous via the spatial factorisation
  $F_T(\varphi)=\int_{\R}\varphi(X_x)L_T^{\,x}(Y)\,dx$
  of~\cite{Fontes-IBMchaos}: it is an $L^{2}(\R)$-valued random
  variable for $\varphi\in C_c^\infty(\R)$.
\item[(C)] \emph{Unmollified full IBM derivative.} The candidate
  $D_\IBM F_T(\varphi)$, with components $D_X F_T(\varphi)$ and
  $D_Y F_T(\varphi)$, is \emph{not} constructed in this paper as an
  $L^{2}(\R)\oplus L^{2}(\R_{+})$-valued operator. The clock
  component is genuinely singular at the level of local time, and
  must be identified, in the unmollified limit, as a
  distribution-valued kernel or via chaos projections.
\end{itemize}

The remainder of this section makes (A) precise; (B) is cited as
background; (C) is the principal analytic target left open. The
scenery derivative is computed in closed form at the mollified level;
the clock derivative is computed at the mollified level and recorded
as singular at $\varepsilon=0$.

\subsection*{Test class and occupation core}

\paragraph{Convention for this section.}
Within Section~\ref{sec:ibm-occ}, $\IBM_s$ denotes the two-sided
composition $X_{Y_s}$ (rather than the canonical
$X_{|Y_s|}$ used elsewhere in the paper); this produces the
$\R$-valued spatial factorisation~\eqref{eq:occ-factor} with the
Brownian local time $L^{x}(Y)$ on $\R$, in line
with~\cite{Fontes-IBMchaos}. The conclusions below transfer to the
canonical $X_{|Y_s|}$ form by replacing $L^{x}(Y)$ on $\R$ with
$L^{x}(|Y|)$ on $\R_{+}$ throughout; we do not carry both notations
in parallel.

Let $\mathscr S_\IBM\coloneqq C_c^{\infty}(\R)$. For
$\varphi\in\mathscr S_\IBM$ and $T>0$, set
\[
  F_T(\varphi)\;\coloneqq\;\int_0^T\varphi(\IBM_s)\,ds.
\]
Since $\varphi$ is bounded and $T$ is finite,
$F_T(\varphi)\in L^{2}(\Omega,\mathcal F_\IBM,\Prob)$.

\begin{definition}[IBM occupation core]\label{def:cocc}
The \emph{IBM occupation core} is
\[
  \mathcal C_{\mathrm{occ}}
  \;\coloneqq\;\mathrm{span}\bigl\{\,F_T(\varphi)-\E F_T(\varphi):
   T>0,\ \varphi\in\mathscr S_\IBM\,\bigr\}\;\subseteq\;
   L^{2}(\Omega,\mathcal F_\IBM,\Prob).
\]
\end{definition}

$\mathcal C_{\mathrm{occ}}$ is a centred linear subspace of the
observable $L^{2}$. We do not claim density and we do not construct a
derivative on $\mathcal C_{\mathrm{occ}}$ directly: at the unmollified
level $F_T(\varphi)\notin\Dom(D_Y)$ as an
$L^{2}(\R_{+})$-valued Malliavin derivative, so
$\mathcal C_{\mathrm{occ}}\not\subseteq\mathcal C_\IBM$. We use
$\mathcal C_{\mathrm{occ}}$ as the unmollified \emph{target}
approximated by the mollified core
$\mathcal C_{\mathrm{occ}}^{\,\varepsilon}\subseteq\mathcal C_\IBM$ of
\cref{def:mollified}.

\subsection*{Spatial factorisation}

The Kac--L\'evy occupation density for the Brownian clock $Y$ gives
\begin{equation}\label{eq:occ-factor}
  F_T(\varphi)
  \;=\;\int_{\R}\varphi(X_x)\,L_T^{\,x}(Y)\,dx,
\end{equation}
where $L_T^{\,x}(Y)$ is the local time of $Y$ at level $x$ up to time
$T$. Identity~\eqref{eq:occ-factor} is standard
(see~\cite{RevuzYor,BorodinSalminen}) and is the structural input
used in~\cite{Fontes-IBMchaos}; we take it as given.

\subsection*{Scenery-mollified occupation core}

The spatial derivative of the unmollified scenery $X$ is distributional
(spatial white noise); we replace $X$ by its mollification, which keeps
both source derivatives classical at every $\varepsilon>0$.

\begin{definition}[Mollified scenery and mollified core]
\label{def:mollified}
Fix a non-negative even mollifier $\rho\in C_c^\infty(\R)$ with
$\int\rho=1$ and write $\rho_\varepsilon(z)=\varepsilon^{-1}\rho(z/\varepsilon)$.
The $\varepsilon$-mollified scenery is
\[
  X^{\varepsilon}_y \;\coloneqq\; (\rho_\varepsilon*X)(y)
  \;=\;\int_{\R}\rho_\varepsilon(y-z)\,X_z\,dz,
  \qquad y\in\R,
\]
which is $\Prob$-a.s.\ $C^{\infty}$ in $y$ on every compact, with
$X^{\varepsilon}_y\to X_y$ locally uniformly in $y$ a.s.\ as
$\varepsilon\to 0$ by Brownian H\"older continuity. The
$\varepsilon$-mollified occupation functional is
\[
  F_{T,\varepsilon}(\varphi)
  \;\coloneqq\;\int_0^T\varphi\bigl(X^{\varepsilon}_{Y_s}\bigr)\,ds,
\]
and the associated mollified core is
\[
  \mathcal C_{\mathrm{occ}}^{\,\varepsilon}
  \;\coloneqq\;\mathrm{span}\bigl\{\,F_{T,\varepsilon}(\varphi)-
      \E F_{T,\varepsilon}(\varphi):T>0,\ \varphi\in\mathscr S_\IBM\,\bigr\}.
\]
\end{definition}

\begin{lemma}[Source regularity of the mollified core]
\label{lem:moll-domain}
For every $\varepsilon>0$, $T>0$ and $\varphi\in\mathscr S_\IBM$,
\[
  F_{T,\varepsilon}(\varphi)\;\in\;\Dom(D_X)\cap\Dom(D_Y),
\]
and hence $\mathcal C_{\mathrm{occ}}^{\,\varepsilon}\subseteq\mathcal C_\IBM$.
\end{lemma}

\begin{proof}
For every $\omega$ and $\varepsilon>0$, $y\mapsto X^\varepsilon_y(\omega)$
is $C^\infty$ with derivatives controlled, on every compact, by
$\sup_K|X|$ and $\|\rho_\varepsilon^{(k)}\|_{L^1}$. Hence
$s\mapsto\varphi(X^\varepsilon_{Y_s(\omega)})$ is the composition of
$C^\infty$ maps integrated over $s\in[0,T]$, and the source chain rule
on the joint Wiener space~\cite[\S1.2]{Nualart} together with Fubini
gives $F_{T,\varepsilon}(\varphi)\in\Dom(D_X)\cap\Dom(D_Y)$. Explicit
kernels are given in \cref{prop:moll-deriv} below; their finiteness in
$L^{2}(\R)$ and $L^{2}(\R_{+})$ respectively follows from $\varphi$
having compact support and $X^\varepsilon,\partial_yX^\varepsilon$ having
finite $L^{2}$-moments on the range of $Y|_{[0,T]}$.
\end{proof}

\subsection*{Explicit source derivatives on the mollified core}

Let $\sigma_x\in L^{2}(\R)$ denote the Cameron--Martin representer of
the scenery evaluation $X\mapsto X_x$; in the two-sided Brownian
parametrisation $\sigma_x=\mathbf 1_{[0,x]}$ for $x\ge 0$ and
$\sigma_x=-\mathbf 1_{[x,0]}$ for $x<0$, so
$\|\sigma_x\|_{L^{2}(\R)}^{2}=|x|$. For the clock,
$D_rY_s=\mathbf 1_{r\le s}\in L^{2}(\R_{+})$.

\begin{proposition}[Source derivatives of the mollified core]
\label{prop:moll-deriv}
Set $(D_XX^\varepsilon_y)(\xi)\coloneqq(\rho_\varepsilon*\sigma_\cdot(\xi))(y)$
and $(\partial_yX^\varepsilon)(y)\coloneqq(\rho_\varepsilon'*X)(y)$. Then
for $\varphi\in\mathscr S_\IBM$, $T>0$, $\varepsilon>0$,
\begin{align}
  D_X F_{T,\varepsilon}(\varphi)(\xi)
  &\;=\;\int_0^T\varphi'\bigl(X^\varepsilon_{Y_s}\bigr)\,
         (D_XX^\varepsilon_{Y_s})(\xi)\,ds,
         \qquad\xi\in\R,
         \label{eq:DX-moll}\\[2pt]
  D_Y F_{T,\varepsilon}(\varphi)(r)
  &\;=\;\int_r^T\varphi'\bigl(X^\varepsilon_{Y_s}\bigr)\,
         (\partial_yX^\varepsilon)(Y_s)\,ds,
         \qquad r\in[0,T].
         \label{eq:DY-moll}
\end{align}
\end{proposition}

\begin{proof}
\eqref{eq:DX-moll}: $X^\varepsilon_y$ is linear in $X$ with
$D_XX^\varepsilon_y(\xi)=\int_\R\rho_\varepsilon(y-z)\sigma_z(\xi)\,dz$.
The source chain rule applied to $\varphi(X^\varepsilon_{Y_s})$ in the
$X$-direction yields $\varphi'(X^\varepsilon_{Y_s})\,D_XX^\varepsilon_{Y_s}(\xi)$;
integration over $s$ commutes by Fubini.

\eqref{eq:DY-moll}: $Y_s$ enters only through $y\mapsto X^\varepsilon_y$
evaluated at $Y_s$. Since $D_rY_s=\mathbf 1_{r\le s}$ and $X^\varepsilon$
has classical spatial derivative
$\partial_yX^\varepsilon=\rho_\varepsilon'*X$, the chain rule gives
$D_r\varphi(X^\varepsilon_{Y_s})=\varphi'(X^\varepsilon_{Y_s})\,
\partial_yX^\varepsilon(Y_s)\,\mathbf 1_{r\le s}$, and integration over
$s$ yields \eqref{eq:DY-moll}.
\end{proof}

\subsection*{Occupation observable energy}

\begin{definition}[Occupation observable energy]\label{def:occ-energy}
The \emph{$\varepsilon$-occupation observable energy} is
\[
  \Energy_\IBM^{\mathrm{occ},\,\varepsilon}
  \;\coloneqq\;
  \overline{\bigl\{(D_XF,\,D_YF):F\in\mathcal C_{\mathrm{occ}}^{\,\varepsilon}\bigr\}}
  ^{\,L^{2}(\R)\oplus L^{2}(\R_{+})}
  \;\subseteq\;\Energy_\IBM.
\]
\end{definition}

We do not claim that
$\overline{\bigcup_{\varepsilon>0}\Energy_\IBM^{\mathrm{occ},\,\varepsilon}}
=\Energy_\IBM$, nor that the energies stabilise as $\varepsilon\to 0$:
identifying the limiting object is part of the analytic programme
listed at the end of this section.

\subsection*{Occupation-core intrinsic derivative and divergence}

Specialising \cref{def:compression} to the occupation core:

\begin{definition}[Occupation derivative and divergence]
\label{def:ibm-occ-D}
For $F\in\mathcal C_{\mathrm{occ}}^{\,\varepsilon}$, the
\emph{occupation-core intrinsic derivative} is
\[
  D_\IBM^{\,\mathrm{occ},\,\varepsilon} F
  \;\coloneqq\;(D_X F,\,D_Y F)\;\in\;\Energy_\IBM^{\mathrm{occ},\,\varepsilon},
\]
explicit from \eqref{eq:DX-moll}--\eqref{eq:DY-moll}. Its adjoint
$\delta_\IBM^{\,\mathrm{occ},\,\varepsilon}=
(D_\IBM^{\,\mathrm{occ},\,\varepsilon})^{*}$
satisfies, on the lifted source-divergence domain,
\begin{equation}\label{eq:occ-delta}
  \delta_\IBM^{\,\mathrm{occ},\,\varepsilon} u
  \;=\;\E\!\bigl[\,\delta_X u_X+\delta_Y u_Y\,\big|\,\mathcal F_\IBM\,\bigr]
  \qquad\text{for }u=(u_X,u_Y)
  \text{ with }u_X\in\Dom(\delta_X),\ u_Y\in\Dom(\delta_Y).
\end{equation}
\end{definition}

\subsection*{Quantitative mollifier estimates}

The convergence theorem below relies on a single explicit second-moment
bound for the scenery mollifier.

\begin{lemma}[Scenery mollifier $L^{2}$-rate, uniform in $x$]
\label{lem:moll-rate}
Let $X$ be the two-sided Brownian scenery and $X^{\varepsilon}=
\rho_\varepsilon*X$ as in \cref{def:mollified}. Then \emph{uniformly}
for every $x\in\R$,
\begin{equation}\label{eq:moll-rate-L2}
  \E\bigl[(X^{\varepsilon}_x-X_x)^{2}\bigr]\;\le\;C_\rho\,\varepsilon,
  \qquad C_\rho\;\coloneqq\;\int_\R|v|\rho(v)\,dv\;<\;\infty,
\end{equation}
and consequently $\E|X^\varepsilon_x-X_x|\le C_\rho^{1/2}\varepsilon^{1/2}$
uniformly in $x$. In particular $X^\varepsilon_0\to 0$ in $L^{2}$ as
$\varepsilon\to 0$, since $X_0=0$ a.s.
\end{lemma}

\begin{proof}
For two-sided Brownian scenery, $\E[(X_z-X_x)^{2}]=|z-x|$ for every
$z,x\in\R$: when $z,x$ have the same sign or one is zero, this is the
standard one-sided Brownian increment variance; when $z>0>x$,
independence of $X_z$ and $X_x$ gives
$\E[(X_z-X_x)^{2}]=\E X_z^{2}+\E X_x^{2}=z+|x|=z-x=|z-x|$ (and
symmetrically for $x>0>z$). Hence the identity holds across the
origin.

Writing $X^\varepsilon_x-X_x=\int_\R\rho_\varepsilon(x-z)(X_z-X_x)\,dz$
and applying Jensen's inequality with the probability measure
$\rho_\varepsilon(x-z)\,dz$,
\[
  \E[(X^\varepsilon_x-X_x)^{2}]
  \;\le\;\int_\R\rho_\varepsilon(x-z)\,\E[(X_z-X_x)^{2}]\,dz
  \;=\;\int_\R\rho_\varepsilon(x-z)|z-x|\,dz.
\]
Substituting $v=(x-z)/\varepsilon$ gives
$\int_\R\rho(v)|v|\varepsilon\,dv=C_\rho\varepsilon$, proving
\eqref{eq:moll-rate-L2}. The bound is uniform in $x\in\R$ since the
right-hand side does not depend on $x$. The $L^{1}$ corollary is
Cauchy--Schwarz, and the $L^{2}$ convergence at $x=0$ is the case
$X_0=0$.
\end{proof}

\begin{lemma}[Scenery-derivative kernel: $L^{2}$ bound and rate]\label{lem:moll-H}
With $H_\varepsilon(y,\xi)\coloneqq(D_XX^\varepsilon_y)(\xi)=(\rho_\varepsilon*\sigma_\cdot(\xi))(y)$:
\begin{enumerate}
\item[(a)] $\|H_\varepsilon(y,\cdot)\|_{L^{2}(\R)}^{2}\le|y|+C_\rho\varepsilon$
  for every $y\in\R$;
\item[(b)] $\|H_\varepsilon(y,\cdot)-\sigma_y\|_{L^{2}(\R)}^{2}\le C'_\rho\varepsilon$
  for every $y\in\R$, with $C'_\rho$ depending only on the support of
  $\rho$.
\end{enumerate}
\end{lemma}

\begin{proof}
(a). Jensen and $\|\sigma_z\|_{L^{2}(\R)}^{2}=|z|$ give
$\int_\R H_\varepsilon(y,\xi)^{2}\,d\xi\le
\int_\R\rho_\varepsilon(y-z)|z|\,dz\le|y|+C_\rho\varepsilon$.

(b). The difference $H_\varepsilon(y,\xi)-\sigma_y(\xi)$ vanishes off the
$O(\varepsilon)$-neighbourhood of the jump locus
$\{\xi\in\R:\xi\,\text{between}\,y\,\text{and a point in supp}\,\rho_\varepsilon(y-\cdot)\}$,
and is uniformly bounded by $1$ there. The $L^{2}(\R,d\xi)$ mass is
therefore $O(\varepsilon)$ uniformly in $y$.
\end{proof}

\begin{lemma}[Spatial derivative of the mollified scenery]\label{lem:moll-eps-inv}
Away from the origin (specifically, for $|y|>c\varepsilon$ where
$\mathrm{supp}\,\rho\subseteq[-c,c]$),
$(\partial_y X^\varepsilon)(y)=(\rho_\varepsilon'*X)(y)$ is centred
Gaussian with
\[
  \Var\bigl((\partial_y X^\varepsilon)(y)\bigr)
  \;=\;\frac{c_\rho^{\,(1)}}{\varepsilon},
  \qquad c_\rho^{\,(1)}\;\coloneqq\;\int_\R\rho(v)^{2}\,dv.
\]
\end{lemma}

\begin{proof}
For $|y|>c\varepsilon$, the support of $z\mapsto\rho_\varepsilon(y-z)$ is
entirely on one side of $0$. Integration by parts (the boundary terms
vanish because $\rho_\varepsilon$ has compact support) gives
$(\partial_y X^\varepsilon)(y)=\int_\R\rho_\varepsilon(y-z)\,dX_z$, an
It\^o integral of a deterministic $L^{2}$ integrand. Hence the
$L^{2}$-isometry yields
$\Var=\int\rho_\varepsilon^{2}(u)\,du=\varepsilon^{-1}\int\rho^{2}$.
\end{proof}

\subsection*{Mollified-to-unmollified convergence in the scenery direction}

\begin{theorem}[Mollified-to-unmollified IBM occupation convergence in the scenery direction]
\label{thm:ibm-occ}
Let $T>0$ and $\varphi\in\mathscr S_\IBM=C_c^\infty(\R)$. Then:
\begin{enumerate}
\item For every $\varepsilon>0$,
  $\mathcal C_{\mathrm{occ}}^{\,\varepsilon}\subseteq\mathcal C_\IBM$,
  and the occupation-core intrinsic derivative
  $D_\IBM^{\,\mathrm{occ},\,\varepsilon}$ is the restriction of the
  minimal observable derivative $D_\IBM$ of \cref{def:ibm-datum} to
  $\mathcal C_{\mathrm{occ}}^{\,\varepsilon}$, with explicit kernels
  \eqref{eq:DX-moll}--\eqref{eq:DY-moll}.
\item On the lifted source-divergence domain, the adjoint
  $\delta_\IBM^{\,\mathrm{occ},\,\varepsilon}$ admits the projected
  formula \eqref{eq:occ-delta}.
\item As $\varepsilon\to 0$,
  $F_{T,\varepsilon}(\varphi)\to F_T(\varphi)$ in $L^{2}(\Omega)$, with
  rate $O(\varepsilon^{1/2})$.
\item As $\varepsilon\to 0$,
  $D_X F_{T,\varepsilon}(\varphi)\to D_X F_T(\varphi)$ in
  $L^{2}(\Omega;L^{2}(\R))$, also at rate $O(\varepsilon^{1/2})$, with
  limit
  \begin{equation}\label{eq:DX-limit}
    D_X F_T(\varphi)(\xi)
    \;=\;\int_{\R}\varphi'(X_x)\,\sigma_x(\xi)\,L_T^{\,x}(Y)\,dx
    \;=\;\int_0^T\varphi'(X_{Y_s})\,\sigma_{Y_s}(\xi)\,ds.
  \end{equation}
\item For every $\varepsilon>0$, $D_Y F_{T,\varepsilon}(\varphi)$ is a
  classical element of $L^{2}(\R_{+})$ given by \eqref{eq:DY-moll}. The
  underlying spatial kernel $\partial_yX^\varepsilon$ has pointwise
  variance $c_\rho^{(1)}\varepsilon^{-1}$ (\cref{lem:moll-eps-inv}), so
  the family $\{D_YF_{T,\varepsilon}(\varphi)\}_{\varepsilon>0}$ is not
  controlled by the uniform-in-$y$ estimates used for the scenery
  component. We do not assert convergence of
  $D_Y F_{T,\varepsilon}(\varphi)$ in $L^{2}(\Omega;L^{2}(\R_{+}))$.
  Any limiting clock object must be interpreted distributionally or
  chaos-by-chaos.
\end{enumerate}
\end{theorem}

\begin{proof}
(i) follows from \cref{lem:moll-domain,prop:moll-deriv}: both $D_\IBM$
and $D_\IBM^{\,\mathrm{occ},\,\varepsilon}$ act as
$F\mapsto D_\Source(\iota F)$ on
$\mathcal C_{\mathrm{occ}}^{\,\varepsilon}\subseteq\mathcal C_\IBM$.
(ii) is \cref{thm:proj-div} applied to
$D_\IBM^{\,\mathrm{occ},\,\varepsilon}$.

\emph{Localisation.} For $R>0$ set
$A_R\coloneqq\{\sup_{s\in[0,T]}|Y_s|\le R\}$. The Brownian reflection
principle gives
\begin{equation}\label{eq:bm-max}
  \Prob(A_R^{c})\;\le\;2\Prob(|Y_T|>R)\;\le\;2\,e^{-R^{2}/(2T)}.
\end{equation}
Both \cref{lem:moll-rate} and \cref{lem:moll-H}\,(b) are uniform in the
spatial variable $y\in\R$; we use the localisation $A_R/A_R^{c}$ to
keep the moment-dependent factor $\|H_\varepsilon\|^{2}\le|y|+C_\rho\varepsilon$
controlled, then send $R\to\infty$ if desired. The
$O(\varepsilon^{1/2})$ rates below hold for every fixed $R>0$.

(iii). By Lipschitz continuity of $\varphi$ and Cauchy--Schwarz,
\[
  |F_{T,\varepsilon}(\varphi)-F_T(\varphi)|^{2}
  \;\le\;\|\varphi'\|_\infty^{2}\,T
   \int_0^T(X^\varepsilon_{Y_s}-X_{Y_s})^{2}\,ds.
\]
\emph{On $A_R$:} $|Y_s|\le R$ for $s\in[0,T]$; conditioning on $Y$ and
applying \cref{lem:moll-rate} (uniform in $y\in[-R,R]$, indeed uniform
in $y\in\R$) gives
$\E_X[(X^\varepsilon_{Y_s}-X_{Y_s})^{2}\mathbf 1_{A_R}]\le C_\rho\varepsilon$.
Hence
$\E[|F_{T,\varepsilon}-F_T|^{2}\mathbf 1_{A_R}]\le C_\rho T^{2}\|\varphi'\|_\infty^{2}\varepsilon$.

\emph{On $A_R^{c}$:} The same Lemma~\ref{lem:moll-rate} bound holds for
all $y\in\R$, so
$\E_X[(X^\varepsilon_{Y_s}-X_{Y_s})^{2}\mathbf 1_{A_R^{c}}]\le C_\rho\varepsilon\,
\mathbf 1_{A_R^{c}}$, giving
$\E[|F_{T,\varepsilon}-F_T|^{2}\mathbf 1_{A_R^{c}}]\le
C_\rho T^{2}\|\varphi'\|_\infty^{2}\varepsilon\,\Prob(A_R^{c})$.

Summing, $\E[|F_{T,\varepsilon}-F_T|^{2}]\le
C_\rho T^{2}\|\varphi'\|_\infty^{2}\varepsilon\,(1+\Prob(A_R^{c}))\le
2C_\rho T^{2}\|\varphi'\|_\infty^{2}\varepsilon$, the claimed
$O(\varepsilon^{1/2})$ rate in $L^{2}(\Omega)$. The estimate is uniform
in $R>0$; the localisation enters only to expose where the
unbounded range of $Y_s$ would otherwise need to be controlled.

(iv). The second equality in \eqref{eq:DX-limit} is the occupation
density formula applied to $x\mapsto\varphi'(X_x)\sigma_x(\xi)$. For
the limit itself, set $H_\varepsilon(y,\xi)\coloneqq(D_XX^\varepsilon_y)(\xi)$
and decompose
\[
  D_X F_{T,\varepsilon}(\varphi)-D_X F_T(\varphi)\;=\;A_\varepsilon+B_\varepsilon,
\]
where
\[
  A_\varepsilon(\xi)
  =\!\int_0^T\!\bigl[\varphi'(X^\varepsilon_{Y_s})-\varphi'(X_{Y_s})\bigr]
   H_\varepsilon(Y_s,\xi)\,ds,
  \quad
  B_\varepsilon(\xi)
  =\!\int_0^T\!\varphi'(X_{Y_s})\bigl[H_\varepsilon(Y_s,\xi)-\sigma_{Y_s}(\xi)\bigr]ds.
\]

\emph{Bound on $A_\varepsilon$, with explicit localisation.} By
Cauchy--Schwarz in $s$ and in $\omega$,
\[
  \|A_\varepsilon\|_{L^{2}(\Omega;L^{2}(\R))}^{2}
  \le T\!\int_0^T\!\E\Bigl[
       |\varphi'(X^\varepsilon_{Y_s})-\varphi'(X_{Y_s})|^{2}\,
       \|H_\varepsilon(Y_s,\cdot)\|_{L^{2}(\R)}^{2}\Bigr]ds.
\]
Conditioning on $Y$, Lipschitz $\varphi'$ and \cref{lem:moll-rate} give
the uniform-in-$y$ bound
$\E_X[|\varphi'(X^\varepsilon_y)-\varphi'(X_y)|^{2}]\le
\|\varphi''\|_\infty^{2}C_\rho\varepsilon$. Split the remaining factor
on $A_R$ and $A_R^{c}$.

On $A_R$, $|Y_s|\le R$, so \cref{lem:moll-H}\,(a) gives
$\|H_\varepsilon(Y_s,\cdot)\|^{2}\le R+C_\rho\varepsilon$ uniformly for
$s\in[0,T]$, hence
\[
  \E\bigl[|\varphi'(X^\varepsilon_{Y_s})-\varphi'(X_{Y_s})|^{2}
   \|H_\varepsilon(Y_s,\cdot)\|^{2}\mathbf 1_{A_R}\bigr]
  \le\|\varphi''\|_\infty^{2}C_\rho(R+C_\rho\varepsilon)\,\varepsilon.
\]
On $A_R^{c}$, the global bound $\|H_\varepsilon(Y_s,\cdot)\|^{2}\le
|Y_s|+C_\rho\varepsilon$ combined with Cauchy--Schwarz and
\eqref{eq:bm-max} gives
\[
  \E[(|Y_s|+C_\rho\varepsilon)\mathbf 1_{A_R^{c}}]
  \le\E[Y_s^{2}]^{1/2}\Prob(A_R^{c})^{1/2}+C_\rho\varepsilon\Prob(A_R^{c})
  \le\sqrt{2T}\,e^{-R^{2}/(4T)}+2C_\rho\varepsilon\,e^{-R^{2}/(2T)}.
\]
Multiplying by the $\varepsilon$-factor from $\varphi'$ Lipschitz, the
$A_R^{c}$ contribution is $O(\varepsilon\,e^{-R^{2}/(4T)})$ for any
$R>0$. Adding the two pieces, for any fixed $R>0$,
\[
  \|A_\varepsilon\|_{L^{2}(\Omega;L^{2}(\R))}^{2}
  \;\le\;T^{2}\|\varphi''\|_\infty^{2}C_\rho\,
   \bigl[(R+C_\rho\varepsilon)\varepsilon
       +\sqrt{2T}\,\varepsilon\,e^{-R^{2}/(4T)}\bigr]
  \;=\;O(\varepsilon).
\]

\emph{Bound on $B_\varepsilon$.} \cref{lem:moll-H}\,(b) is uniform in
$y\in\R$, so no localisation is needed for this term:
\[
  \|B_\varepsilon\|_{L^{2}(\Omega;L^{2}(\R))}^{2}
  \le\|\varphi'\|_\infty^{2}T\!\int_0^T\!
       \E\|H_\varepsilon(Y_s,\cdot)-\sigma_{Y_s}\|_{L^{2}(\R)}^{2}\,ds
  \le C'_\rho\,\|\varphi'\|_\infty^{2}\,T^{2}\,\varepsilon.
\]
Adding the two contributions,
$\|D_X F_{T,\varepsilon}(\varphi)-D_X F_T(\varphi)\|_{L^{2}(\Omega;L^{2}(\R))}
=O(\varepsilon^{1/2})$.

(v). For each $\varepsilon>0$, $\partial_yX^\varepsilon=\rho_\varepsilon'*X$
is a centred Gaussian random function of $y$, $\Prob$-a.s.\ continuous
on every compact, and $D_Y F_{T,\varepsilon}(\varphi)$ defined by
\eqref{eq:DY-moll} is a well-defined element of $L^{2}(\R_{+})$ by
\cref{lem:moll-domain}. \Cref{lem:moll-eps-inv} gives the pointwise
variance $\Var(\partial_yX^\varepsilon(y))=c_\rho^{(1)}\varepsilon^{-1}$
for $|y|>c\varepsilon$. The variance scaling of $\partial_yX^\varepsilon$
is the analytic obstruction; the uniform-in-$y$ control used for the
scenery component fails for the clock kernel. Establishing sharp
divergence or renormalised convergence of
$D_Y F_{T,\varepsilon}(\varphi)$ requires a careful account of
correlations along the clock path, possible zeros of $\varphi'$, and
clock revisits, and belongs to the next-stage clock analysis. We assert
neither convergence nor divergence in $L^{2}(\Omega;L^{2}(\R_{+}))$ here.
\end{proof}

\begin{corollary}[Clock derivative: the first analytic obstruction]
\label{cor:unmollified}
\Cref{thm:ibm-occ}~(iii)--(iv) shows that the
observable-compression construction is not formal: on the IBM
occupation core, the mollified intrinsic derivative converges in the
observable functional ($F_{T,\varepsilon}\to F_T$ in $L^{2}(\Omega)$)
and in the scenery-energy direction ($D_X F_{T,\varepsilon}\to D_X F_T$
in $L^{2}(\Omega;L^{2}(\R))$), with closed-form limit
\eqref{eq:DX-limit} via the spatial factorisation. The clock direction
sits outside this picture: the spatial kernel
$\partial_yX^\varepsilon=\rho_\varepsilon'*X$ has pointwise variance
$c_\rho^{(1)}\varepsilon^{-1}$ (\cref{lem:moll-eps-inv}), so no $L^{2}$
limit is asserted here; the pointwise variance scaling indicates that
the clock component cannot be handled by the same classical
compactness argument as the scenery component. Any limiting clock
object must be interpreted distributionally or chaos-by-chaos --- as a
distribution-valued kernel against $L^{2}(\R_{+})$ test integrands, or
through source-chaos projections in the bidegree atlas
of~\cite{Fontes-IBMchaos,NualartVives}. This is the first analytic
obstruction to a fully classical intrinsic IBM derivative and the
natural entry point for the second-order observable energy of
Section~\ref{sec:ibm-source-2nd}.
\end{corollary}

\subsection*{Why this matters}

\Cref{thm:ibm-occ}~(iii)--(iv) is the first piece of evidence that the
abstract observable-compression construction is realised concretely on
a non-trivial subspace of $\mathcal C_\IBM$, with the unmollified
scenery component surviving the $\varepsilon\to 0$ limit in the
expected $L^{2}$-sense. \Cref{cor:unmollified} isolates the remaining
analytic gap as a single, structural feature of the composite
$X\circ Y$: the clock direction probes the spatial white noise of the
scenery, an object that lives outside any classical $L^{2}$-valued
Malliavin theory. This is exactly where the source-chaos atlas
of~\cite{Fontes-IBMchaos} and the second-order observable energy of
Section~\ref{sec:ibm-source-2nd} must enter, and is the natural
analytic obstruction the next theorem in this programme has to
resolve.

\subsection*{Connection to source bidegrees}

\begin{remark}[Bidegree consistency]\label{rem:bidegree}
The lifted occupation derivative
$J D_\IBM^{\,\mathrm{occ},\,\varepsilon} F$ sits inside
$L^{2}(\R)\oplus L^{2}(\R_{+})$ and decomposes against the joint Wiener
bidegree atlas
$\bigoplus_{m,n\ge 0}\Hilb^{X}_{m}\otimes\Hilb^{Y}_{n}$
of~\cite{Fontes-IBMchaos}. Mode by mode:
\begin{itemize}
\item bidegrees $(m,0)$, $m\ge 1$, recover pure scenery contributions
  driven by \eqref{eq:DX-moll};
\item bidegrees $(0,n)$, $n\ge 1$, recover pure clock contributions
  driven by \eqref{eq:DY-moll};
\item mixed bidegrees $(m,n)$, $m,n\ge 1$, are
  \emph{consistency constraints} linking $D_X F$ and $D_Y F$ through
  the spatial factorisation~\eqref{eq:occ-factor}.
\end{itemize}
We do not reprove the bidegree expansion; it is the content
of~\cite{Fontes-IBMchaos} and is cited as background.
\end{remark}

\paragraph{Transition: two distinct IBM functionals.}
The occupation-core construction of this section concerns the pathwise
functional $F_T(\varphi)=\int_0^{T}\varphi(X_{Y_s})\,ds$. The Fourier-mode
calculation below concerns the Brownian-time semigroup observable
$u(t,x)=\E[f(x+X_{|Y_t|})]$. These are related but distinct: the
former supplies a concrete observable-energy core (\cref{thm:ibm-occ});
the latter identifies the PDE shadow expected at the second-order
level. The calculation that follows is \emph{not} derived from the
first-order observable compression; it is a separate diagnostic of
what the second-order observable energy of
Section~\ref{sec:second-order} must reproduce.

\subsection*{Fourier-mode PDE-shadow diagnostic}

We record a PDE-shadow diagnostic identifying, at the level of
Fourier-mode test functions, the second-order signature that any
intrinsic IBM It\^o formula would have to reproduce. \emph{We do not
derive the intrinsic IBM It\^o formula here; the Fourier computation
identifies what the second-order observable energy must reproduce in
the semigroup setting.}

\paragraph{Convention.} $X=(X_s)_{s\in\R}$ is the standard two-sided
Brownian scenery ($\Var(X_s)=|s|$, generator $\tfrac12\Delta$), and
$Y=(Y_t)_{t\ge 0}$ is an independent standard one-sided Brownian clock
from $Y_0=0$. Set
\[
  Z_t\;\coloneqq\;x+X_{|Y_t|},\qquad t\ge 0,
\]
the Brownian-time process started at $x$. The use of $|Y_t|$ here (in
place of $Y_t$ on its own filtration) is for clean PDE benchmarking;
the abstract observable-compression construction does not depend on
this convention.

\paragraph{Characteristic function.}
For $k\in\R$ and $t>0$,
\[
  u_k(t,x)\;\coloneqq\;\E\bigl[e^{ikZ_t}\bigr]
  \;=\;e^{ikx}\,\E\bigl[e^{ikX_{|Y_t|}}\bigr].
\]
Conditioning on $|Y_t|=s$ gives $\E[e^{ikX_s}]=e^{-k^{2}s/2}$, and
$|Y_t|$ has half-normal density
$p_t(s)=\sqrt{2/(\pi t)}\,e^{-s^{2}/(2t)}$ for $s\ge 0$. Hence
\begin{equation}\label{eq:uk-int}
  u_k(t,x)
  \;=\;e^{ikx}\,\sqrt{\tfrac{2}{\pi t}}
   \int_0^\infty e^{-k^{2}s/2-s^{2}/(2t)}\,ds.
\end{equation}
Completing the square,
$-k^{2}s/2-s^{2}/(2t)=-(s+k^{2}t/2)^{2}/(2t)+k^{4}t/8$, and
substituting $u=(s+k^{2}t/2)/\sqrt{t}$,
\begin{equation}\label{eq:uk-closed}
  u_k(t,x)
  \;=\;2\,e^{ikx}\,e^{k^{4}t/8}\,\Phi^{c}\!\bigl(k^{2}\sqrt{t}/2\bigr),
  \qquad
  \Phi^{c}(a)\coloneqq\!\int_a^\infty\!\frac{e^{-u^{2}/2}}{\sqrt{2\pi}}\,du.
\end{equation}
At $t\to 0^{+}$, $\Phi^{c}(0)=1/2$ recovers $u_k(0^{+},x)=e^{ikx}=f_k(x)$.

\paragraph{PDE identity at the Fourier-mode level.}
Differentiate \eqref{eq:uk-closed} in $t$. With $a=k^{2}\sqrt t/2$
and $b=k^{4}t/8$, $(\Phi^{c})'(a)=-\phi(a)$ for
$\phi(a)=(2\pi)^{-1/2}e^{-a^{2}/2}$, and the cancellation
$a^{2}/2=b$ gives $e^{b}\phi(a)=(2\pi)^{-1/2}$. Therefore
\begin{equation}\label{eq:uk-pde}
  \partial_t u_k(t,x)
  \;=\;e^{ikx}\Bigl[\,2\,e^{b}\cdot\tfrac{k^{4}}{8}\,\Phi^{c}(a)
       \;-\;2\,e^{b}\phi(a)\cdot\tfrac{k^{2}}{4\sqrt{t}}\Bigr]
  \;=\;\frac{k^{4}}{8}\,u_k(t,x)\;-\;\frac{k^{2}\,e^{ikx}}{2\sqrt{2\pi t}}.
\end{equation}
Translating $-k^{2}e^{ikx}=\Delta f_k$ and $k^{4}u_k=\Delta^{2}u_k$
(with $\Delta=\partial_x^{2}$),
\begin{equation}\label{eq:az-fourier}
  \partial_t u_k(t,x)
  \;=\;\frac{1}{8}\,\Delta^{2}u_k(t,x)
       \;+\;\frac{\Delta f_k(x)}{2\sqrt{2\pi t}}.
\end{equation}
Fourier synthesis on smooth bounded $f$ yields
\begin{equation}\label{eq:az-general}
  \partial_t u(t,x)
  \;=\;\frac{1}{8}\,\Delta^{2}u(t,x)
       \;+\;\frac{\Delta f(x)}{2\sqrt{2\pi t}},
  \qquad u(0,x)=f(x).
\end{equation}
Under the stated convention, \eqref{eq:az-general} is the Brownian-time
PDE of~\cite{AllouZheng}; the constants $1/8$ and $1/(2\sqrt{2\pi})$
are derived directly from \eqref{eq:uk-int}.

\paragraph{Convention sensitivity.} The bi-Laplacian coefficient $1/8$
and the boundary coefficient $1/(2\sqrt{2\pi})$ depend on the scenery
variance normalisation. With $\Var(X_s)=\alpha s$, the coefficient of
$\Delta^{2}u$ becomes $\alpha^{2}/8$ and that of $\Delta f$ becomes
$\alpha/(2\sqrt{2\pi t})$. We have used $\alpha=1$ (standard Brownian
motion); rescaling $X$ rescales both constants accordingly.

\paragraph{Interpretation.}
\Cref{eq:az-general} is \emph{not} a new PDE theorem and not the
intrinsic IBM It\^o formula. It identifies the second-order shadow the
energy-first programme must reproduce:
\begin{itemize}
\item the bi-Laplacian $\tfrac18\Delta^{2}u$ on the right is what a
  successful second-order observable energy
  $\Energy_\IBM^{(2)}$ (\cref{rem:bilap-where}) should generate as the
  $\Carre$-form contribution;
\item the singular $t^{-1/2}$ term $\Delta f/(2\sqrt{2\pi t})$ is the
  boundary/local-time signature attributable to $|Y_t|$ near the origin
  of the half-normal density.
\end{itemize}
The first-order observable compression of \cref{thm:ibm-occ}
constructs $D_\IBM$ on the mollified occupation core. The
Fourier-mode calculation \eqref{eq:uk-int}--\eqref{eq:az-general}
shows what the next layer --- the second-order observable energy ---
must reproduce: the $k^{4}$ multiplier corresponding to $\Delta^{2}$,
together with the singular $t^{-1/2}$ boundary term.

\paragraph{Two levels, one source structure.}
For pathwise occupation functionals $F_T(\varphi)$, mixed bidegrees
$(m,n)$ with $m,n\ge 1$ obstruct a closed scalar PDE
(\cref{rem:bidegree}). For Brownian-time semigroup observables
$u(t,x)=\E[f(x+X_{|Y_t|})]$, the averaged law admits the closed
fourth-order PDE \eqref{eq:az-general}. Both phenomena reflect the
same composite source structure $X\circ Y$ but at different levels:
the occupation-functional layer keeps the joint $(X,Y)$ filtration
explicit and resolves the mixed sector chaos-by-chaos
in~\cite{Fontes-IBMchaos}; the semigroup layer averages over $Y$ first
and exposes a closed PDE in $x$. The Fourier diagnostic above operates
strictly at the semigroup level.

\subsection*{What remains for the full IBM It\^o formula}

The occupation-core construction is the first non-trivial IBM
instantiation of the observable-compression framework: it produces a
closed first-order operator on a non-empty, explicit subdomain of
$\mathcal C_\IBM$, with closed-form source derivatives at the mollified
level. The gap between \cref{thm:ibm-occ} and a full intrinsic IBM
It\^o formula splits into four analytic problems.
\begin{enumerate}
\item Either prove density of $\mathcal C_{\mathrm{occ}}$ in
  $L^{2}(\mathcal F_\IBM)$, or commit to the smaller occupation-level
  observable calculus carried by the closure of
  $\mathcal C_{\mathrm{occ}}$.
\item Identify the unmollified clock derivative
  $D_Y F_T(\varphi)=\lim_{\varepsilon\to 0}D_Y F_{T,\varepsilon}(\varphi)$
  as a distribution-valued kernel against $L^{2}(\R_{+})$ test
  integrands, and prove closability of the resulting operator on
  $\mathcal C_{\mathrm{occ}}$.
\item Construct the second-order observable energy
  $\Energy_\IBM^{(2)}$ on (the closure of) the occupation core, in the
  sense of Section~\ref{sec:ibm-source-2nd}.
\item Verify that the second-order $\Carre$-form, averaged against the
  law of $Y$, reproduces the bi-Laplacian/local-time PDE signature
  identified by~\cite{AllouZheng}.
\end{enumerate}
The It\^o identity itself remains the target of \cref{conj:ibm-ito}.

\section{Second-order observable energy}\label{sec:second-order}

We now formalise the operator-energy object that a future intrinsic
IBM It\^o formula must use: a \emph{candidate second-order observable
energy}, built by compressing a second-order source derivative through
the observable filtration. The construction mirrors
Section~\ref{sec:compat} but at second order. We treat it as a
candidate datum, not a completed generator: \cref{thm:ibm-occ}
supplies the first-order operator $D_\IBM$ on the mollified occupation
core, and the Fourier diagnostic \eqref{eq:az-general} indicates that
any successful intrinsic IBM It\^o formula must access a $\Delta^{2}$
contribution plus a singular $t^{-1/2}$ boundary term. The present
section defines the abstract object expected to carry that
information.

\subsection*{Abstract definition}

Let $(\Hilb_\Source,\Energy_\Source,D_\Source)$ be the source energy
datum of Section~\ref{sec:compat}. Suppose $D_\Source^{(2)}$ is a closed
densely defined second-order source derivative
\[
  D_\Source^{(2)}\colon\Dom(D_\Source^{(2)})\subseteq\Hilb_\Source
   \longrightarrow\Energy_\Source^{(2)},
\]
mapping into some second-order source energy space
$\Energy_\Source^{(2)}$. In the Gaussian setting, $D_\Source^{(2)}$ is
the iteration $D_\Source\otimes D_\Source$ on the second source
Sobolev space, and $\Energy_\Source^{(2)}$ is the appropriate tensor
product; the choice is part of the datum.

\begin{definition}[Second-order observable smooth core]
\label{def:obs-core-2}
$\mathcal C_\Obs^{(2)}\coloneqq
 \bigl\{F\in\Hilb_\Obs:\iota F\in\Dom(D_\Source^{(2)})\bigr\}$.
\end{definition}

\begin{definition}[Candidate second-order compression operator and energy]
\label{def:compression-2}
For $F\in\mathcal C_\Obs^{(2)}$, set
$D_\Obs^{(2),0}F\coloneqq D_\Source^{(2)}(\iota F)
\in\Energy_\Source^{(2)}$.
The \emph{candidate second-order observable energy} is
\[
  \Energy_\Obs^{(2)}
  \;\coloneqq\;
  \overline{\bigl\{\,D_\Source^{(2)}(\iota F):F\in\mathcal C_\Obs^{(2)}\,\bigr\}}
   ^{\,\Energy_\Source^{(2)}}
  \;\subseteq\;\Energy_\Source^{(2)},
\]
with isometric inclusion
$J^{(2)}\colon\Energy_\Obs^{(2)}\hookrightarrow\Energy_\Source^{(2)}$.
\end{definition}

The next proposition is the second-order analogue of
\cref{thm:compression}.

\begin{proposition}[Second-order observable compression]
\label{prop:compression-2}
Assume $D_\Source^{(2)}$ is closed and densely defined and that
$\mathcal C_\Obs^{(2)}$ is dense in $\Hilb_\Obs$. Then
$D_\Obs^{(2),0}\colon\mathcal C_\Obs^{(2)}\subseteq\Hilb_\Obs
\to\Energy_\Obs^{(2)}$ is closed on the maximal pullback domain
$\mathcal C_\Obs^{(2)}$; we denote it $D_\Obs^{(2)}$. Any restriction
to a subcore $\mathcal C'\subseteq\mathcal C_\Obs^{(2)}$ dense in
$\Hilb_\Obs$ is closable, and the closure is contained in
$D_\Obs^{(2)}$.
\end{proposition}

\begin{proof}
Identical structurally to \cref{thm:compression}: if $F_n\in
\mathcal C_\Obs^{(2)}$ with $F_n\to F$ in $\Hilb_\Obs$ and
$D_\Obs^{(2),0}F_n\to h$ in $\Energy_\Obs^{(2)}$, the isometric
inclusions $\iota$ and $J^{(2)}$ give $\iota F_n\to\iota F$ in
$\Hilb_\Source$ and $J^{(2)}D_\Obs^{(2),0}F_n=D_\Source^{(2)}(\iota F_n)
\to J^{(2)}h$ in $\Energy_\Source^{(2)}$. Closedness of
$D_\Source^{(2)}$ yields $\iota F\in\Dom(D_\Source^{(2)})$, i.e.,
$F\in\mathcal C_\Obs^{(2)}$, and $D_\Obs^{(2),0}F=h$. Density of
$\mathcal C_\Obs^{(2)}$ in $\Hilb_\Obs$ ensures the adjoint is
densely defined; closability of subcore restrictions is the same
argument applied to sequences in $\mathcal C'$.
\end{proof}

\Cref{prop:compression-2} produces a closed operator; it does not
produce a generator. Identifying the action of $D_\Obs^{(2)}$ on a
useful test class and converting it into an It\^o-type correction is
the next analytic problem.

\subsection*{IBM specialisation}

For IBM with source $\Source=(X,Y)$, the natural second-order source
derivative is the four-tuple
\begin{equation}\label{eq:DS2-IBM}
  D_\Source^{(2)}F
  \;=\;\bigl(D_X^{\,2}F,\;D_XD_YF,\;D_YD_XF,\;D_Y^{\,2}F\bigr),
\end{equation}
acting on the joint Sobolev space, with values in the symmetrised tensor
$\Energy_\Source^{(2)}\subseteq L^{2}(\R^{2})\oplus L^{2}(\R\times\R_{+})
\oplus L^{2}(\R_{+}\times\R)\oplus L^{2}(\R_{+}^{2})$.
Each block is the iterated partial Malliavin derivative on the
corresponding Gaussian factor.

For the mollified occupation functional $F_{T,\varepsilon}(\varphi)$ of
\cref{def:mollified}, each block of \eqref{eq:DS2-IBM} is classical:
$X^\varepsilon$ is $C^{\infty}$ in its spatial argument, so all
relevant joint derivatives are well-defined. The clock-block components
$D_XD_Y$, $D_YD_X$, $D_Y^{\,2}$ inherit the singular spatial kernel
$\partial_yX^\varepsilon$ via the source chain rule and propagate its
$\varepsilon^{-1}$ pointwise-variance scaling (\cref{lem:moll-eps-inv})
into the unmollified limit; for fixed $\varepsilon>0$ they are
well-defined.

\begin{definition}[Candidate second-order IBM observable energy]
\label{def:second-order-ibm}
Let $\mathcal C_\IBM^{(2)}\coloneqq\{F\in L^{2}(\mathcal F_\IBM):
\iota F\in\Dom(D_\Source^{(2)})\}$. The \emph{candidate second-order
IBM observable energy} is
\[
  \Energy_\IBM^{(2)}
  \;\coloneqq\;\overline{\bigl\{D_\Source^{(2)}(\iota F):F\in\mathcal
   C_\IBM^{(2)}\bigr\}}^{\,\Energy_\Source^{(2)}}.
\]
The second-order IBM compression operator $D_\IBM^{(2)}$ is the closed
pullback of $D_\Source^{(2)}$ on $\mathcal C_\IBM^{(2)}$, provided the
latter is dense in $L^{2}(\mathcal F_\IBM)$.
\end{definition}

\Cref{def:second-order-ibm} is a candidate datum, not a generator: the
closure exists abstractly by \cref{prop:compression-2} once the density
hypothesis on $\mathcal C_\IBM^{(2)}$ is verified, but identifying
$\Energy_\IBM^{(2)}$ concretely and producing a closed-form action of
$D_\IBM^{(2)}$ on cylindrical IBM functionals are open.

\subsection*{Connection to source bidegrees}

\begin{remark}[Bidegree consistency at second order]\label{rem:bidegree-2}
The four blocks of \eqref{eq:DS2-IBM} act on the joint Wiener bidegree
atlas of~\cite{Fontes-IBMchaos} as follows: $D_X^{\,2}F$ sees the
second scenery chaos, in particular the $(m,n)$-component with
$m\ge 2$; $D_Y^{\,2}F$ sees the second clock chaos, $n\ge 2$;
$D_XD_YF$ and $D_YD_XF$ see the mixed scenery--clock interaction,
supported on $(m,n)$ with $m,n\ge 1$. The mixed-sector $(2,1)$,
$(1,2)$, and $(2,2)$ kernels computed in~\cite{Fontes-IBMchaos} thus
serve as \emph{consistency requirements} for a candidate
$D_\IBM^{(2)}$: it must reproduce those kernels in the lifted
$J^{(2)}$-view. We do not reprove the bidegree expansion; the
consistency requirement is one-way.
\end{remark}

\subsection*{Connection to the Fourier PDE shadow}

The diagnostic \eqref{eq:az-general} of Section~\ref{sec:ibm-occ}
requires that, in the semigroup-averaged view, the second-order
observable energy produce a $k^{4}$ multiplier (equivalently a
$\tfrac18\Delta^{2}u$ contribution) plus the singular boundary term
$\Delta f/(2\sqrt{2\pi t})$. The $\Carre^{(2)}$-form
\[
  \Carre^{(2)}_\IBM(F,G)
  \;\coloneqq\;\langle D_\IBM^{(2)}F,\,D_\IBM^{(2)}G\rangle_{\Energy_\IBM^{(2)}}
\]
built from \cref{def:second-order-ibm} is the natural candidate to
carry this information once averaged against the law of $Y$. The
present section defines the operator-energy object; producing the
closed action and matching the Fourier constants $1/8$ and
$1/(2\sqrt{2\pi})$ are the next analytic problems.

\subsection*{Scope discipline}

We do not derive the intrinsic IBM It\^o formula here, nor do we claim
the bi-Laplacian correction follows from $D_\IBM^{(2)}$ as constructed.
\Cref{def:second-order-ibm} supplies a candidate second-order
observable energy; \cite{Fontes-IBMchaos} supplies the consistency
requirements at each bidegree; \eqref{eq:az-general} supplies the
target Fourier signature. Closing the gap between these three is the
next analytic target.

\subsection*{Remaining analytic tasks at second order}

Extending the first-order list of Section~\ref{sec:tasks}:
\begin{enumerate}
\item Verify density of $\mathcal C_\IBM^{(2)}$ in
  $L^{2}(\mathcal F_\IBM)$, ideally on (the closure of) the mollified
  occupation core $\mathcal C_{\mathrm{occ}}^{\,\varepsilon}$.
\item Identify $\Energy_\IBM^{(2)}$ concretely as a closed subspace of
  $\Energy_\Source^{(2)}$ on a useful test class.
\item Match the bidegree projections of~\cite{Fontes-IBMchaos} at the
  $(2,0)$, $(1,1)$, $(0,2)$, $(2,1)$, $(1,2)$, $(2,2)$ sectors.
\item Show that the averaged $\Carre^{(2)}_\IBM$ reproduces the
  $\Delta^{2}$ and $t^{-1/2}$-boundary signatures of
  \eqref{eq:az-general}.
\item Combine with \cref{conj:ibm-ito} to identify the correction
  $C_\IBM(f;t)$.
\end{enumerate}

\section{An operator-It\^o transfer principle}\label{sec:transfer}
\label{sec:ibm-ito}

We layer the framework one step further. Combining the first-order
observable divergence $\delta_\Obs$ (\cref{thm:proj-div}) with the
second-order observable correction sketched in
Section~\ref{sec:second-order}, we state and prove an abstract
\emph{conditional} transfer theorem showing how a source-level It\^o
expansion descends to an observable-level one. We then specialise to
IBM, recap the Fourier-mode diagnostic, and read off what the
source-resolved chaos atlas says about the remaining pieces.

\subsection*{The conditional transfer theorem}

Let $(\Hilb_\Source,\Energy_\Source,D_\Source)$ and
$(\Hilb_\Obs,\Energy_\Obs,D_\Obs)$ be the source and observable energy
data of Section~\ref{sec:compat}, with isometric inclusions $\iota$ and
$J$ and observable projection $P_\Obs=\iota^*=\E[\,\cdot\mid\mathcal F_\Obs]$.
Suppose an observable process $U=(U_t)_{t\ge 0}$ and a smooth class of
test functions admit a \emph{source-level} It\^o-type expansion
\begin{equation}\label{eq:source-ito}
  f(U_t)-f(U_0)\;=\;\delta_\Source\bigl(A_f(t)\bigr)\;+\;R_f(t),
  \qquad A_f(t)\in\Dom(\delta_\Source),
\end{equation}
on the joint source $\sigma$-algebra $\mathcal F_\Source$, with
remainder $R_f(t)\in L^{2}(\mathcal F_\Source)$.

\begin{theorem}[Conditional operator-It\^o transfer]\label{thm:transfer}
Assume \eqref{eq:source-ito} holds with $A_f(t)\in J(\Energy_\Obs)$,
so that there exists $B_f(t)\in\Energy_\Obs$ with $JB_f(t)=A_f(t)$, and
assume $B_f(t)$ lies in the lifted source-divergence domain
$\{u\in\Energy_\Obs:Ju\in\Dom(\delta_\Source)\}$. Suppose further that
$f(U_t)-f(U_0)$ is $\Prob$-a.s.\ $\mathcal F_\Obs$-measurable. Then
\begin{equation}\label{eq:obs-ito}
  f(U_t)-f(U_0)\;=\;\delta_\Obs\bigl(B_f(t)\bigr)\;+\;P_\Obs R_f(t),
\end{equation}
with $\delta_\Obs(B_f(t))=P_\Obs\,\delta_\Source(A_f(t))$.
\end{theorem}

\begin{proof}
By \cref{thm:proj-div} applied to $B_f(t)$,
$\delta_\Obs(B_f(t))=P_\Obs\,\delta_\Source(JB_f(t))
=P_\Obs\,\delta_\Source(A_f(t))$. Apply $P_\Obs$ to both sides of
\eqref{eq:source-ito}: since $f(U_t)-f(U_0)$ is
$\mathcal F_\Obs$-measurable, $P_\Obs[f(U_t)-f(U_0)]=f(U_t)-f(U_0)$;
linearity of $P_\Obs$ on the right gives
$\delta_\Obs(B_f(t))+P_\Obs R_f(t)$.
\end{proof}

\Cref{thm:transfer} does \emph{not} prove an It\^o formula for $U$ on
its own. It converts a source-level identity \eqref{eq:source-ito}
into an observable-level identity \eqref{eq:obs-ito}, conditional on
the source divergence integrand falling in the lifted observable image
$J(\Energy_\Obs)$. The practical content is a separation of duties:
the source side supplies the It\^o expansion; the compression
construction (\cref{thm:compression,thm:proj-div}) supplies the
observable lift; the correction term is the observable projection of
the source remainder.

\subsection*{IBM specialisation: target identity}

For IBM with source $\Source=(X,Y)$, the joint Wiener space carries
the classical Skorohod divergence $\delta_\Source(u_X,u_Y)=\delta_Xu_X
+\delta_Yu_Y$ and the source second-order derivatives $D_X^{2}$,
$D_XD_Y$, $D_YD_X$, $D_Y^{2}$ of \eqref{eq:DS2-IBM}. The intrinsic IBM
It\^o identity that \cref{thm:transfer} would specialise to has the
schematic form
\begin{equation}\label{eq:ibm-target}
  f(\IBM_t)-f(\IBM_0)
  \;=\;\delta_\IBM\bigl(B_f(t)\bigr)\;+\;C_\IBM(f;t),
  \qquad C_\IBM(f;t)\;=\;P_\IBM\,R_f(t),
\end{equation}
where $B_f(t)\in\Energy_\IBM$ is the first-order observable integrand
and $C_\IBM(f;t)$ is the observable second-order correction obtained
from the closed second-order IBM energy $\Energy_\IBM^{(2)}$ of
\cref{def:second-order-ibm}. \Cref{eq:ibm-target} is the
\emph{target identity}, not a theorem in full generality: it requires
both the source-level expansion \eqref{eq:source-ito} for $\IBM$ and
the verification that $A_f(t)\in J(\Energy_\IBM)$ on a useful test
class. The schematic version is recorded below.

\begin{conjecture}[Schematic IBM It\^o]\label{conj:ibm-ito}
Suppose $D_\IBM$ admits a cylindrical reduction
(\cref{rem:two-layer}) and the second-order datum
$(\Energy_\IBM^{(2)},D_\IBM^{(2)})$ is closed and densely defined.
Then for every $f$ in a suitable smooth class and every $t\ge 0$,
\eqref{eq:ibm-target} holds with
\begin{enumerate}
\item $B_f(t)\in\Energy_\IBM$ a representation integrand associated
  to $f$ and to the energy structure of $\IBM$ on $[0,t]$;
\item the correction $C_\IBM(f;t)$ built from a contraction of
  $D_\IBM^{(2)}$ against a smooth functional of $f$, with source lift
  $J^{\otimes 2}\,C_\IBM(f;t)$ compatible with the
  bi-Laplacian/local-time PDE signature identified
  in~\cite{AllouZheng};
\item $C_\IBM(f;t)$ is \emph{not} of the form
  $\tfrac12\int_{0}^{t}f''(\IBM_{s})\,d\langle\IBM\rangle_{s}$, since
  $\IBM$ does not supply the semimartingale quadratic-variation
  structure $\langle\IBM\rangle$ needed for a classical It\^o
  correction of that form (in its own filtration $\IBM$ is not a
  semimartingale, so no such $d\langle\IBM\rangle$ exists);
  quadratic variation is the wrong organising invariant, and the
  observable energy order (fourth, not second) is the right one
  (cf.~the energy-order discussion after
  \cref{thm:energy-first-ito-moll}).
\end{enumerate}
\end{conjecture}

\begin{remark}[Status of \cref{conj:ibm-ito}]
The first-order pieces ($\delta_\IBM, B_f(t)$) are objects supplied by
the construction of Section~\ref{sec:ibm-source} and made concrete on
the mollified occupation core by \cref{thm:ibm-occ}; the second-order
correction $C_\IBM(f;t)$ presupposes the second-order datum of
Section~\ref{sec:second-order}. The conjecture commits to a precise
shape via \cref{thm:transfer} but not to a specific functional form.
\end{remark}

\subsection*{Fourier-mode second-order match}

Under the convention $\Var(X_s)=s$, the Brownian-time semigroup
observable $u_k(t,x)=\E[e^{ik(x+X_{|Y_t|})}]$ admits an explicit
second-order PDE shadow with the bi-Laplacian Fourier multiplier
$k^{4}/8$. We derive it directly from the half-normal density of
$|Y_t|$ via an ordinary differential equation in $t$, independently
of the closed-form expression \eqref{eq:uk-closed}.

\begin{proposition}[Fourier-mode bi-Laplacian shadow]
\label{prop:fourier-2nd}
Let $a=k^{2}/2$ and
\[
  q(t)\;\coloneqq\;\int_0^\infty e^{-as}\,p_t(s)\,ds,
  \qquad p_t(s)=\sqrt{\tfrac{2}{\pi t}}\,e^{-s^{2}/(2t)},\quad s\ge 0,
\]
so that $u_k(t,x)=e^{ikx}q(t)$. Then $q$ satisfies
\begin{equation}\label{eq:q-ode}
  q'(t)\;=\;\frac{a^{2}}{2}\,q(t)\;-\;\frac{a}{\sqrt{2\pi t}}
  \;=\;\frac{k^{4}}{8}\,q(t)\;-\;\frac{k^{2}}{2\sqrt{2\pi t}},
\end{equation}
and equivalently
\begin{equation}\label{eq:az-prop}
  \partial_t u_k(t,x)
  \;=\;\frac{1}{8}\,\Delta^{2}u_k(t,x)
       \;+\;\frac{\Delta f_k(x)}{2\sqrt{2\pi t}},
  \qquad f_k(x)=e^{ikx}.
\end{equation}
\end{proposition}

\begin{proof}
The Gaussian density $\phi_t(s)=(2\pi t)^{-1/2}e^{-s^{2}/(2t)}$
satisfies the heat equation $\partial_t\phi_t=\tfrac12\partial_s^{2}\phi_t$
pointwise; since $p_t=\sqrt{2}\,\phi_t|_{s\ge 0}$, the same identity
holds for $p_t$ at every $s>0$. Therefore
\[
  q'(t)\;=\;\int_0^\infty e^{-as}\,\partial_t p_t(s)\,ds
   \;=\;\tfrac{1}{2}\int_0^\infty e^{-as}\,\partial_s^{2}p_t(s)\,ds.
\]
Integrate by parts twice in $s$. Boundary terms at $s=\infty$ vanish by
super-exponential decay of $p_t$; at $s=0$ we have
$\partial_s p_t(0)=-\sqrt{2/(\pi t)}\cdot(0/t)=0$ and
$p_t(0)=\sqrt{2/(\pi t)}$. Hence
\[
  \int_0^\infty e^{-as}\,\partial_s^{2}p_t(s)\,ds
  \;=\;a\!\int_0^\infty e^{-as}\,\partial_s p_t(s)\,ds
  \;=\;a\bigl(-p_t(0)+a\,q(t)\bigr)
  \;=\;a^{2}q(t)-a\,\sqrt{\tfrac{2}{\pi t}}.
\]
Dividing by $2$ and using $\sqrt{2}/(2\sqrt{\pi t})=1/\sqrt{2\pi t}$
gives the first equality of \eqref{eq:q-ode}; substituting
$a=k^{2}/2$ gives the second, since $a^{2}/2=k^{4}/8$ and
$a/\sqrt{2\pi t}=k^{2}/(2\sqrt{2\pi t})$.

For \eqref{eq:az-prop}, multiply \eqref{eq:q-ode} by $e^{ikx}$:
$\partial_t u_k(t,x)=\frac{k^{4}}{8}u_k(t,x)-\frac{k^{2}e^{ikx}}{2\sqrt{2\pi t}}$.
Since $\Delta e^{ikx}=-k^{2}e^{ikx}$ and $\Delta^{2}e^{ikx}=k^{4}e^{ikx}$,
the right-hand side is $\tfrac18\Delta^{2}u_k(t,x)+\Delta f_k(x)/(2\sqrt{2\pi t})$.
\end{proof}

Fourier-synthesis on smooth bounded $f$ gives, with $u(t,x)=\E[f(x+X_{|Y_t|})]$,
\begin{equation}\label{eq:az-recap}
  \partial_t u(t,x)
  \;=\;\frac{1}{8}\,\Delta^{2}u(t,x)
       \;+\;\frac{\Delta f(x)}{2\sqrt{2\pi t}}.
\end{equation}
The constants $1/8$ and $1/(2\sqrt{2\pi})$ are convention-dependent:
rescaling $\Var(X_s)=\alpha s$ rescales them to $\alpha^{2}/8$ and
$\alpha/(2\sqrt{2\pi t})$.

\paragraph{Second-order target.}
\Cref{prop:fourier-2nd} is the first concrete second-order target for
the intrinsic IBM observable energy: any closed second-order
observable energy producing an IBM It\^o formula via
\cref{thm:transfer} must reproduce the $k^{4}/8$ multiplier on Fourier
modes and the singular $t^{-1/2}\Delta f$ boundary term. We do
\emph{not} derive \eqref{eq:az-prop} from $D_\IBM^{(2)}$ of
\cref{def:second-order-ibm}; the proposition is a PDE-shadow
identity at the semigroup level, and we use it as the consistency
target that any future closed-form action of $D_\IBM^{(2)}$, averaged
against the law of $Y$, must satisfy. \Cref{eq:az-prop} operates at
the semigroup level $u=\E[f(\cdot)]$; \cref{conj:ibm-ito} operates
pathwise on $f(\IBM_t)$.

\paragraph{Source-bidegree reading.}
In the source atlas of~\cite{Fontes-IBMchaos}, the $k^{4}$ Fourier
multiplier is the signature of second-order scenery variation
transported through the Brownian clock --- the $(2,0)$, $(2,1)$, and
$(1,1)$ sectors interacting with the clock distribution --- and the
boundary $t^{-1/2}\Delta f$ term records the clock's local-time
singularity at the origin of the half-normal density.

\subsection*{Semigroup-level IBM It\^o correction}

\paragraph{Level structure.}
The paper occupies three proved layers and one open target:
\begin{itemize}
\item[\textbf{(L1)}] \emph{Proved abstract operator theorem.}
  Observable compression and projected divergence
  (\cref{thm:compression,thm:proj-div}).
\item[\textbf{(L2)}] \emph{Proved / controlled IBM core.}
  Mollified occupation core; $F_{T,\varepsilon}\to F_T$ in $L^{2}$ and
  $D_X F_{T,\varepsilon}\to D_X F_T$ in scenery energy; the clock
  direction recorded as an obstruction
  (\cref{thm:ibm-occ}, \cref{cor:unmollified}).
\item[\textbf{(L3)}] \emph{Deterministic semigroup benchmark.}
  The Brownian-time PDE correction (\cref{thm:bt-correction} below).
\item[\textbf{(L4)}] \emph{Future unmollified quartic-dual
  pathwise / operator-It\^o theorem.} Still conditional and open
  (\cref{conj:ibm-ito}, conditional form of \cref{thm:transfer});
  the correction is expected to live in the quartic-energy dual
  $\Energy_\IBM^{(4),-}$ of \eqref{eq:rigged-4}.
\end{itemize}

The present subsection records~(L3): the deterministic PDE
benchmark that any pathwise / operator-It\^o formula for IBM must
reproduce after full expectation. The relevant operator-energy
input is the second operator-derivative energy $\Energy_\IBM^{(2)}$
of \cref{def:second-order-ibm}, whose carr\'e / tensor structure
$\Energy_\IBM^{(4)}=\Energy_\IBM^{(2)}\otimes\Energy_\IBM^{(2)}$
produces the fourth-order observable shadow; the correction
$C_\IBM(f;t)$ itself is expected in the dual
$\Energy_\IBM^{(4),-}$, not in classical $L^{2}(\Omega)$. The identity below is the known Brownian-time
semigroup PDE; we give a direct derivation from the half-normal
integral representation, so that constants and convention are
explicit. This is a deterministic PDE statement, source-neutral with
respect to the operator-energy framework --- the framework's role is
to use it as a benchmark, not to derive it.

\begin{theorem}[Brownian-time semigroup correction]
\label{thm:bt-correction}
Let $X$ be standard two-sided Brownian scenery with $\Var(X_s)=s$ and
$Y$ an independent standard Brownian clock. For $f\in C_c^{\infty}(\R)$
(or, more generally, $f\in\mathscr S(\R)$), define
\[
  u(t,x)\;\coloneqq\;\E[f(x+X_{|Y_t|})],
  \qquad t\ge 0,\ x\in\R.
\]
Then $u\in C^{\infty}((0,\infty)\times\R)$, $u(0+,x)=f(x)$, and for
every $t>0$,
\begin{equation}\label{eq:bt-correction}
  \partial_t u(t,x)
  \;=\;\frac{1}{8}\,\partial_x^{4}u(t,x)
       \;+\;\frac{1}{2\sqrt{2\pi t}}\,\partial_x^{2}f(x).
\end{equation}
\end{theorem}

\begin{proof}
Conditioning on $|Y_t|$ and using independence of $X$ and $Y$,
\begin{equation}\label{eq:bt-integral}
  u(t,x)\;=\;\int_0^\infty (P_s f)(x)\,p_t(s)\,ds,
\end{equation}
where $P_s f(x)=\E[f(x+X_s)]=(G_s*f)(x)$ is the heat semigroup with
variance $s$, so $\partial_s P_s f=\tfrac12\,\partial_x^{2} P_s f$,
and $p_t(s)=\sqrt{2/(\pi t)}\,e^{-s^{2}/(2t)}$ is the half-normal
density.

The Gaussian density $\phi_t(s)=(2\pi t)^{-1/2}e^{-s^{2}/(2t)}$
satisfies the heat equation $\partial_t\phi_t=\tfrac12\,\partial_s^{2}
\phi_t$ pointwise; since $p_t=\sqrt 2\,\phi_t|_{s\ge 0}$, the same
identity holds for $p_t$ at every $s>0$, with $\partial_s p_t(0)=0$
and $p_t(0)=\sqrt{2/(\pi t)}$. Therefore
\[
  \partial_t u(t,x)
  \;=\;\int_0^\infty(P_s f)(x)\,\partial_t p_t(s)\,ds
   \;=\;\tfrac12\int_0^\infty(P_s f)(x)\,\partial_s^{2}p_t(s)\,ds.
\]
Integrate by parts twice in $s$; boundary terms at $s=\infty$ vanish
by rapid decay of $p_t$ and boundedness of $P_s f$:
\[
\begin{aligned}
  \tfrac12\!\int_0^\infty\!(P_s f)(x)\,\partial_s^{2}p_t(s)\,ds
  &= -\tfrac12\!\int_0^\infty\!\partial_s(P_s f)(x)\,\partial_s p_t(s)\,ds
   \quad(\partial_s p_t(0)=0)\\
  &= -\tfrac14\!\int_0^\infty\!\partial_x^{2}P_s f(x)\,\partial_s p_t(s)\,ds
   \quad(\partial_s P_s=\tfrac12\,\partial_x^{2} P_s)\\
  &= -\tfrac14\bigl[\partial_x^{2}P_s f(x)\,p_t(s)\bigr]_0^\infty
   +\tfrac14\!\int_0^\infty\!\partial_s\partial_x^{2}P_s f(x)\,p_t(s)\,ds\\
  &= \tfrac14\,\partial_x^{2}f(x)\,p_t(0)
   +\tfrac18\!\int_0^\infty\!\partial_x^{4}P_s f(x)\,p_t(s)\,ds.
\end{aligned}
\]
The boundary term at $s=0$ uses $P_0 f=f$ and $p_t(0)=\sqrt{2/(\pi t)}$,
so $\tfrac14\partial_x^{2}f(x)\,p_t(0)=\partial_x^{2}f(x)/(2\sqrt{2\pi t})$.
The remaining integral equals $\tfrac18\,\partial_x^{4}u(t,x)$ by the
representation \eqref{eq:bt-integral}. Smoothness of $u$ on
$(0,\infty)\times\R$ and the initial condition $u(0+,x)=f(x)$ follow
from standard regularising properties of the Gaussian heat kernel
$G_s$ and the half-normal density $p_t$.
\end{proof}

\paragraph{Connection to the operator-It\^o transfer.}
\Cref{thm:bt-correction} identifies the deterministic correction that
any closed second-order observable IBM energy must reproduce on
Brownian-time semigroup observables. In the notation of
\cref{thm:transfer}, the operator-It\^o transfer produces an
observable correction $P_\IBM R_f(t)$ that is itself
$\mathcal F_\IBM$-measurable; the right-hand side of
\eqref{eq:bt-correction} is its \emph{full-expectation} reduction
$\E[P_\IBM R_f(t)]=\E[R_f(t)]$ for the Brownian-time class. We do
\emph{not} derive \eqref{eq:bt-correction} from $D_\IBM^{(2)}$ of
\cref{def:second-order-ibm}; \cref{thm:bt-correction} is a closed PDE
identity at the deterministic level, while a closed-form action of
$D_\IBM^{(2)}$ would be needed for the pathwise observable expansion.

\paragraph{Deterministic semigroup identity.}
Integrating \eqref{eq:bt-correction} in $t$ on $(0,t]$ and using
$u(0+,x)=f(x)$ gives the deterministic semigroup identity
\begin{equation}\label{eq:bt-target}
  u(t,x)-f(x)
  \;=\;\int_0^t\Bigl[\,\tfrac18\,\partial_x^{4}u(s,x)
       \;+\;\tfrac{1}{2\sqrt{2\pi s}}\,\partial_x^{2}f(x)\,\Bigr]\,ds,
\end{equation}
valid for $f\in\mathscr S(\R)$. The right-hand side contains no
divergence term: \eqref{eq:bt-target} operates strictly at the
deterministic PDE / semigroup level.

\paragraph{Benchmark for the next theorem.}
\begin{center}
\fbox{\begin{minipage}{0.92\linewidth}
\emph{Any pathwise / operator-It\^o formula for IBM must reproduce,
after full expectation, the integrated deterministic correction}
\begin{equation}\label{eq:bt-benchmark}
\begin{split}
  C^{\mathrm{det}}_\IBM(f;t,x)
   &\;\coloneqq\;\int_0^{t}\!\Bigl[\tfrac{1}{8}\partial_x^{4}u(s,x)
   +\tfrac{\partial_x^{2}f(x)}{2\sqrt{2\pi s}}\Bigr]\,ds\\
   &\;=\;u(t,x)-f(x),
\end{split}
\end{equation}
\emph{for $t>0$, $f\in\mathscr S(\R)$, with
$u(t,x)=\E[f(x+X_{|Y_t|})]$ under the convention $\Var(X_s)=s$.}
\end{minipage}}
\end{center}

\paragraph{Relation to the future unmollified quartic-dual
pathwise/operator-It\^o target.}
A future unmollified pathwise It\^o identity for IBM would have the
schematic form
\begin{equation}\label{eq:pathwise-target}
  f(x+X_{|Y_t|})-f(x)
  \;=\;\delta_\IBM(B_f(t))\;+\;C_\IBM(f;t),
\end{equation}
with $B_f(t)$ a first-order observable integrand and $C_\IBM(f;t)$
the observable correction; the proved obstruction
\cref{prop:var-divergence} rules out $C_\IBM(f;t)\in L^{2}(\Omega)$,
so $C_\IBM(f;t)$ is expected in the quartic-energy dual
$\Energy_\IBM^{(4),-}$ of \eqref{eq:rigged-4}, paired against test
integrands in $\Energy_\IBM^{(4),+}$. \Cref{eq:pathwise-target} is
\emph{not} proved here; it remains the subject of
\cref{conj:ibm-ito} and the conditional operator-It\^o transfer
principle \cref{thm:transfer}. Passing from such a pathwise
identity to the deterministic Brownian-time semigroup means taking
the \emph{full} expectation over the source randomness: the centred
divergence term $\delta_\IBM(B_f(t))$ should vanish, while
$\E[C_\IBM(f;t)]$ should reduce to the deterministic correction in
\eqref{eq:bt-target}. The required operator-energy input is the
second operator-derivative energy $\Energy_\IBM^{(2)}$ of
\cref{def:second-order-ibm}, whose tensor square produces the
fourth-order observable shadow; the semigroup statement does not
require this input.

\paragraph{Source-resolved chaos perspective.}
The source-resolved chaos paper~\cite{Fontes-IBMchaos} explains why
the fourth-order signature in \eqref{eq:bt-correction} should be
read as a statistical / operator shadow of pure and mixed source
bidegrees --- $(2,0)$, $(0,2)$, $(1,1)$, $(2,1)$, $(1,2)$ ---
\emph{not} as a Markov generator of $\IBM$: $\IBM$ is not Markov in
its own filtration, and the mixed sectors $(m,n)$ with $m,n\ge 1$
obstruct a closed scalar PDE for pathwise occupation functionals.
The present paper turns that interpretation into an operator-energy
benchmark: \cref{thm:bt-correction} concerns the \emph{averaged}
Brownian-time semigroup, where the average over $Y$ collapses the
mixed sectors and exposes the closed PDE \eqref{eq:bt-correction},
and \eqref{eq:bt-benchmark} records the deterministic correction that
the closed second-order observable energy must reproduce. Equivalently,
the Fourier-mode diagnostic of \cref{prop:fourier-2nd} identifies what
$C_\IBM(f;t)$ in \eqref{eq:pathwise-target} must reduce to after full
expectation: the $k^{4}/\Delta^{2}$ multiplier together with the
$t^{-1/2}\Delta f$ boundary term.

\subsection*{Current status}

\begin{itemize}
\item first-order observable compression: proved abstractly
  (\cref{thm:compression,thm:proj-div});
\item mollified IBM occupation derivative: constructed explicitly
  (\cref{thm:ibm-occ}~(i)--(ii));
\item scenery convergence: controlled at rate $O(\varepsilon^{1/2})$
  (\cref{thm:ibm-occ}~(iii)--(iv));
\item clock component: singular, no $L^{2}$ limit asserted
  (\cref{cor:unmollified});
\item abstract operator-It\^o transfer principle: proved conditionally
  (\cref{thm:transfer});
\item Brownian-time Fourier diagnostic: identifies the target
  correction \eqref{eq:az-recap};
\item candidate second-order observable energy: defined
  (\cref{def:second-order-ibm,prop:compression-2});
\item the closed second-order IBM observable energy and the full
  proof of \cref{conj:ibm-ito}: remain the next analytic theorem.
\end{itemize}

\section{A semigroup-level IBM correction}\label{sec:sg-correction}

We now name and isolate the deterministic correction that
\cref{thm:bt-correction} produces, and record one consistency
consequence for any future unmollified quartic-dual pathwise IBM
It\^o formula. The content of
this section is a packaging of \cref{thm:bt-correction} with explicit
named symbols, plus a single conditional corollary; no new analytic
work beyond \cref{thm:bt-correction} is performed.

\subsection*{The semigroup correction}

Let $X$ be standard two-sided Brownian scenery with $\Var(X_s)=s$ and
$Y$ an independent standard Brownian clock. For $f\in\mathscr S(\R)$,
define the Brownian-time semigroup observable
\begin{equation}\label{eq:sg-u}
  u(t,x)\;=\;\E[f(x+X_{|Y_t|})]
   \;=\;\int_0^\infty(P_s f)(x)\,p_t(s)\,ds,
\end{equation}
where $P_s=e^{s\Delta/2}$ is the heat semigroup of variance $s$ and
$p_t(s)=\sqrt{2/(\pi t)}\,e^{-s^{2}/(2t)}$, $s\ge 0$, is the half-normal
density.

\begin{definition}[Semigroup correction density and integrated correction]
\label{def:sg-correction}
For $f\in\mathscr S(\R)$ and $t>0$, the \emph{semigroup correction
density} is
\begin{equation}\label{eq:sg-density}
  \mathcal C_\IBM^{\mathrm{sg}}(f;t,x)
  \;\coloneqq\;\frac{1}{8}\,\partial_x^{4}u(t,x)
   \;+\;\frac{1}{2\sqrt{2\pi t}}\,\partial_x^{2}f(x),
\end{equation}
and the \emph{integrated semigroup correction} is
\begin{equation}\label{eq:sg-integrated}
  C_\IBM^{\mathrm{sg}}(f;t,x)
  \;\coloneqq\;\int_0^t\mathcal C_\IBM^{\mathrm{sg}}(f;s,x)\,ds.
\end{equation}
\end{definition}

\begin{theorem}[Semigroup-level IBM correction]\label{thm:sg-correction}
For $f\in\mathscr S(\R)$ and $t>0$,
\begin{equation}\label{eq:sg-main}
  u(t,x)-f(x)\;=\;C_\IBM^{\mathrm{sg}}(f;t,x),
\end{equation}
equivalently $\partial_t u(t,x)=\mathcal C_\IBM^{\mathrm{sg}}(f;t,x)$.
The integrand defining $C_\IBM^{\mathrm{sg}}$ is integrable on
$(0,t]$: the boundary singularity
$s^{-1/2}\partial_x^{2}f(x)$ is of integrable type, since
$\int_0^t s^{-1/2}\,ds=2\sqrt t<\infty$ and $\partial_x^{2}f$ is
bounded for $f\in\mathscr S(\R)$.
\end{theorem}

\begin{proof}
The pointwise-in-$t$ identity
$\partial_t u(t,x)=\mathcal C_\IBM^{\mathrm{sg}}(f;t,x)$ is precisely
\cref{thm:bt-correction}. Integrate over $(0,t]$ and use the initial
condition $u(0+,x)=f(x)$ to obtain \eqref{eq:sg-main}. For
integrability of the integrand near $s=0$: the first term
$\tfrac18\partial_x^{4}u(s,x)$ is bounded uniformly on compact sets in
$(s,x)$ by the smoothing of the Gaussian heat kernel applied to
$f\in\mathscr S(\R)$ together with the half-normal kernel; the second
term satisfies $\int_0^{t}(2\sqrt{2\pi s})^{-1}\,ds=\sqrt{t/(2\pi)}<\infty$.
\end{proof}

\subsection*{Expectation-level interpretation}

\Cref{thm:sg-correction} constructs the deterministic correction
$C_\IBM^{\mathrm{sg}}(f;t,x)$. It is not the pathwise correction
$C_\IBM(f;t)$ of \eqref{eq:pathwise-target}; it fixes the
expectation-level value that any pathwise correction must take on the
Brownian-time semigroup class:
\[
  \E[C_\IBM(f;t)]\;=\;C_\IBM^{\mathrm{sg}}(f;t,x)
\]
under the assumption that the pathwise identity exists and the
divergence term is centred. The next corollary records this as a
hard consistency requirement.

\begin{corollary}[Conditional pathwise-to-semigroup consistency]
\label{cor:pathwise-to-sg}
Suppose, on a domain of test functions $f\in\mathscr S(\R)$ and time
horizons $t>0$, a pathwise / operator-It\^o identity
of the form \eqref{eq:pathwise-target} holds:
\[
  f(x+X_{|Y_t|})-f(x)\;=\;\delta_\IBM(B_f(t))+C_\IBM(f;t),
\]
with $\E[\delta_\IBM(B_f(t))]=0$ (centred divergence). Then necessarily
\begin{equation}\label{eq:consistency}
  \E[C_\IBM(f;t)]\;=\;C_\IBM^{\mathrm{sg}}(f;t,x)\;=\;u(t,x)-f(x).
\end{equation}
\end{corollary}

\begin{proof}
Take full expectation of \eqref{eq:pathwise-target}:
$\E[f(x+X_{|Y_t|})]-f(x)=\E[\delta_\IBM(B_f(t))]+\E[C_\IBM(f;t)]
=\E[C_\IBM(f;t)]$. The left-hand side is $u(t,x)-f(x)$, which by
\cref{thm:sg-correction} equals $C_\IBM^{\mathrm{sg}}(f;t,x)$.
\end{proof}

\Cref{cor:pathwise-to-sg} is a theorem of \emph{consistency}, not of
existence: it does not produce a pathwise identity, but constrains the
expectation-level value of any such pathwise identity should one
exist. In particular, any candidate $C_\IBM(f;t)$ produced from a
closed-form action of $D_\IBM^{(2)}$ on the Brownian-time semigroup
class must satisfy \eqref{eq:consistency}.

\subsection*{Connection to the second-order observable energy}

The remaining open problem is to identify $C_\IBM(f;t)$ intrinsically
from a closed-form action of the second-order observable energy
$\Energy_\IBM^{(2)}$ (\cref{def:second-order-ibm}).
\Cref{thm:sg-correction} pins down the deterministic value that this
second-order correction must have after full expectation;
\cref{cor:pathwise-to-sg} converts the benchmark into a hard
consistency requirement on any candidate $C_\IBM(f;t)$. Producing
such a candidate from $D_\IBM^{(2)}$ on a suitable test class, and
verifying \eqref{eq:consistency}, is the natural next analytic step;
the pathwise correction itself remains open.

\section{The proved IBM It\^o formula at the mollified level}\label{sec:moll-pathwise}

This section is the calibration locus of the paper: we prove the
mollified pathwise IBM It\^o formula and its fourth-order semigroup
consistency. The open problem is the unmollified pathwise correction
in the quartic observable-energy dual, addressed in
\cref{sec:toward-unmoll}.

The proved IBM It\^o formula combines four already-established
ingredients: the mollified pathwise It\^o--Tanaka identity
(\cref{prop:moll-tanaka}), the operator-divergence reading of the
martingale term (\cref{thm:transfer} specialised to IBM), the
explicit mollified pathwise correction
$C_{\IBM,\varepsilon}(f;t)$, and its expectation-level convergence
to the fourth-order Brownian-time PDE correction
(\cref{thm:moll-consistency} combined with \cref{thm:bt-correction}).
We state the combined theorem first, then verify its constituent
pieces in the subsections that follow.

\begin{theorem}[Energy-first It\^o formula for mollified IBM]
\label{thm:energy-first-ito-moll}
Let $X$ be standard two-sided Brownian scenery with $\Var(X_s)=|s|$,
$Y$ an independent standard Brownian clock, $\rho$ an even
non-negative compactly supported mollifier, $X^\varepsilon=
\rho_\varepsilon*X$, and
\[
  \IBM_t^{\,\varepsilon}\;\coloneqq\;x+X^\varepsilon_{|Y_t|},
  \qquad t\ge 0,\;x\in\R.
\]
For $f\in\mathscr S(\R)$, set $g_\varepsilon(r)\coloneqq
f(x+X^\varepsilon_r)$, $r\in\R$, which is $\Prob_X$-a.s.\ $C^\infty$
in $r$. Then for every $\varepsilon>0$ and $t\ge 0$:

\smallskip
\noindent\textbf{(i) Mollified pathwise It\^o--Tanaka identity.}
\begin{equation}\label{eq:efi-ito}
  f(\IBM_t^{\,\varepsilon})-f(\IBM_0^{\,\varepsilon})
  \;=\;\delta_Y\bigl(B_{f,\varepsilon}(t)\bigr)
   \;+\;C_{\IBM,\varepsilon}(f;t),
\end{equation}
with clock-divergence integrand
\begin{equation}\label{eq:efi-B}
  B_{f,\varepsilon}(t)(s)
   \;=\;g_\varepsilon'(|Y_s|)\,\mathrm{sgn}(Y_s)\,
        \mathbf 1_{[0,t]}(s),
\end{equation}
and mollified pathwise correction
\begin{equation}\label{eq:efi-C}
  C_{\IBM,\varepsilon}(f;t)
   \;=\;g_\varepsilon'(0)\,L_t^{0}(Y)
   \;+\;\tfrac12\!\int_0^t g_\varepsilon''(|Y_s|)\,ds,
\end{equation}
where $L_t^{0}(Y)$ is the Brownian local time of $Y$ at $0$ and the
clock derivatives are explicit:
\begin{equation}\label{eq:efi-gpp}
  g_\varepsilon'(r)
   \;=\;f'(x+X^\varepsilon_r)\,\partial_rX^\varepsilon_r,\qquad
  g_\varepsilon''(r)
   \;=\;f''(x+X^\varepsilon_r)\,(\partial_rX^\varepsilon_r)^{2}
   \;+\;f'(x+X^\varepsilon_r)\,\partial_r^{2}X^\varepsilon_r.
\end{equation}

\smallskip
\noindent\textbf{(ii) Operator-divergence reading of the martingale
term.} The term $\delta_Y(B_{f,\varepsilon}(t))$ in \eqref{eq:efi-ito}
is the source clock Skorohod divergence of the integrand
\eqref{eq:efi-B}; it coincides pathwise with the conditional It\^o
integral $\int_0^t g_\varepsilon'(|Y_s|)\,\mathrm{sgn}(Y_s)\,dY_s$, is
a continuous $L^{2}$-martingale in $Y$ (conditional on $X$), and
satisfies $\E[\delta_Y(B_{f,\varepsilon}(t))]=0$. In the language of
the conditional operator-It\^o transfer principle
(\cref{thm:transfer}), \eqref{eq:efi-ito} is the mollified instance
of the schematic pathwise target \eqref{eq:pathwise-target}, with the
clock-source divergence carrying the martingale part.

\smallskip
\noindent\textbf{(iii) Fourth-order expectation-level consistency.}
With $u(t,x)=\E[f(x+X_{|Y_t|})]$,
\begin{equation}\label{eq:efi-fourth}
  \lim_{\varepsilon\to 0}\E\bigl[C_{\IBM,\varepsilon}(f;t)\bigr]
   \;=\;\int_0^t\!\Bigl[\,\tfrac{1}{8}\,\partial_x^{4}u(s,x)
        \;+\;\tfrac{1}{2\sqrt{2\pi s}}\,\partial_x^{2}f(x)\,\Bigr]\,ds
   \;=\;u(t,x)-f(x).
\end{equation}
\end{theorem}

\begin{proof}
Part (i) is \cref{prop:moll-tanaka} below, repackaged using
\eqref{eq:efi-C}--\eqref{eq:efi-gpp}. For (ii), the integrand
$B_{f,\varepsilon}(t)$ is $Y$-adapted with continuous paths in $s$
and we verify conditional square-integrability for each $\omega_X$
as follows. By \cref{def:mollified}, $X^\varepsilon$ is
$\Prob_X$-a.s.\ $C^{\infty}$ in $r$, hence the path
$r\mapsto\partial_r X^\varepsilon_r$ is locally bounded a.s.\ on
every compact $r$-interval. On $[0,t]$ the random variable
$|Y_s|$ ranges over a compact subinterval of $\R_{+}$ (the path of
$|Y|$ on $[0,t]$ is continuous and so attains a finite maximum).
Combined with the Schwartz bounds on $f$ (so $f'$ is bounded on
$\R$), this gives
\[
  \E_Y\!\int_0^{t}g_\varepsilon'(|Y_s|)^{2}\,ds
   \;\le\;\|f'\|_\infty^{2}\,
    \E_Y\!\int_0^{t}\bigl(\partial_r X^\varepsilon_{|Y_s|}\bigr)^{2}\,ds
   \;<\;\infty\quad\text{a.s.\ in }\omega_X,
\]
since $\partial_r X^\varepsilon$ is continuous in $r$ and $|Y|$ is
bounded on $[0,t]$ a.s. The source clock Skorohod divergence on
$Y$-adapted square-integrable integrands coincides with the It\^o
integral~\cite{Nualart}, giving the identification
$\delta_Y(B_{f,\varepsilon}(t))=\int_0^t g_\varepsilon'(|Y_s|)
\mathrm{sgn}(Y_s)\,dY_s$ and the centring
$\E[\delta_Y(B_{f,\varepsilon}(t))]=0$. Part (iii) is
\cref{thm:moll-consistency} combined with the integrated form
\eqref{eq:sg-integrated} of \cref{thm:sg-correction}, which equates
$C_\IBM^{\mathrm{sg}}(f;t,x)$ with the time integral on the right of
\eqref{eq:efi-fourth}.
\end{proof}

\paragraph{The proved IBM It\^o formula of the paper.}
\Cref{thm:energy-first-ito-moll} is the proved IBM It\^o formula of
this paper. It is pathwise at the mollified scenery level and
fourth-order correct after expectation. The theorem does \emph{not}
assert convergence of $C_{\IBM,\varepsilon}(f;t)$ in
$L^{2}(\Omega)$ as $\varepsilon\to 0$, and the unmollified pathwise
correction $C_\IBM(f;t)$ is not constructed here as an
$L^{2}(\Omega)$ random variable; the obstruction result of
\cref{sec:toward-unmoll} (\cref{prop:var-divergence}) explains why
the natural unmollified target is instead the dual of the quartic
composite observable energy
$\Energy_\IBM^{(4)}=\Energy_\IBM^{(2)}\otimes\Energy_\IBM^{(2)}$ in
the rigged hierarchy \eqref{eq:rigged-4}.

\paragraph{Why the correction is quartic at the energy level.}
The clock It\^o formula behind \eqref{eq:efi-ito} is second order in
the Brownian-time variable $r$: it is the ordinary It\^o--Tanaka
formula applied to $g_\varepsilon(|Y_t|)$. The coefficients of that
second-order expansion, $g_\varepsilon'$ and $g_\varepsilon''$
displayed in \eqref{eq:efi-gpp}, are however first and second
\emph{spatial} derivatives of a Brownian scenery. The PDE shadow of
\eqref{eq:efi-C} in expectation, displayed on the right of
\eqref{eq:efi-fourth}, is therefore fourth order in the spatial
variable: the bi-Laplacian term $\tfrac18\,\partial_x^{4}u$ is the
shadow of the $f'(x+X^\varepsilon_r)\,\partial_r^{2}X^\varepsilon$
contribution to $g_\varepsilon''$ after averaging over the
half-normal clock distribution, and the boundary term
$\partial_x^{2}f/(2\sqrt{2\pi s})$ is its half-normal boundary
correction. This is the canonical example of energy-order
classification replacing semimartingale quadratic variation:
classical $\langle\cdot\rangle$-tracking sees only the second-order
clock term of \eqref{eq:efi-ito} and misses the fourth-order spatial
PDE shadow that is the actual It\^o correction for IBM.

\subsection*{Mollified pathwise It\^o--Tanaka identity}

For $\varepsilon>0$, $x\in\R$, and $t\ge 0$, set
$Z_t^\varepsilon=x+X^\varepsilon_{|Y_t|}$. For $f\in\mathscr S(\R)$,
define the conditionally smooth function in the clock variable
\[
  g_\varepsilon(r)\;\coloneqq\;f(x+X^\varepsilon_r),
  \qquad r\in\R,
\]
which is $\Prob_X$-a.s.\ $C^\infty$ in $r$.

\begin{proposition}[Mollified pathwise It\^o-Tanaka identity]
\label{prop:moll-tanaka}
For every $\varepsilon>0$, $t\ge 0$, $f\in\mathscr S(\R)$,
\begin{equation}\label{eq:moll-tanaka}
  f(Z_t^\varepsilon)-f(Z_0^\varepsilon)
  \;=\;M_{\varepsilon,t}^{Y}\;+\;g_\varepsilon'(0)\,L_t^{0}(Y)
       \;+\;\tfrac12\!\int_0^t g_\varepsilon''(|Y_s|)\,ds,
\end{equation}
where $L_t^{0}(Y)$ is the Brownian local time of $Y$ at $0$ and
\begin{equation}\label{eq:moll-mart}
  M_{\varepsilon,t}^{Y}\;=\;\int_0^t g_\varepsilon'(|Y_s|)\,
   \mathrm{sgn}(Y_s)\,dY_s
\end{equation}
is a continuous $L^{2}$-martingale in $Y$ (conditional on $X$). The
derivatives in \eqref{eq:moll-tanaka} are explicit:
\[
  g_\varepsilon'(r)=f'(x+X^\varepsilon_r)\,\partial_rX^\varepsilon_r,
  \quad
  g_\varepsilon''(r)=f''(x+X^\varepsilon_r)\,(\partial_rX^\varepsilon_r)^{2}
   +f'(x+X^\varepsilon_r)\,\partial_r^{2}X^\varepsilon_r.
\]
\end{proposition}

\begin{proof}
Tanaka's formula gives $d|Y_t|=\mathrm{sgn}(Y_t)dY_t+dL_t^{0}(Y)$ and
$\langle|Y|\rangle_t=t$. For each $\omega_X$, $g_\varepsilon\in C^\infty(\R)$,
so the It\^o-Tanaka formula applied to $g_\varepsilon(|Y_t|)$ yields,
using $g_\varepsilon'(|Y_s|)dL_s^{0}=g_\varepsilon'(0)dL_s^{0}$
(since $dL_s^{0}$ is supported on $\{Y_s=0\}$),
\[
  g_\varepsilon(|Y_t|)-g_\varepsilon(0)
  =\!\int_0^t\!\! g_\varepsilon'(|Y_s|)\mathrm{sgn}(Y_s)\,dY_s
   +g_\varepsilon'(0)L_t^{0}(Y)
   +\tfrac12\!\int_0^t\!\! g_\varepsilon''(|Y_s|)\,ds.
\]
The left-hand side equals $f(Z_t^\varepsilon)-f(Z_0^\varepsilon)$ since
$Z_t^\varepsilon=x+X^\varepsilon_{|Y_t|}$ and $|Y_0|=0$. The explicit
formulas for $g_\varepsilon',g_\varepsilon''$ are the chain rule.
\end{proof}

Define the \emph{mollified pathwise correction}
\begin{equation}\label{eq:moll-correction}
  C_{\IBM,\varepsilon}(f;t)
  \;\coloneqq\;
  g_\varepsilon'(0)\,L_t^{0}(Y)\;+\;\tfrac12\!\int_0^t\!\!
   g_\varepsilon''(|Y_s|)\,ds.
\end{equation}
Both contributions are pathwise well-defined random variables for each
$\varepsilon>0$. The first is a clock local-time term, the second is
the conditional It\^o second-order term.

\subsection*{Expectation-level convergence}

Set $v^\varepsilon(r,x)\coloneqq\E_X[f(x+X^\varepsilon_r)]$ and
$U^\varepsilon(t,x)\coloneqq\E[f(x+X^\varepsilon_{|Y_t|})]
=\int_0^\infty v^\varepsilon(r,x)\,p_t(r)\,dr$.

\begin{theorem}[Expectation-level consistency at $\varepsilon>0$]
\label{thm:moll-consistency}
For every $f\in\mathscr S(\R)$, $\varepsilon>0$, and $t>0$,
\begin{equation}\label{eq:moll-exp-id}
  \E[C_{\IBM,\varepsilon}(f;t)]
   \;=\;U^\varepsilon(t,x)-v^\varepsilon(0,x).
\end{equation}
Consequently, as $\varepsilon\to 0$,
\begin{equation}\label{eq:moll-limit}
  \lim_{\varepsilon\to 0}\,\E[C_{\IBM,\varepsilon}(f;t)]
   \;=\;u(t,x)-f(x)
   \;=\;C_\IBM^{\mathrm{sg}}(f;t,x).
\end{equation}
\end{theorem}

\begin{proof}
Take expectation in \eqref{eq:moll-tanaka}. The It\^o integral
$M_{\varepsilon,t}^{Y}$ has zero $Y$-conditional mean, so
$\E[M_{\varepsilon,t}^{Y}]=0$. Therefore
$\E[f(Z_t^\varepsilon)-f(Z_0^\varepsilon)]=\E[C_{\IBM,\varepsilon}(f;t)]$.
By definition $\E[f(Z_t^\varepsilon)]=U^\varepsilon(t,x)$ and
$\E[f(Z_0^\varepsilon)]=\E[f(x+X^\varepsilon_0)]=v^\varepsilon(0,x)$,
giving \eqref{eq:moll-exp-id}.

For \eqref{eq:moll-limit}: by Lipschitz $f$ and the $L^{1}$ corollary
of \cref{lem:moll-rate} (uniform in $r\in\R$),
$\E|f(x+X^\varepsilon_r)-f(x+X_r)|\le\|f'\|_\infty\,
C_\rho^{1/2}\,\varepsilon^{1/2}$. Averaging against the law of
$|Y_t|$ gives $U^\varepsilon(t,x)\to u(t,x)$. At $r=0$,
\cref{lem:moll-rate} together with $X_0=0$ a.s.\ gives
$X^\varepsilon_0\to 0$ in $L^{2}$, so $v^\varepsilon(0,x)
=\E[f(x+X^\varepsilon_0)]\to f(x)$ by Lipschitz $f$. The second
equality in \eqref{eq:moll-limit} is \cref{thm:sg-correction}.
\end{proof}

\paragraph{Local-time term has zero expectation.}
The mollifier $\rho$ is non-negative and \emph{even}
(\cref{def:mollified}), so $\rho_\varepsilon$ is even and
$\rho_\varepsilon'$ is odd. A reflection-symmetry argument then
shows
\begin{equation}\label{eq:lt-zero-mean}
  \E\bigl[g_\varepsilon'(0)\,L_t^{0}(Y)\bigr]\;=\;0
  \qquad\text{for every }\varepsilon>0,\ t>0.
\end{equation}
Indeed, $L_t^{0}(Y)$ is $Y$-measurable and $g_\varepsilon'(0)$ is
$X$-measurable; by independence $\E[g_\varepsilon'(0)L_t^{0}(Y)]
=\E_X[g_\varepsilon'(0)]\,\E_Y[L_t^{0}(Y)]$. The two-sided Brownian
scenery has the symmetry $X_z\stackrel{d}{=}X_{-z}$ as random fields;
under the corresponding reflection, $X^\varepsilon_0
=\int\rho_\varepsilon(z)X_z\,dz$ is invariant (even integrand) while
$\partial_rX^\varepsilon_0=-\int\rho_\varepsilon'(z)X_z\,dz$ picks up a
sign (odd integrand). Therefore $\E_X[h(X^\varepsilon_0)
\partial_rX^\varepsilon_0]=0$ for every Borel $h$ with finite
expectation, and in particular $\E_X[g_\varepsilon'(0)]
=\E_X[f'(x+X^\varepsilon_0)\partial_rX^\varepsilon_0]=0$, giving
\eqref{eq:lt-zero-mean}.

In view of \eqref{eq:lt-zero-mean}, the whole expectation
\eqref{eq:moll-exp-id} is carried by the interior term:
\[
  \E[C_{\IBM,\varepsilon}(f;t)]
  \;=\;\tfrac12\int_0^t\E\bigl[g_\varepsilon''(|Y_s|)\bigr]\,ds
  \;=\;U^\varepsilon(t,x)-v^\varepsilon(0,x).
\]

\paragraph{Bulk and boundary structure of the interior expectation
(interpretive).}
The interior term in \eqref{eq:moll-correction} admits a heat-semigroup
reading. Conditioning on $X$ and using
$v^\varepsilon(r,x)=(P_{\sigma_\varepsilon^{2}(r)}f)(x)$ with
$\sigma_\varepsilon^{2}(r)\coloneqq\Var(X^\varepsilon_r)$,
$\E_X[g_\varepsilon''(r)]=\partial_r^{2}v^\varepsilon(r,x)$. For
$r>c\varepsilon$ (bulk), $(\sigma_\varepsilon^{2})'(r)=1$ exactly, so
$\partial_r v^\varepsilon=\tfrac12\Delta v^\varepsilon$ and
$\partial_r^{2}v^\varepsilon=\tfrac14\Delta^{2}v^\varepsilon
\xrightarrow{\varepsilon\to 0}\tfrac14\,\partial_x^{4}P_rf(x)$ pointwise
in $r$ on $(0,\infty)$. Heuristically, integrating against
$p_s(r)\,dr$ and combining with the outer $\tfrac12$ in
\eqref{eq:moll-correction} matches the bulk contribution
$\tfrac18\,\partial_x^{4}u(s,x)$ of \eqref{eq:bt-correction}. On the
boundary layer $r\in[0,c\varepsilon]$ near the half-normal origin,
$(\sigma_\varepsilon^{2})''(r)=O(\varepsilon^{-1})$ on the
scale-$\varepsilon$ transition to the bulk slope --- the mechanism by
which the boundary term $\partial_x^{2}f(x)/(2\sqrt{2\pi s})$ of
\eqref{eq:bt-correction} is identified at the expectation level
through the half-normal integration-by-parts identity used in
\cref{thm:bt-correction}. This bulk/boundary decomposition is
\emph{interpretive}: we do not separately prove the boundary-layer
limit. The global convergence
$\E[C_{\IBM,\varepsilon}(f;t)]\to C_\IBM^{\mathrm{sg}}(f;t,x)$ is
already established by \cref{thm:moll-consistency}; the decomposition
above only explains \emph{through which channel} bulk and boundary
contributions appear.

\paragraph{Local-time versus boundary correction: not the same object.}
The boundary term $\partial_x^{2}f(x)/(2\sqrt{2\pi s})$ in
\eqref{eq:bt-correction} is \emph{not} the expectation of the
Tanaka local-time term in this symmetric mollified model:
\eqref{eq:lt-zero-mean} shows that expectation is zero. Rather, the
boundary correction appears in the $\varepsilon\to 0$ limit through
the half-normal integration-by-parts boundary layer of the interior
second-derivative term, as identified above. The two phenomena ---
pathwise Tanaka local time and the semigroup boundary correction
--- coexist in the same identity \eqref{eq:moll-tanaka} but
contribute through different channels: the local-time term is
pathwise present but expectation-invisible, while the boundary
correction is the expectation-level boundary layer of the interior
term.

\subsection*{Operator-divergence interpretation}

In the notation of the operator-It\^o transfer principle
(\cref{thm:transfer}), $M_{\varepsilon,t}^{Y}$ is the clock-source
divergence integral $\delta_Y(B_{f,\varepsilon}(t))$ of
\eqref{eq:efi-B}: this is the operator-divergence reading recorded
as part~(ii) of \cref{thm:energy-first-ito-moll}. The observable
projection of $\delta_Y(B_{f,\varepsilon}(t))$ onto $\mathcal F_\IBM$,
if it converges in the unmollified limit, is the candidate
$\delta_\IBM(B_f(t))$ of the schematic pathwise target
\eqref{eq:pathwise-target}. We do not assert this projection limit
exists.

\subsection*{Status of the pathwise programme}

\begin{itemize}
\item \emph{Mollified pathwise IBM It\^o formula}: \textbf{proved}
  (\cref{thm:energy-first-ito-moll}~(i)--(ii); constituent pieces
  \cref{prop:moll-tanaka} and the operator-divergence reading).
\item \emph{Fourth-order expectation-level consistency}: \textbf{proved}
  (\cref{thm:energy-first-ito-moll}~(iii); constituent pieces
  \cref{thm:moll-consistency} and \cref{thm:sg-correction}). In
  particular, the local-time term is expectation-invisible
  (\eqref{eq:lt-zero-mean}), so the entire expectation is carried by
  the interior $g_\varepsilon''$-integral.
\item \emph{Raw $L^{2}(\Omega)$ unmollified pathwise correction}
  $C_\IBM(f;t)$: \textbf{obstructed}. The first-scenery-chaos
  candidate $\Psi_\varepsilon$ associated with the
  $\partial_r^{2}X^\varepsilon$ content of $g_\varepsilon''$ has
  $\Var(\Psi_\varepsilon)\to\infty$
  (\cref{prop:var-divergence}); ordinary Wick / Hermite renormalisation
  acts in $\Hilb_0^{X}\oplus\Hilb_2^{X}$ and structurally does not
  address this $\Hilb_1^{X}$ singularity
  (\cref{prop:wick-insufficient}).
\item \emph{Quartic observable-energy dual correction}: \textbf{open
  analytic target}. The unmollified pathwise correction $C_\IBM(f;t)$
  is expected to live in the dual of the quartic composite observable
  energy $\Energy_\IBM^{(4)}=\Energy_\IBM^{(2)}\otimes
  \Energy_\IBM^{(2)}$, in the rigged hierarchy of
  \eqref{eq:rigged-4}, with $\E[C_\IBM]=C_\IBM^{\mathrm{sg}}$ forced
  by \cref{cor:pathwise-to-sg}.
\end{itemize}

\section{Toward the unmollified pathwise correction}\label{sec:toward-unmoll}

We now investigate whether the mollified pathwise correction
$C_{\IBM,\varepsilon}(f;t)$ converges as $\varepsilon\to 0$ as a
random variable, beyond the expectation-level convergence of
\cref{thm:moll-consistency}. The conclusion of this section is
negative at the raw level: the centred fluctuation
$\tilde C_{\IBM,\varepsilon}\coloneqq C_{\IBM,\varepsilon}-\E[C_{\IBM,\varepsilon}]$
has variance that diverges as $\varepsilon\to 0$, so
$C_{\IBM,\varepsilon}$ does \emph{not} converge in $L^{2}(\Omega)$.
We then identify the natural energy-contraction / source-chaos
renormalisation that the unmollified pathwise correction should be
built from --- realised in the present Gaussian Brownian-scenery
source atlas as Wick / Hermite centering --- mark its convergence as
open, and place the analytic gap in the source-resolved chaos
coordinates.

\subsection*{Mean / fluctuation decomposition}

Set
\begin{equation}\label{eq:decomp}
  \tilde C_{\IBM,\varepsilon}(f;t)
   \;\coloneqq\;C_{\IBM,\varepsilon}(f;t)-\E[C_{\IBM,\varepsilon}(f;t)],
\end{equation}
so $C_{\IBM,\varepsilon}=\E[C_{\IBM,\varepsilon}]+\tilde C_{\IBM,\varepsilon}$
and $\E[\tilde C_{\IBM,\varepsilon}]=0$. By \cref{thm:moll-consistency},
the mean component converges to $C_\IBM^{\mathrm{sg}}(f;t,x)$. The
fluctuation $\tilde C_{\IBM,\varepsilon}$ is the new object to test.

\subsection*{Variance scaling of the constituents}

For $|y|>c\varepsilon$, the It\^o-isometry calculation of
\cref{lem:moll-eps-inv} extends to the second spatial derivative:
\begin{equation}\label{eq:var-2deriv}
  \Var\bigl(\partial_y^{2}X^\varepsilon_y\bigr)
  \;=\;\int_\R\rho_\varepsilon'(y-z)^{2}\,dz
  \;=\;\frac{c_\rho^{(2)}}{\varepsilon^{3}},
  \qquad c_\rho^{(2)}\coloneqq\int_\R\rho'(v)^{2}\,dv,
\end{equation}
since $\partial_y^{2}X^\varepsilon_y=\int\rho_\varepsilon'(y-z)\,dX_z$ in
the It\^o-integral form valid away from the origin. The pointwise
variance of $g_\varepsilon''(r)$ is therefore dominated, for
$f\in\mathscr S(\R)$ with $f'\not\equiv 0$ and $r>c\varepsilon$, by
the $f'(x+X^\varepsilon_r)\partial_r^{2}X^\varepsilon_r$ contribution
of order $\|f'\|_\infty^{2}c_\rho^{(2)}\varepsilon^{-3}$. The
$f''(\cdot)(\partial_rX^\varepsilon)^{2}$ contribution gives only the
subdominant $\varepsilon^{-2}$ scaling (Gaussian fourth moment).

\begin{lemma}[Brownian local-time Sobolev regularity]\label{lem:lt-sobolev}
Let $Y$ be standard Brownian motion, $t>0$, and write
$L\coloneqq L_t^{\cdot}(|Y|)$ for the spatial local time of the
reflected Brownian motion $|Y|$ on $[0,t]$, supported in $[0,M_t]$
with $M_t=\sup_{[0,t]}|Y_s|>0$ a.s.
\begin{enumerate}
\item For every $s<\tfrac12$, $L\in H^{s}_{\mathrm{loc}}(\R)$ a.s.
\item For every compact interval $J\subset(0,\infty)$, on the event
  $\{M_t>\sup J\}$ (which has positive $\Prob_Y$-probability for
  every $t>0$ and every finite $\sup J$),
  $L\notin H^{1}_{\mathrm{loc}}(J)$ a.s.; in particular
  $L\notin H^{1}(\R)$ a.s.\ on $\{M_t>\sup J\}$.
\end{enumerate}
\end{lemma}

\begin{proof}
By Tanaka's formula,
$|Y_t|=\int_0^{t}\mathrm{sgn}(Y_s)\,dY_s+L_t^{0}(Y)$ is a continuous
semimartingale, and we may invoke the Brownian semimartingale local
time apparatus directly on $|Y|$.

(i) The local time $r\mapsto L_t^{r}(|Y|)$ is $1/2^{-}$ H\"older
continuous on $\R$ a.s.; this is the classical regularity for
continuous-semimartingale local times in the Brownian case
(see~\cite{RevuzYor}, Ch.~VI, Theorem VI.1.7 for $Y$, applied to
$|Y|$ via the Tanaka decomposition). Consequently
$L\in H^{s}_{\mathrm{loc}}(\R)$ for every $s<1/2$ a.s.

(ii) By the Bouleau--Yor theorem~\cite{BouleauYor} applied to the
continuous semimartingale $|Y|$, the spatial local time
$r\mapsto L_t^{r}(|Y|)$ has nontrivial spatial quadratic variation
on every compact subinterval of the visited range. Concretely,
along refining partitions of $J\subset(0,\infty)$,
\[
  \liminf_{|\Delta|\to 0}\sum_{[x_k,x_{k+1}]\subset J}
   \bigl(L_t^{x_{k+1}}(|Y|)-L_t^{x_k}(|Y|)\bigr)^{2}\;>\;0
\]
a.s.\ on $\{M_t>\sup J\}$. Only the nontriviality of the spatial
quadratic variation is used in what follows; no exact value of the
limit is required. If $L\in H^{1}_{\mathrm{loc}}(J)$ on a
positive-$\Prob_Y$-probability subset of $\{M_t>\sup J\}$, then on
that subset $L$ is absolutely continuous on $J$ with
$L^{2}$-derivative; Cauchy--Schwarz then gives
\[
  \sum_{[x_k,x_{k+1}]\subset J}\bigl(L_t^{x_{k+1}}-L_t^{x_k}\bigr)^{2}
   \;\le\;|\Delta|\,\|L'\|_{L^{2}(J)}^{2}\;\to\;0
\]
along refining partitions, contradicting the displayed nontrivial
liminf. Hence $L\notin H^{1}_{\mathrm{loc}}(J)$ a.s.\ on
$\{M_t>\sup J\}$.
\end{proof}

\begin{theorem}[$L^{2}$-divergence of the first-chaos projection]
\label{prop:var-divergence}
Let $f\in\mathscr S(\R)$, $x\in\R$, $t>0$. Assume the
nondegeneracy condition
\begin{equation}\label{eq:nondegen}
  C(\eta)\;\coloneqq\;\E[f'(x+X_\eta)]\;=\;(P_\eta f')(x),
  \qquad\eta>0,\qquad\text{is not identically zero on }(0,\infty).
\end{equation}
Then
\begin{equation}\label{eq:var-diverges}
  \liminf_{\varepsilon\to 0}\,
   \Var\bigl(\mathrm{Proj}_1^{X}\,C_{\IBM,\varepsilon}(f;t)\bigr)\;=\;+\infty,
\end{equation}
and in particular $C_{\IBM,\varepsilon}(f;t)$ does not converge in
$L^{2}(\Omega)$ as $\varepsilon\to 0$. The hypothesis
\eqref{eq:nondegen} holds in particular whenever $f'(x)\ne 0$
(since $C(\eta)\to f'(x)$ as $\eta\to 0^{+}$).
\end{theorem}

\begin{proof}
Set $U_r=X^\varepsilon_r$, $V_r=\partial_rX^\varepsilon_r$,
$W_r=\partial_r^{2}X^\varepsilon_r$, all centred Gaussians in
$\Hilb_1^{X}$. In the bulk ($r>c\varepsilon$),
$\sigma_V^{2}=c_\rho^{(1)}/\varepsilon$, $\Cov(U_r,V_r)=\tfrac12$,
and
$\Cov(U_r,W_r)=-\sigma_V^{2}=-c_\rho^{(1)}/\varepsilon$ via
$\partial_r\Cov(U_r,V_r)=\sigma_V^{2}+\Cov(U_r,W_r)=0$. We use the
white-noise representation of the two-sided Brownian scenery: each
of $U_r$, $V_r$, $W_r$ is $W(\cdot)$ for a deterministic kernel in
$L^{2}(\R,dw)$, with $D_wV_r=\rho_\varepsilon(r-w)$,
$D_wW_r=\rho_\varepsilon'(r-w)$, and $D_wU_r=\xi_{r,\varepsilon}(w)$,
the smoothed-Heaviside kernel of $U_r$ on the two-sided scenery.

\emph{Step 1: First-chaos kernel decomposition.}
The pathwise correction $C_{\IBM,\varepsilon}(f;t)$ has two
summands: the boundary local-time piece $g_\varepsilon'(0)L_t^{0}(Y)$
and the interior integral
$\tfrac12\!\int_0^{t}g_\varepsilon''(|Y_s|)\,ds$. Since
$g_\varepsilon'(0)=f'(x+X_0^\varepsilon)V_0$ is the product of a
smooth function of one first-chaos element with another first-chaos
element, $f'(x+X_0^\varepsilon)$ has a non-trivial zeroth-chaos
component $\E[f'(x+X_0^\varepsilon)]$, and the boundary piece
\emph{does} contribute to $\mathrm{Proj}_1^{X}$ in general. We do
not discard it. Conditional on $Y$, the Malliavin kernel of
$\mathrm{Proj}_1^{X}C_{\IBM,\varepsilon}(f;t)$ splits as
\begin{equation}\label{eq:Kdecomp}
  K_\varepsilon(w\mid Y)
   \;=\;K_\varepsilon^{\mathrm{int}}(w\mid Y)
   \;+\;K_\varepsilon^{\mathrm{bdy}}(w\mid Y),
\end{equation}
with $K_\varepsilon^{\mathrm{bdy}}(w\mid Y)=L_t^{0}(Y)\,\E_X[D_wg_\varepsilon'(0)]$
($Y$-measurable scalar times an $X$-deterministic kernel) and
$K_\varepsilon^{\mathrm{int}}(w\mid Y)=\E_X[D_w\tfrac12\!\int_0^{t}g_\varepsilon''(|Y_s|)\,ds\mid Y]$.

\emph{Step 2: Spatial support of $K_\varepsilon^{\mathrm{bdy}}$.}
From $g_\varepsilon'(0)=f'(x+X_0^\varepsilon)V_0$,
\[
  \E_X[D_wg_\varepsilon'(0)]
   \;=\;\E_X[f''(x+X_0^\varepsilon)V_0]\,D_wX_0^\varepsilon
   \;+\;\E_X[f'(x+X_0^\varepsilon)]\,D_wV_0.
\]
Both kernels are supported in $\{|w|\le c\varepsilon\}$: indeed,
$D_wV_0=\rho_\varepsilon(0-w)=\rho_\varepsilon(w)$ has compact
support in $|w|\le c\varepsilon$, and
$D_wX_0^\varepsilon=\xi_{0,\varepsilon}(w)=\int\rho_\varepsilon(z)\eta_z(w)\,dz$
(with $\eta_z(w)=\mathrm{sgn}(z)\mathbf 1_{w\in[0,z]_{*}}$) vanishes
for $|w|>c\varepsilon$ because $\rho_\varepsilon(z)$ is supported in
$|z|\le c\varepsilon$ and $\eta_z(w)=0$ whenever $|w|>|z|$. Hence
\begin{equation}\label{eq:Kbdy-support}
  \mathrm{supp}\,K_\varepsilon^{\mathrm{bdy}}(\cdot\mid Y)
   \;\subseteq\;[-c\varepsilon,c\varepsilon],\qquad\forall\,\varepsilon>0.
\end{equation}

\emph{Step 3: Interior kernel via Stein/Malliavin.}
For $g_\varepsilon''(r)=f''(x+U_r)V_r^{2}+f'(x+U_r)W_r$,
\begin{equation}\label{eq:Dgpp}
\begin{aligned}
  \E[D_wg_\varepsilon''(r)]
   &=\bigl(\E[f'''(x+U_r)V_r^{2}]+\E[f''(x+U_r)W_r]\bigr)\xi_{r,\varepsilon}(w)\\
   &\quad+2\E[f''(x+U_r)V_r]\,\rho_\varepsilon(r-w)
   +\E[f'(x+U_r)]\,\rho_\varepsilon'(r-w).
\end{aligned}
\end{equation}
Stein for jointly Gaussian $(U,V)$ gives, for any smooth $h$,
$\E[h(x+U)V]=\Cov(U,V)\E[h'(x+U)]$ and
$\E[h(x+U)V^{2}]=\sigma_V^{2}\E[h(x+U)]+\Cov(U,V)^{2}\E[h''(x+U)]$.
In the bulk,
\begin{align*}
  \E[f'''(x+U_r)V_r^{2}]+\E[f''(x+U_r)W_r]
   &=(\sigma_V^{2}+\Cov(U_r,W_r))\E[f'''(x+U_r)]\\
   &\qquad+\tfrac14\E[f^{(5)}(x+U_r)]
   =\tfrac14\E[f^{(5)}(x+U_r)],
\end{align*}
$2\E[f''(x+U_r)V_r]=\E[f'''(x+U_r)]$, so the two $\varepsilon^{-1}$
contributions cancel exactly. Writing
$A_\varepsilon(r)=\tfrac14\E[f^{(5)}(x+U_r)]$,
$B_\varepsilon(r)=\E[f'''(x+U_r)]$,
$C_\varepsilon(r)=\E[f'(x+U_r)]$
(bounded uniformly in $(r,\varepsilon)$, with limits $A,B,C$, in
particular $C(r)=(P_r f')(x)$ as in \eqref{eq:nondegen}),
\begin{equation}\label{eq:Dkernel}
  \E[D_wg_\varepsilon''(r)]
  =A_\varepsilon(r)\xi_{r,\varepsilon}(w)
   +B_\varepsilon(r)\rho_\varepsilon(r-w)
   +C_\varepsilon(r)\rho_\varepsilon'(r-w).
\end{equation}
The $O(\varepsilon)$ boundary-layer corrections at $r\in[0,c\varepsilon]$
contribute $O(\varepsilon^{1/2})$ in $L^{2}(\R,dw)$-norm below and
do not affect the conclusion. Setting $L\coloneqq L_t^{\cdot}(|Y|)$
(supported in $[0,M_t]$ with $M_t=\sup_{[0,t]}|Y_s|>0$ a.s.),
integration against $L\,dr$ gives
\begin{equation}\label{eq:T123}
  K_\varepsilon^{\mathrm{int}}(w\mid Y)
   =\tfrac12\bigl(T_1+T_2+T_3\bigr)(w),
\end{equation}
where
$T_1(w)=\!\int_0^{\infty}\!A_\varepsilon(r)L(r)\xi_{r,\varepsilon}(w)\,dr$,
$T_2(w)=(\rho_\varepsilon*B_\varepsilon L)(w)$,
$T_3(w)=-\partial_w(\rho_\varepsilon*C_\varepsilon L)(w)$, with
$\|T_j\|_{L^{2}(\R,dw)}^{2}=\Var_X(I_j^\varepsilon\mid Y)$ for
$I_j^\varepsilon$ the corresponding random variables
($I_1^\varepsilon=\int AL\,U_r\,dr$,
$I_2^\varepsilon=\int BL\,V_r\,dr$,
$I_3^\varepsilon=\int CL\,W_r\,dr$).

\emph{Step 4: Spatial localisation away from the boundary.}
By hypothesis \eqref{eq:nondegen} and the real-analyticity of
$C(\eta)=(P_\eta f')(x)$ on $(0,\infty)$, the zero set of $C$ is
discrete; choose a compact interval $J=[a,b]\subset(0,\infty)$ with
$a>0$, $\inf_{\eta\in J}|C(\eta)|\ge c_J>0$, and a non-empty closed
subinterval $J^{\circ}=[a',b']\subset(a,b)$. Fix a smooth cutoff
$\chi\in C_c^{\infty}(\R)$ with $\chi\equiv 1$ on $J^{\circ}$,
$\mathrm{supp}\,\chi\subset J$, and $0\le\chi\le 1$. Restrict to
$\varepsilon$ small enough that $c\varepsilon<a/2$. Then, by Step~2,
$\mathrm{supp}\,K_\varepsilon^{\mathrm{bdy}}\cap\mathrm{supp}\,\chi=\emptyset$,
so $\chi\cdot K_\varepsilon^{\mathrm{bdy}}\equiv 0$ and, by
\eqref{eq:Kdecomp},
\begin{equation}\label{eq:localised}
  \Var_X(\mathrm{Proj}_1^{X}C_{\IBM,\varepsilon}\mid Y)
  =\|K_\varepsilon(\cdot\mid Y)\|_{L^{2}(\R)}^{2}
  \;\ge\;\|\chi K_\varepsilon^{\mathrm{int}}(\cdot\mid Y)\|_{L^{2}(\R)}^{2}
  =\tfrac14\|\chi(T_1+T_2+T_3)\|_{L^{2}}^{2}.
\end{equation}

\emph{Step 5: Localised $\|\chi T_1\|$ and $\|\chi T_2\|$ are bounded
uniformly.}
Since $|\chi|\le 1$, $\|\chi T_j\|_{L^{2}}\le\|T_j\|_{L^{2}}$, and:
$\|T_1\|_{L^{2}}^{2}=\Var_X(I_1^\varepsilon\mid Y)
\le(\|A\|_\infty t)^{2}(M_t+c\varepsilon)$ via
$|\Cov(U_r,U_{r'})|\le M_t+c\varepsilon$ and
$\int L=t$; $\|T_2\|_{L^{2}}^{2}=\|\rho_\varepsilon*B_\varepsilon L\|_{L^{2}}^{2}
\le\|B_\varepsilon L\|_{L^{2}}^{2}\le\|B\|_\infty^{2}\|L\|_\infty^{2}M_t$
by Young's inequality. Both bounds are uniform in $\varepsilon$.

\emph{Step 6: $\|\chi T_3\|\to\infty$ on $\{M_t>b\}$.}
Set $u_\varepsilon\coloneqq\rho_\varepsilon*(C_\varepsilon L)$. Since
$C_\varepsilon\to C$ uniformly on compact subsets of $(0,\infty)$
(heat-semigroup continuity in $\eta$), $L$ has compact support in
$[0,M_t]$, and $\|\rho_\varepsilon*f-f\|_{L^{2}}\to 0$ for
$f\in L^{2}$, we obtain $u_\varepsilon\to CL$ in $L^{2}(\R)$, and
hence $\chi u_\varepsilon\to\chi CL$ in $L^{2}(\R)$.

\emph{Claim: $\chi CL\notin H^{1}(\R)$ a.s.\ on $\{M_t>b\}$.}
On $J^{\circ}$, $\chi\equiv 1$ and $|C|\ge c_J>0$; if
$\chi CL\in H^{1}(\R)$, then $\chi CL|_{J^{\circ}}=CL|_{J^{\circ}}
\in H^{1}_{\mathrm{loc}}(J^{\circ})$, and dividing by the smooth
nonvanishing $C$ on $J^{\circ}$ gives $L\in H^{1}_{\mathrm{loc}}(J^{\circ})$.
But \cref{lem:lt-sobolev}(ii) applied to the compact subinterval
$J^{\circ}$ yields $L\notin H^{1}_{\mathrm{loc}}(J^{\circ})$ a.s.\ on
$\{M_t>\sup J^{\circ}\}=\{M_t>b'\}\supseteq\{M_t>b\}$, a
contradiction. Hence $\chi CL\notin H^{1}(\R)$ a.s.\ on $\{M_t>b\}$.

For any $R>0$, $\widehat{\chi u_\varepsilon}\to\widehat{\chi CL}$
in $L^{2}(\R)$ by Plancherel transfer of $L^{2}$-convergence; with
$|\widehat{\chi u_\varepsilon}(k)|\le\|\chi u_\varepsilon\|_{L^{1}(\R)}\le
\|\chi\|_{L^{1}}\sup_\varepsilon\|u_\varepsilon\|_{L^{\infty}}<\infty$
uniform in $\varepsilon,k$, dominated convergence on $\{|k|\le R\}$
gives
\[
  \liminf_{\varepsilon\to 0}\|(\chi u_\varepsilon)'\|_{L^{2}}^{2}
   \;=\;\liminf_{\varepsilon\to 0}\frac{1}{2\pi}\!\int_\R k^{2}|\widehat{\chi u_\varepsilon}|^{2}dk
   \;\ge\;\frac{1}{2\pi}\!\int_{|k|\le R}\!k^{2}|\widehat{\chi CL}|^{2}dk.
\]
Letting $R\to\infty$ and using $\chi CL\notin H^{1}$ gives
$\liminf_{\varepsilon\to 0}\|(\chi u_\varepsilon)'\|_{L^{2}}=+\infty$
a.s.\ on $\{M_t>b\}$. Now
$\chi T_3=-\chi\,u_\varepsilon'=-(\chi u_\varepsilon)'+\chi'u_\varepsilon$,
and $\|\chi'u_\varepsilon\|_{L^{2}}\to\|\chi'CL\|_{L^{2}}<\infty$
(smooth $\chi'$ times $L^{2}$-convergent $u_\varepsilon$). Therefore
\[
  \|\chi T_3\|_{L^{2}}\;\ge\;\|(\chi u_\varepsilon)'\|_{L^{2}}-\|\chi'u_\varepsilon\|_{L^{2}}\;\xrightarrow[\varepsilon\to 0]{}\;+\infty
  \quad\text{a.s.\ on }\{M_t>b\}.
\]

\emph{Step 7: Conclusion.}
By the reverse triangle inequality applied to \eqref{eq:localised}
and Steps~5--6,
\[
  \|K_\varepsilon\|_{L^{2}}
   \;\ge\;\tfrac12\,\|\chi(T_1+T_2+T_3)\|_{L^{2}}
   \;\ge\;\tfrac12\bigl(\|\chi T_3\|_{L^{2}}-\|\chi T_1\|_{L^{2}}-\|\chi T_2\|_{L^{2}}\bigr)
   \;\xrightarrow[\varepsilon\to 0]{}\;+\infty
\]
a.s.\ on $\{M_t>b\}$, which has positive $\Prob_Y$-probability for
every $t>0$ (e.g.\ because $\{Y_t>b\}\subset\{M_t>b\}$ and
$\Prob(Y_t>b)>0$). Hence
$\Var_X(\mathrm{Proj}_1^{X}C_{\IBM,\varepsilon}\mid Y)\to+\infty$
a.s.\ on $\{M_t>b\}$, and Fatou over $Y$ gives
$\liminf\E_Y[\Var_X(\cdot\mid Y)]=+\infty$; since
$\Var(\mathrm{Proj}_1^{X}C_{\IBM,\varepsilon})\ge\E_Y[\Var_X(\cdot\mid Y)]$,
\eqref{eq:var-diverges} follows. $L^{2}(\Omega)$-convergence of
$C_{\IBM,\varepsilon}(f;t)$ would force a uniformly bounded variance
and so fails.
\end{proof}

The key analytic point of \cref{prop:var-divergence} is twofold.
First, the naive $\varepsilon^{-1}$-divergent first-chaos
contributions from the two summands of $g_\varepsilon''(r)$
\emph{cancel exactly} in the Malliavin kernel \eqref{eq:Dkernel},
leaving a residual $T_3$ summand with $\rho_\varepsilon'$-kernel and
coefficient $C_\varepsilon L$. Second, a spatial cutoff $\chi$
supported away from the boundary $w=0$ suppresses the first-chaos
contribution of the local-time term $g_\varepsilon'(0)L_t^{0}(Y)$
(whose Malliavin kernel is supported in $\{|w|\le c\varepsilon\}$,
\eqref{eq:Kbdy-support}) and isolates the interior $T_3$ piece; the
localised divergence
$\|\chi T_3\|_{L^{2}}\to\infty$ is driven by
$\chi CL\notin H^{1}(\R)$, in turn forced by the Bouleau--Yor spatial
quadratic variation of the clock local time
(\cref{lem:lt-sobolev}).

\paragraph{The obstruction is quartic, not merely topological.}
The divergence above is not only a topology issue; it reflects the
\emph{order} of the IBM observable energy. Brownian motion has
quadratic energy and produces the second-order heat operator
$\tfrac12\Delta$. IBM composes a Brownian scenery with a Brownian
clock, so the observable correction $C_\IBM$ is quartic at the
energy level: $C_{\IBM,\varepsilon}$ is second order in the clock
variable (via $\tfrac12\int g_\varepsilon''(|Y_s|)\,ds$), but because
its coefficients are derivatives of the Brownian scenery the
semigroup shadow of the integrated correction is fourth order in the
spatial variable. The appearance of $\Delta^{2}u$ in
\eqref{eq:bt-correction} is the PDE shadow of this quartic composite
observable energy, not of a second-order Markov generator.
Accordingly, the candidate second-order observable energy
$\Energy_\IBM^{(2)}$ of \cref{def:second-order-ibm} should be read as
a second-order \emph{operator-derivative} object whose
$\Carre^{(2)}$-form, once averaged against the clock, produces a
fourth-order / quartic PDE-shadow signature. The
$\partial_r^{2}X^\varepsilon$ singularity isolated by
\cref{prop:var-divergence} is the rank-one footprint of this
quartic structure in the first scenery chaos.

\subsection*{Energy-contraction renormalisation candidate}

The subtraction below should not be read as intrinsically Gaussian.
At the operator level, the divergent object is the
\emph{diagonal contraction} of two first-order energy directions: the
self-pairing of the energy-derivative kernel
$\partial_rX^\varepsilon_r$ against itself, measured in the source
energy space. In any Gaussian source atlas, this contraction is
realised as the expectation of the squared first-chaos element. In the
present Brownian-scenery realisation, the diagonal contraction is
\begin{equation}\label{eq:contraction}
  \mathrm{Contr}_\varepsilon(r)
  \;\coloneqq\;\E\bigl[(\partial_rX^\varepsilon_r)^{2}\bigr]
  \;=\;\frac{c_\rho^{(1)}}{\varepsilon}
  \qquad\text{for }|r|>c\varepsilon,
\end{equation}
so Wick / Hermite centering is precisely the Brownian-coordinate
expression of an operator-energy contraction subtraction. We use the
operator language below; the Gaussian Wick formula \eqref{eq:wick-sq}
is the explicit coordinate realisation in this paper.

Define the contraction-renormalised square
\begin{equation}\label{eq:wick-sq}
  :(\partial_rX^\varepsilon_r)^{2}:\;\coloneqq\;
  (\partial_rX^\varepsilon_r)^{2}\,-\,\mathrm{Contr}_\varepsilon(r)
  \;=\;(\partial_rX^\varepsilon_r)^{2}\,-\,\frac{c_\rho^{(1)}}{\varepsilon};
\end{equation}
in the present Gaussian source coordinates this is the usual
Wick / Hermite square in $\Hilb_2^{X}$. The
\emph{contraction-renormalised} second-derivative density is
\begin{equation}\label{eq:wick-g}
  :g_\varepsilon''(r):\;\coloneqq\;
   f''(x+X^\varepsilon_r):(\partial_rX^\varepsilon_r)^{2}:
   \;+\;f'(x+X^\varepsilon_r)\,\partial_r^{2}X^\varepsilon_r,
\end{equation}
and the \emph{operator-contraction-renormalised correction candidate}
is
\begin{equation}\label{eq:wick-correction}
  \widehat C_{\IBM,\varepsilon}(f;t)
  \;\coloneqq\;
   g_\varepsilon'(0)\,L_t^{0}(Y)\;+\;\tfrac12\!\int_0^t\,
   :g_\varepsilon''(|Y_s|):\,ds.
\end{equation}
By construction, $:g_\varepsilon''(r):$ has the contraction
$\mathrm{Contr}_\varepsilon(r)$ subtracted from the square component;
its expectation differs from that of the raw $g_\varepsilon''(r)$ by the
explicit deterministic counterterm
$-\mathrm{Contr}_\varepsilon(r)\,\E_X[f''(x+X^\varepsilon_r)]$.

We do not prove convergence of $\widehat C_{\IBM,\varepsilon}$;
\eqref{eq:wick-correction} is offered as the natural
contraction-renormalised candidate. Whether
$\widehat C_{\IBM,\varepsilon}(f;t)$ converges in $L^{2}(\Omega)$, in
some weaker topology, or only chaos-by-chaos remains open.

\subsection*{Source-chaos coordinates}

The source-resolved chaos atlas of~\cite{Fontes-IBMchaos} suggests
where convergence of $\widehat C_{\IBM,\varepsilon}$ should be sought.
The Wick square $:(\partial_rX^\varepsilon)^{2}:$ lives in the second
scenery Wiener chaos $\Hilb^{X}_{2}$; the kernel
$\partial_r^{2}X^\varepsilon$ lives in $\Hilb^{X}_{1}$; multiplication
by smooth functions of $X^\varepsilon$ couples to other scenery
sectors; integration against the half-normal weight of $|Y_s|$ couples
to the clock sectors $\Hilb^{Y}_{n}$. The bidegree projections of
$\widehat C_{\IBM,\varepsilon}$ onto $(2,0)$, $(2,1)$, $(1,1)$,
$(1,2)$, $(0,2)$ are the natural objects in which
$\varepsilon\to 0$ convergence might be checked first, in line with
the consistency requirements of \cref{rem:bidegree-2}.

\subsection*{Status and next analytic target}

\begin{itemize}
\item \emph{Raw centred correction $\tilde C_{\IBM,\varepsilon}$}:
  $\Var(\mathrm{Proj}_1^{X}\,C_{\IBM,\varepsilon}(f;t))\to+\infty$ as
  $\varepsilon\to 0$ \textbf{(proved)}, for every $f\in\mathscr S(\R)$
  with $(P_\cdot f')(x)\not\equiv 0$ on $(0,\infty)$
  (\cref{prop:var-divergence}). The raw $L^{2}(\Omega)$-target for
  the unmollified pathwise correction is therefore ruled out at the
  first scenery-chaos level.
\item \emph{Contraction-renormalised candidate
  $\widehat C_{\IBM,\varepsilon}$}: \textbf{defined}
  (\eqref{eq:wick-correction}). Its $L^{2}$-convergence is
  \textbf{open}; the Wick / Hermite subtraction acts in
  $\Hilb_0^{X}\oplus\Hilb_2^{X}$ and structurally does not touch the
  divergent $\rho_\varepsilon'$-kernel of \eqref{eq:Dkernel}
  (\cref{prop:wick-insufficient}).
\item \emph{Finite-bidegree projections of
  $\widehat C_{\IBM,\varepsilon}$}: convergence at each
  $(m,n)$-sector is the natural next analytic step, drawing on the
  kernels of~\cite{Fontes-IBMchaos}.
\item \emph{Unmollified pathwise correction $C_\IBM(f;t)$}: open. Any
  future definition must (i) reduce in full expectation to
  $C_\IBM^{\mathrm{sg}}(f;t,x)$ (\cref{cor:pathwise-to-sg}) and
  (ii) be constructed in a topology weaker than $L^{2}(\Omega)$,
  since the raw $L^{2}$-limit of $C_{\IBM,\varepsilon}$ does not
  exist (\cref{prop:var-divergence}); a quartic observable-energy
  dual is the structural target (\eqref{eq:rigged-4}).
\end{itemize}

\section{Renormalised pathwise correction: variance after Gaussian contraction subtraction}\label{sec:renorm-correction}

The contraction-renormalised candidate $\widehat C_{\IBM,\varepsilon}$
of \eqref{eq:wick-correction} subtracts the diagonal contraction
$\mathrm{Contr}_\varepsilon(r)=c_\rho^{(1)}/\varepsilon$ from the
squared first-energy direction $(\partial_rX^\varepsilon)^{2}$ ---
realised concretely in the present Gaussian source atlas as the
Wick / Hermite square subtraction \eqref{eq:wick-sq}. We now check
whether that subtraction is sufficient to give an $L^{2}$-bounded
candidate, and conclude that it is not: the linear-in-Gaussian term
$\partial_r^{2}X^\varepsilon$ lives in the first scenery chaos
$\Hilb_1^{X}$ with $\Var\sim\varepsilon^{-3}$, and no subtraction of
products of Gaussian first-chaos elements lowers its rank-1 chaos
mass. The unmollified pathwise correction $C_\IBM(f;t)$ must therefore
be constructed in a topology weaker than $L^{2}(\Omega)$, or
chaos-by-chaos with test integrands.

\subsection*{Cross-covariance audit at the bulk}

For $|r|>c\varepsilon$, the bulk covariance matrix of the triple
$(X^\varepsilon_r,\,\partial_rX^\varepsilon_r,\,\partial_r^{2}X^\varepsilon_r)$
is
\[
  \mathrm{Cov}\;=\;
  \begin{pmatrix}
   r+O(\varepsilon) & 1/2 & -c_\rho^{(1)}/\varepsilon\\[2pt]
   1/2 & c_\rho^{(1)}/\varepsilon & 0\\[2pt]
   -c_\rho^{(1)}/\varepsilon & 0 & c_\rho^{(2)}/\varepsilon^{3}
  \end{pmatrix},
\]
the only $\varepsilon^{-1}$ cross-correlation being
\[
  \Cov(X^\varepsilon_r,\partial_r^{2}X^\varepsilon_r)
   \;=\;-\,c_\rho^{(1)}/\varepsilon,
\]
obtained by twice differentiating
$\Var(X^\varepsilon_r)=\sigma_\varepsilon^{2}(r)$ and using
$(\sigma_\varepsilon^{2})'(r)=1$ for $|r|>c\varepsilon$.
This cross-covariance contributes the divergent piece
$-c_\rho^{(1)}\varepsilon^{-1}\E[f''(x+X^\varepsilon_r)]$ to
$\E[f'(\cdot)\partial_r^{2}X^\varepsilon]$; in the unrenormalised
$g_\varepsilon''$ it is cancelled exactly by the $\sigma_V^{2}$ part
of $\E[f''(\cdot)V^{2}]$, leaving
$\E[g_\varepsilon''(r)]=\tfrac14\E[f^{(4)}(x+X^\varepsilon_r)]$
finite. The Wick / Hermite subtraction in $\mathcal R_\varepsilon$
(the contraction subtraction in the Gaussian source coordinates)
removes the $\sigma_V^{2}$ part of $(\partial_rX^\varepsilon)^{2}$
\emph{only}, so an additional deterministic counterterm
$+\mathrm{Contr}_\varepsilon(r)\,\E[f''(x+X^\varepsilon_r)]$ is
needed to preserve $\E[\mathcal R_\varepsilon]=\E[g_\varepsilon'']$.
The cross-covariance produces no stochastic counterterm that would
tame the $\varepsilon^{-3}$ variance of
$\partial_r^{2}X^\varepsilon$.

\subsection*{Variance after Gaussian contraction subtraction}

\begin{proposition}[Gaussian square-contraction subtraction is insufficient]
\label{prop:wick-insufficient}
Under the hypotheses of \cref{prop:var-divergence}, the contraction
subtraction
$(\partial_rX^\varepsilon)^{2}\mapsto:(\partial_rX^\varepsilon)^{2}:$
does not modify the $\partial_r^{2}X^\varepsilon$ content of the
integrand $g_\varepsilon''$. Consequently the
first-scenery-chaos singularity in
$\mathrm{Proj}_1^{X}C_{\IBM,\varepsilon}$ carried by
$\partial_r^{2}X^\varepsilon$ (identified in
\cref{prop:var-divergence}) is not removed by the subtraction, and
the contraction renormalisation is necessary but not sufficient for
$L^{2}$-tightness of $\widehat C_{\IBM,\varepsilon}$. In the Gaussian
Brownian-scenery source atlas, the subtraction is the Wick / Hermite
square centering used in the present paper.
\end{proposition}

\begin{proof}
The contraction subtraction
$(\partial_rX^\varepsilon_r)^{2}\mapsto:(\partial_rX^\varepsilon_r)^{2}:
=(\partial_rX^\varepsilon_r)^{2}-c_\rho^{(1)}/\varepsilon$ alters only
the $V^{2}$ factor in
$g_\varepsilon''(r)=f''(x+X^\varepsilon_r)V_r^{2}+f'(x+X^\varepsilon_r)W_r$
(with $V_r=\partial_rX^\varepsilon_r$, $W_r=\partial_r^{2}X^\varepsilon_r$).
The term $f'(x+X^\varepsilon_r)W_r$ --- which carries the
$C_\varepsilon(r)\rho_\varepsilon'(r-w)$ contribution to the
first-chaos Malliavin kernel \eqref{eq:Dkernel} of
\cref{prop:var-divergence}, with coefficient
$C_\varepsilon(r)=\E_X[f'(x+X^\varepsilon_r)]=(P_{\sigma_\varepsilon^{2}(r)}f')(x)$,
and is the source of the residual first-chaos $L^{2}$-divergence
identified in Step~4 of the proof of \cref{prop:var-divergence} ---
is untouched by the contraction subtraction.

The subtraction does add a counterterm to $\widehat C_{\IBM,\varepsilon}$
of the form $-\tfrac{c_\rho^{(1)}}{2\varepsilon}\int_0^t
f''(x+X^\varepsilon_{|Y_s|})\,ds$, a smooth functional of $X^\varepsilon$
alone. Its first-chaos kernel (via Stein on $f''(x+X^\varepsilon)$)
is built from the $U_r$-fiber $\xi_{r,\varepsilon}(w)$ of
$\Hilb_1^{X}$, not from the $W_r$-fiber $\rho_\varepsilon'(r-w)$; in
the notation of \eqref{eq:Dkernel} and \eqref{eq:T123}, it modifies
$A_\varepsilon$ and contributes to the $T_1$-summand, but does not
modify $T_3=-\partial_w(\rho_\varepsilon*C_\varepsilon L)$, which is
the localized divergent piece of \cref{prop:var-divergence}. Hence
the residual $L^{2}$-divergence of
$\mathrm{Proj}_1^{X}\widehat C_{\IBM,\varepsilon}$ persists; the
contraction subtraction acts in the
$\Hilb_0^{X}\oplus\Hilb_2^{X}$ even-chaos sectors and does not
address the rank-$1$ chaos mass carried by the $\rho_\varepsilon'$-kernel
in \eqref{eq:Dkernel}.
\end{proof}

\paragraph{Caveat on $L^{2}$-divergence of $\widehat C_{\IBM,\varepsilon}$.}
\Cref{prop:wick-insufficient} shows that the Wick / Hermite
subtraction cannot cancel the $\rho_\varepsilon'$-driven $T_3$
contribution to the localised first-chaos kernel
\eqref{eq:T123}, but it does not by itself prove
$\Var(\widehat C_{\IBM,\varepsilon})\to\infty$ unconditionally:
an algebraic cancellation between the residual $T_3$ piece and the
Stein counterterm generated by the contraction subtraction would
require structural conditions on $f$ beyond \eqref{eq:nondegen},
which we do not analyse. The structural content is that the
contraction subtraction acts in the wrong chaos sector to address
the $\rho_\varepsilon'(r-w)$ kernel, so no $L^{2}$-tightness can be
expected from it alone.

\paragraph{What is needed instead.}
\Cref{prop:wick-insufficient} identifies the analytic obstacle
precisely: $\partial_r^{2}X^\varepsilon$ is a centred Gaussian in
$\Hilb_1^{X}$ whose variance tracks the spatial-derivative-of-white-noise
it converges to as $\varepsilon\to 0$. No contraction subtraction of
squared first-energy directions (no Wick subtraction of Gaussian
squares, in the present source coordinates) addresses this rank-one
chaos mass. The renormalisation must operate at the level of a
\emph{quartic composite observable energy}, not at the level of a
single Gaussian square: the second-order clock derivative
$g_\varepsilon''(|Y_s|)$, paired with second-order scenery covariances,
produces the fourth-order $\Delta^{2}$ signature of \eqref{eq:bt-correction},
and the appropriate target topology is the dual of a fourth-order
energy space rather than of $L^{2}$. Schematically, with
$\Energy_\IBM^{(2)}$ the candidate second-order observable energy of
\cref{def:second-order-ibm} and $\Energy_\IBM^{(4)}=\Energy_\IBM^{(2)}\otimes
\Energy_\IBM^{(2)}$ its tensor square (the quartic composite energy
underlying the $\Delta^{2}$ shadow), one expects a rigged hierarchy
\begin{equation}\label{eq:rigged-4}
  \Energy_\IBM^{(4),+}
   \;\subset\;\Energy_\IBM^{(4)}
   \;\subset\;\Energy_\IBM^{(4),-},
\end{equation}
in which the unmollified pathwise correction $C_\IBM(f;t)$ should
live in the negative space $\Energy_\IBM^{(4),-}$ and be paired with
test integrands from $\Energy_\IBM^{(4),+}$. Two analytic routes
remain to construct such a pairing:
\begin{enumerate}
\item[(i)] interpret $\partial_r^{2}X$ in the unmollified limit as a
  Hida / Schwartz distribution and read the
  $f'(X)\partial_r^{2}X$ contribution as a distribution-valued kernel
  pairing with test integrands in $\Energy_\IBM^{(4),+}$ --- the
  fourth-order analogue of the Tanaka-style interpretation already
  required for the unmollified clock derivative in
  \cref{cor:unmollified};
\item[(ii)] Stratonovich-regularise the time integral
  $\int_0^t f'(x+X^\varepsilon_{|Y_s|})\,\partial_r^{2}
   X^\varepsilon_{|Y_s|}\,ds$ via integration by parts in the clock
  variable, transferring the spatial $r$-derivative onto
  Brownian-clock local-time data and producing a quartic-energy
  pairing.
\end{enumerate}
Neither route is pursued here. The renormalisation programme is
therefore not Wick-subtraction of a second-chaos object, but
renormalisation of a \emph{quartic composite observable energy} whose
PDE shadow is $\Delta^{2}$.

\subsection*{Finite-bidegree projections: anisotropy of the divergence}

The source-resolved chaos atlas of~\cite{Fontes-IBMchaos} projects
$\widehat C_{\IBM,\varepsilon}$ onto sectors
$\Hilb_m^{X}\otimes\Hilb_n^{Y}$. The proof of
\cref{prop:wick-insufficient} shows that the divergence is
\emph{anisotropic} across bidegrees:
\begin{itemize}
\item \emph{$(2,n)$-bidegrees}, which contain the contraction-renormalised
  square $:(\partial_rX^\varepsilon)^{2}:$ in $\Hilb_2^{X}$ --- the
  Wick / Hermite square in the Gaussian source atlas: conjectured
  $L^{2}$-bounded and convergent chaos-by-chaos as
  $\varepsilon\to 0$, in line with the kernel-by-kernel framework of
  \cite{Fontes-IBMchaos}.
\item \emph{$(1,n)$-bidegrees}, which contain $\partial_r^{2}X^\varepsilon$
  in $\Hilb_1^{X}$: structurally unaddressed by the contraction
  subtraction (\cref{prop:wick-insufficient}); convergence
  expected in the fourth-order / quartic observable-energy dual
  topology $\Energy_\IBM^{(4),-}$ of \eqref{eq:rigged-4}, of which
  Hida-distribution and Stratonovich-regularised topologies are
  candidate realisations.
\end{itemize}
We do not prove either statement; together they identify the precise
location of the analytic obstacle in the chaos atlas.

\subsection*{Semigroup consistency is preserved}

Whatever topology eventually houses $C_\IBM(f;t)$, the consistency
\eqref{eq:consistency} from \cref{cor:pathwise-to-sg} forces
$\E[C_\IBM(f;t)]=C_\IBM^{\mathrm{sg}}(f;t,x)$. The schematic
decomposition
\begin{equation}\label{eq:final-decomp}
  C_\IBM(f;t)\;=\;C_\IBM^{\mathrm{sg}}(f;t,x)
   \;+\;C_\IBM^{\mathrm{fluc}}(f;t),
  \qquad\E[C_\IBM^{\mathrm{fluc}}]=0,
\end{equation}
is therefore the structural form to expect of any future
construction; we do not assert that $C_\IBM^{\mathrm{fluc}}$ exists in
the present analytical setting.

\subsection*{Status}

\begin{itemize}
\item \emph{Raw centred correction $\tilde C_{\IBM,\varepsilon}$}:
  $\Var(\mathrm{Proj}_1^{X}C_{\IBM,\varepsilon}(f;t))\to+\infty$
  \textbf{(proved)} for every $f$ with $(P_\cdot f')(x)\not\equiv 0$
  (\cref{prop:var-divergence}); raw $L^{2}(\Omega)$ unmollified
  pathwise correction is ruled out.
\item \emph{Contraction-renormalised $\widehat C_{\IBM,\varepsilon}$
  (i.e., Wick / Hermite square in the Gaussian source atlas)}: the
  subtraction acts in $\Hilb_0^{X}\oplus\Hilb_2^{X}$ and structurally
  does not address the $\partial_r^{2}X^\varepsilon$ first-chaos
  singularity of \cref{prop:var-divergence}; necessary but
  insufficient (\cref{prop:wick-insufficient}).
\item \emph{$(2,n)$-bidegree projections}: conjectured $L^{2}$-bounded
  / chaos-by-chaos convergent; not proved here.
\item \emph{$(1,n)$-bidegree projections}: structurally unaddressed by
  contraction subtraction; convergence expected in the fourth-order
  observable-energy dual topology $\Energy_\IBM^{(4),-}$, realised
  concretely as Hida-distributional or Stratonovich-regularised.
\item \emph{Unmollified pathwise correction $C_\IBM(f;t)$}: open. Must
  satisfy $\E[C_\IBM]=C_\IBM^{\mathrm{sg}}$ and be constructed via a
  renormalisation of the quartic composite observable energy
  $\Energy_\IBM^{(4)}$, in the fourth-order observable-energy dual
  topology of \eqref{eq:rigged-4} (Wick / Hermite centering is the
  Brownian-coordinate realisation of the second-chaos part of this
  quartic renormalisation; the rank-one $\partial_r^{2}X^\varepsilon$
  contribution requires a genuinely fourth-order treatment).
\end{itemize}

\section{Open analytic tasks}\label{sec:tasks}

The compression construction of Section~\ref{sec:compat} supplies the
abstract framework; \cref{def:ibm-datum} supplies the abstract IBM
datum and \cref{thm:ibm-occ} supplies the first concrete instantiation
on the mollified occupation core. Carrying this through to the
intrinsic IBM It\^o formula reduces to the following analytic problems,
extending the four listed at the end of Section~\ref{sec:ibm-occ}.

\begin{enumerate}
\item \emph{Density of $\mathcal C_\IBM$.} Verify that
  $\mathcal C_\IBM=\{F\in L^{2}(\mathcal F_\IBM):\iota F\in\Dom(D_X)\cap\Dom(D_Y)\}$
  is dense in $L^{2}(\mathcal F_\IBM)$. The mollified occupation core
  $\bigcup_{\varepsilon>0}\mathcal C_{\mathrm{occ}}^{\,\varepsilon}$
  furnishes a non-empty subspace; whether its $L^{2}$-closure exhausts
  $L^{2}(\mathcal F_\IBM)$ is open, and depends on the regularity of
  Wiener compositions $X_{|Y_{t}|}$ at pointwise evaluations.

\item \emph{Concrete description of $\Energy_\IBM$.} The compression
  construction defines $\Energy_\IBM\subseteq L^{2}(\R)\oplus L^{2}(\R_{+})$
  as a closed subspace, but does not give a usable inner-product
  realisation. Candidate descriptions: occupation-time-weighted norms,
  or the spatial decomposition of IBM occupation
  functionals~\cite{Fontes-IBMchaos}. The first concrete sub-realisation
  is $\Energy_\IBM^{\mathrm{occ},\,\varepsilon}$ of
  \cref{def:occ-energy}.

\item \emph{Unmollified clock derivative.} Identify
  $\lim_{\varepsilon\to 0} D_Y F_{T,\varepsilon}(\varphi)$ as a
  distribution-valued kernel against $L^{2}(\R_{+})$ test integrands,
  via the Tanaka representation of $L_T^{\,x}(Y)$; prove closability of
  the resulting operator on $\mathcal C_{\mathrm{occ}}$. This is the
  central analytic content not yet provided by \cref{thm:ibm-occ}.

\item \emph{Cylindrical reduction for $D_\IBM$.} Establish a chain or
  product rule for $D_\IBM$ via a cylindrical reduction with random
  representers; the natural candidate is the joint Brownian--Malliavin
  representer projected onto $\mathcal F_\IBM$.

\item \emph{Second-order construction.} Verify density of
  $\mathcal C_\IBM^{(2)}$ on (the closure of) the occupation core and
  complete the construction of $\Energy_\IBM^{(2)}$ and $D_\IBM^{(2)}$
  as in Section~\ref{sec:ibm-source-2nd}.

\item \emph{Recover the PDE shadow.} Show that, averaged against the
  law of $Y$, the second-order $\Carre$-form built from
  $D_\IBM^{(2)}$ matches the bi-Laplacian/local-time PDE signature
  identified by~\cite{AllouZheng}.

\item \emph{Unmollified pathwise correction in the quartic
  observable-energy dual.} The raw $L^{2}(\Omega)$-target for the
  unmollified pathwise correction $C_\IBM(f;t)$ is ruled out by the
  proved first-chaos divergence (\cref{prop:var-divergence}). The
  remaining analytic task is to construct $C_\IBM(f;t)$ in the dual
  $\Energy_\IBM^{(4),-}$ of the quartic composite observable energy
  $\Energy_\IBM^{(4)}=\Energy_\IBM^{(2)}\otimes\Energy_\IBM^{(2)}$
  (rigged hierarchy \eqref{eq:rigged-4}), and to prove the
  corresponding energy-dual It\^o identity --- a refinement of the
  schematic \cref{conj:ibm-ito} in which $C_\IBM(f;t)$ is paired
  against test integrands from $\Energy_\IBM^{(4),+}$ rather than
  treated as an element of $L^{2}(\Omega)$. We do not assert the
  existence of such a dual topology; constructing one (e.g.\ via
  Hida-distribution or Stratonovich-regularisation routes) is the
  natural open analytic problem.

\item \emph{Reconciliation with companion papers.} Compare the
  source-resolved variance budget of~\cite{Fontes-IBMchaos} and the
  numerical decompositions of~\cite{Fontes-IBMnumerical} with the
  energy correction of \cref{conj:ibm-ito} term by term on each
  bidegree; cross-check the second-order energy against the
  PDE-generator framework of~\cite{Fontes-PDEgenerators}.
\end{enumerate}

\section*{Next analytic targets}

The viewpoint of this note: \emph{the source Malliavin calculus is the
atlas; the intrinsic operator derivative, defined by observable
compression, is the invariant object.} The ladder of results in this
paper is:
\begin{enumerate}
\item the abstract first-order observable compression
  (\cref{thm:compression,thm:proj-div}) gives $D_\Obs$ and $\delta_\Obs$
  by closing the source pullback through the observable filtration;
\item the mollified IBM occupation functionals $F_{T,\varepsilon}(\varphi)$
  of \cref{def:mollified} supply a concrete non-empty first-order IBM
  core, with \cref{thm:ibm-occ}~(i)--(ii) realising the abstract
  construction explicitly;
\item the scenery component converges in the unmollified limit:
  $F_{T,\varepsilon}\to F_T$ in $L^{2}(\Omega)$ and
  $D_XF_{T,\varepsilon}\to D_XF_T$ in $L^{2}(\Omega;L^{2}(\R))$ at rate
  $O(\varepsilon^{1/2})$ (\cref{thm:ibm-occ}~(iii)--(iv));
\item the clock component is the first singular obstruction
  (\cref{cor:unmollified}), traced to the $\varepsilon^{-1}$ pointwise
  variance of $\partial_yX^\varepsilon$ (\cref{lem:moll-eps-inv});
\item the candidate second-order observable energy
  $\Energy_\IBM^{(2)}$ of \cref{def:second-order-ibm} is the next
  operator-energy object needed for an IBM It\^o formula, with
  \cref{prop:compression-2} establishing its closability;
\item the Fourier-mode diagnostic \eqref{eq:az-general} identifies the
  PDE shadow that this second-order energy must reproduce: the
  bi-Laplacian $\tfrac18\Delta^{2}u$ together with the singular
  boundary term $\Delta f/(2\sqrt{2\pi t})$.
\end{enumerate}
The next analytic theorem is to construct the unmollified IBM
correction $C_\IBM(f;t)$ as an element of the dual
$\Energy_\IBM^{(4),-}$ of the quartic composite observable-energy
space $\Energy_\IBM^{(4)}=\Energy_\IBM^{(2)}\otimes\Energy_\IBM^{(2)}$
in the rigged hierarchy \eqref{eq:rigged-4}, and to prove the
corresponding energy-dual It\^o identity. Its expectation should
recover the integrated Brownian-time semigroup correction
$\E[C_\IBM(f;t)]=u(t,x)-f(x)=C^{\mathrm{det}}_\IBM(f;t,x)$ of
\eqref{eq:bt-benchmark}. The proved first-chaos $L^{2}$-obstruction
\cref{prop:var-divergence} rules out $C_\IBM(f;t)$ as an ordinary
$L^{2}(\Omega)$ random variable, so the construction must take
place in the quartic energy-dual topology; once available, it
would convert the schematic pathwise It\^o \cref{conj:ibm-ito} into
a theorem in that topology. The component tasks are listed in
Section~\ref{sec:tasks} and Section~\ref{sec:second-order}.

\begin{thebibliography}{99}

\bibitem{AllouZheng}
H.~Allouba and W.~Zheng,
\emph{Brownian-time processes: the PDE connection and the
half-derivative generator},
Ann. Probab. \textbf{29} (2001), no.~4, 1780--1795.

\bibitem{BogachevGaussian}
V.~I.~Bogachev,
\emph{Gaussian Measures},
Mathematical Surveys and Monographs, vol.~62,
American Mathematical Society, Providence, RI, 1998.

\bibitem{BorodinSalminen}
A.~N.~Borodin and P.~Salminen,
\emph{Handbook of Brownian Motion --- Facts and Formulae},
2nd ed., Probability and Its Applications,
Birkh\"auser Verlag, Basel, 2002.

\bibitem{BouleauHirsch}
N.~Bouleau and F.~Hirsch,
\emph{Dirichlet Forms and Analysis on Wiener Space},
de Gruyter Studies in Mathematics, vol.~14,
Walter de Gruyter \& Co., Berlin, 1991.

\bibitem{BouleauYor}
N.~Bouleau and M.~Yor,
\emph{Sur la variation quadratique des temps locaux de certaines
semimartingales},
C.~R.~Acad.~Sci.~Paris S\'er.~I Math. \textbf{292} (1981), no.~9,
491--494.

\bibitem{Fontes-IBMchaos}
R.~Fontes,
\emph{Source-resolved chaos expansions for iterated Brownian motion},
companion paper, in preparation, 2026.

\bibitem{Fontes-IBMnumerical}
R.~Fontes,
\emph{A numerical variance budget for iterated Brownian motion},
companion paper, in preparation, 2026.

\bibitem{Fontes-OperatorDerivative}
R.~Fontes,
\emph{The operator derivative in continuous stochastic calculus:
a Hilbert energy space framework},
preprint, 2026.

\bibitem{Fontes-PDEgenerators}
R.~Fontes,
\emph{PDE generators and energy corrections for Brownian-time processes},
companion paper, in preparation, 2026.

\bibitem{FukushimaOshimaTakeda}
M.~Fukushima, Y.~Oshima, and M.~Takeda,
\emph{Dirichlet Forms and Symmetric Markov Processes},
2nd ed., de Gruyter Studies in Mathematics, vol.~19,
Walter de Gruyter \& Co., Berlin, 2011.

\bibitem{Janson}
S.~Janson,
\emph{Gaussian Hilbert Spaces},
Cambridge Tracts in Mathematics, vol.~129,
Cambridge University Press, Cambridge, 1997.

\bibitem{Nualart}
D.~Nualart,
\emph{The Malliavin Calculus and Related Topics},
2nd ed., Probability and Its Applications,
Springer-Verlag, Berlin, 2006.

\bibitem{NualartPardoux}
D.~Nualart and E.~Pardoux,
\emph{Stochastic calculus with anticipating integrands},
Probab. Theory Related Fields \textbf{78} (1988), no.~4, 535--581.

\bibitem{NualartVives}
D.~Nualart and J.~Vives,
\emph{Chaos expansions and local times},
Publ. Mat. \textbf{36} (1992), no.~2B, 827--836.

\bibitem{ReedSimonI}
M.~Reed and B.~Simon,
\emph{Methods of Modern Mathematical Physics. I: Functional Analysis},
revised and enlarged ed., Academic Press, New York, 1980.

\bibitem{RevuzYor}
D.~Revuz and M.~Yor,
\emph{Continuous Martingales and Brownian Motion},
3rd ed., Grundlehren der Mathematischen Wissenschaften, vol.~293,
Springer-Verlag, Berlin, 1999.

\end{thebibliography}

\end{document}
